% Leçon 32 — Grammaire : compl. d’état — Chinois 1ère A — Cours signé OnBuch+
% Programme officiel MINESEC. Compiler avec : tectonic ZHA32-grammaire-compl-detat.tex
\newcommand{\DOCMATIERE}{Chinois}
\newcommand{\DOCNIVEAU}{1ère A}
\newcommand{\DOCMODULE}{Module 4 : Activités économiques et monde du travail}
\newcommand{\DOCLECON}{ZHA32}
\newcommand{\DOCTITRE}{Leçon 32 — Grammaire : compl. d’état}
% OnBuch+ — préambule « cours signés » ENRICHI (compilé avec tectonic / XeLaTeX)
% Système visuel « Pop » d'OnBuch+ : fond crème (jamais blanc), bordures encre
% épaisses, ombres « dures » décalées, titres Archivo Black en capitales.
% Version MATHÉMATIQUES : graphes (pgfplots/tikz), boîtes pédagogiques
% (définition, propriété, méthode, attention, le savais-tu, exercice résolu,
% exercice, TP, bilan), page de garde et signature OnBuch+ ancrée en bas.
%
% Macros attendues dans main.tex AVANT \input{preamble} :
%   \DOCMATIERE  \DOCNIVEAU  \DOCMODULE  \DOCLECON (numéro)  \DOCTITRE
\documentclass[11pt]{article}

\usepackage{fontspec}
\usepackage[french]{babel} % français = langue principale ; le chinois via \zh{..} / textebox (xeCJK)
\usepackage{xeCJK}
\usepackage{pifont} % \ding{51} \ding{55} utilisés par les modèles
\usepackage[a4paper,margin=2cm,top=3.4cm,bottom=2.8cm,headheight=1.4cm,footskip=1.3cm]{geometry}
\usepackage{amsmath,amssymb,amsfonts}
\usepackage{mathtools} % \xmapsto, \coloneqq... souvent utilisés par les modèles
\usepackage{cancel} % simplifications barrées (bilans de forces)
% \abs{z} (module, valeur absolue) et \norm{u} (norme), utilisés par les modèles
\providecommand{\abs}{}\let\abs\relax\DeclarePairedDelimiter{\abs}{\lvert}{\rvert}
\providecommand{\norm}{}\let\norm\relax\DeclarePairedDelimiter{\norm}{\lVert}{\rVert}
\usepackage[table]{xcolor} % [table] : fournit aussi \rowcolors (alternance de lignes)
\usepackage{pagecolor}
\usepackage{tikz}
\usetikzlibrary{arrows.meta,positioning,calc,shapes.geometric,shapes.misc,shapes.arrows,decorations.pathmorphing,decorations.pathreplacing,decorations.markings,patterns,shadows,mindmap,trees,fit,backgrounds,matrix,intersections,angles,quotes}
\usepackage{pgfplots}
\usepackage[version=4]{mhchem} % équations nucléaires \ce{^{235}_{92}U}
\usepackage[european,siunitx]{circuitikz} % schémas électriques
\pgfplotsset{compat=1.18}
\usepgfplotslibrary{fillbetween,groupplots} % aires entre courbes (calcul intégral)
\usepackage{enumitem}
\usepackage{fancyhdr}
% [most] charge toutes les bibliothèques tcolorbox (listings, minted, raster,
% documentation...) et gonfle inutilement la mémoire TeX sur un document long
% (~300 boîtes) : on ne charge que ce qui est réellement utilisé.
\usepackage[skins,breakable]{tcolorbox}
\usepackage{graphicx}
\usepackage{colortbl}
\usepackage{booktabs}
\usepackage{tabularx}
\usepackage{diagbox} % case d'en-tête diagonale des tableaux à double entrée
% Colonnes extensibles centrée / gauche / droite (C, L, R), souvent utilisées par les modèles
\newcolumntype{C}{>{\centering\arraybackslash}X}
\newcolumntype{L}{>{\raggedright\arraybackslash}X}
\newcolumntype{R}{>{\raggedleft\arraybackslash}X}
\usepackage{array}
\usepackage{multicol}
\usepackage{multirow}
\usepackage{siunitx}
% parse-numbers=false : accepte aussi \SI{2,5 \times 10^{-3}}{...}
\sisetup{locale=FR,output-decimal-marker={,},group-minimum-digits=4,parse-numbers=false}
\DeclareSIUnit\ohm{\ensuremath{\symup{\Omega}}} % U+2126 absent de la police
\DeclareSIUnit\division{div} % réglages d'oscilloscope (ms/div, V/div)
\DeclareSIUnit\tr{tr} % tours (vitesse de rotation, ex. \SI{1500}{\tr\per\minute})
\DeclareSIUnit\molar{M} % concentration molaire (ex. \SI{3}{\molar}, \SI{1}{\micro\molar})
\DeclareSIUnit\an{an} % année (le modèle confond parfois avec \year, un compteur TeX) — ex. \SI{5}{\kilo\meter\per\an}
\DeclareSIUnit\FCFA{FCFA} % franc CFA (ex. \SI{500}{\FCFA\per\kilo\gram})
\DeclareSIUnit\degreeBrix{\textdegree Brix} % teneur en sucre d'un sirop/jus (réfractométrie)
\DeclareSIUnit\month{mois} % le modèle utilise parfois \month, un compteur TeX, comme unité de durée
\usepackage{qrcode}
% \degree : commande fréquemment utilisée par les modèles pour un angle ou une
% température en dehors de \SI{}{} (non fournie par les paquets chargés).
\providecommand{\degree}{^{\circ}}
% \qed (fin de démonstration), utilisé par les modèles sans amsthm
\providecommand{\qed}{\hfill$\square$}
% \barre : double barre des valeurs interdites dans un tableau de variations (tkz-tab)
\providecommand{\barre}{\ensuremath{\|}}
% environnement proof (amsthm non chargé) utilisé par les modèles
\ifdefined\proof\else\newenvironment{proof}[1][Démonstration]{\par\noindent\textit{#1.}\ }{\qed\par}\fi
\usepackage{lastpage}
\usepackage{hyperref}
\hypersetup{hidelinks}

% ============================================================
% Palette « Pop » OnBuch+ — mêmes hex que la classe P (plus_ui.dart)
% ============================================================
\definecolor{poporange}{HTML}{F59321}
\definecolor{popdark}{HTML}{E07A0C}
\definecolor{popcream}{HTML}{FFF6E8}
\definecolor{popink}{HTML}{1C1714}
\definecolor{poppurple}{HTML}{7A5AE0}
\definecolor{popgreen}{HTML}{2E9E6A}
\definecolor{poppink}{HTML}{E0457B}
\definecolor{popblue}{HTML}{2F7DE1}
\definecolor{popgold}{HTML}{C98A12}
\definecolor{popmuted}{HTML}{8A7F76}
\definecolor{popline}{HTML}{EADFD2}
\definecolor{popgray}{HTML}{8A7F76}
\colorlet{popgrey}{popgray}
% Alias tolérants (le modèle nomme parfois les couleurs différemment)
\colorlet{popwhite}{white}
\colorlet{popblack}{popink}
\colorlet{popred}{poppink}
\colorlet{popyellow}{popgold}
% Teintes claires (fonds de tableaux/encadrés) — le modèle nomme parfois
% les couleurs avec un suffixe L (ex. popblueL) sans les avoir définies.
\colorlet{poporangeL}{poporange!20}
\colorlet{popdarkL}{popdark!20}
\colorlet{poppurpleL}{poppurple!20}
\colorlet{popgreenL}{popgreen!20}
\colorlet{poppinkL}{poppink!20}
\colorlet{popblueL}{popblue!20}
\colorlet{popgoldL}{popgold!20}
\colorlet{popredL}{popred!20}
\colorlet{popgrayL}{popgray!20}
% Teintes claires pour fonds de tableaux / zones
\colorlet{popblueL}{popblue!10!white}
\colorlet{poporangeL}{poporange!14!white}
\colorlet{popgreenL}{popgreen!12!white}
\colorlet{poppurpleL}{poppurple!12!white}
\colorlet{poppinkL}{poppink!12!white}
\colorlet{popgoldL}{popgold!14!white}

\pagecolor{popcream}

% ============================================================
% Typographie
% ============================================================
\defaultfontfeatures{Ligatures=TeX}
\setmainfont{PlusJakartaSans-Medium.ttf}[Path=../../../fonts/,BoldFont=PlusJakartaSans-Bold.ttf]
% Symboles, API et emojis absents de la police principale : repli sur DejaVu Sans (caractères actifs)
\newfontface\symfont{DejaVuSans.ttf}[Path=../../../fonts/]
\makeatletter
\newcommand\symactive[1]{\catcode#1=13 \begingroup\lccode`\~=#1 \lowercase{\endgroup\def~}{{\symfont\char#1}}}
\newcommand\symblank[1]{\catcode#1=13 \begingroup\lccode`\~=#1 \lowercase{\endgroup\def~}{}}
\newcommand\symhyph[1]{\catcode#1=13 \begingroup\lccode`\~=#1 \lowercase{\endgroup\def~}{\mbox{-}}}
\newcount\symc
\newcommand\symrange[2]{\symc=#1 \@whilenum\symc<#2 \do{\expandafter\symactive\expandafter{\the\symc}\advance\symc by 1 }}
\newcommand\symblankrange[2]{\symc=#1 \@whilenum\symc<#2 \do{\expandafter\symblank\expandafter{\the\symc}\advance\symc by 1 }}
\makeatother
\symrange{"0250}{"02FF}\symactive{"2713}\symactive{"2714}\symactive{"2717}\symactive{"2718}\symactive{"2610}\symactive{"2611}\symactive{"2612}
\symrange{"2600}{"27BF}
\catcode"2705=13 \begingroup\lccode`\~="2705 \lowercase{\endgroup\def~}{{\symfont\char"2713}}\symblank{"26BD}\symblank{"26A1}
\symrange{"2460}{"257F}\symactive{"223C}\symactive{"2248}\symactive{"2260}
\symactive{"01F8}\symactive{"01F9}\symactive{"1E3E}\symactive{"1E3F}
\symhyph{"2011}\symhyph{"2E1A}\symblankrange{"1F300}{"1FAFF}\symblank{"FE0F}
% Caractères chinois : WenQuanYi Zen Hei (sous-ensemble GB2312 + pinyin), gras/italique simulés
\setCJKmainfont{WenQuanYiZenHei-subset.ttf}[Path=../../../fonts/,AutoFakeBold=3,AutoFakeSlant=0.2]
\xeCJKsetup{CJKspace=false}
% Pinyin : voyelles à caron (ǎ ǐ ǒ ǔ ǖ ǘ ǚ ǜ) absentes de la police principale -> caractères actifs
\newfontface\pyfont{WenQuanYiZenHei-subset.ttf}[Path=../../../fonts/]
\newcommand\pyactive[1]{\catcode#1=13 \begingroup\lccode`\~=#1 \lowercase{\endgroup\def~}{{\pyfont\char#1}}}
\pyactive{"01CD}\pyactive{"01CE}\pyactive{"01CF}\pyactive{"01D0}\pyactive{"01D1}\pyactive{"01D2}\pyactive{"01D3}\pyactive{"01D4}\pyactive{"01D5}\pyactive{"01D6}\pyactive{"01D7}\pyactive{"01D8}\pyactive{"01D9}\pyactive{"01DA}\pyactive{"01DB}\pyactive{"01DC}
% Maths en sans-serif (Fira Math) avec chiffres et lettres droites de Plus
% Jakarta, pour que les formules se fondent dans le texte.
\usepackage[math-style=ISO,bold-style=ISO]{unicode-math}
\setmathfont{FiraMath-Regular.otf}[Path=../../../fonts/]
\setmathfont{PlusJakartaSans-Medium.ttf}[Path=../../../fonts/,range={up/num,up/{Latin,latin},\mathup}]
\usepackage{icomma} % virgule décimale sans espace en mode math
% Caractères mathématiques Unicode absents des polices de texte (Plus Jakarta,
% Archivo Black) : rendus par leur équivalent mathématique, en texte comme en
% formule — sinon ils s'affichent en carré vide (ex. « Arithmétique dans ℤ »).
\usepackage{newunicodechar}
\newunicodechar{ℝ}{\ensuremath{\mathbb{R}}}
\newunicodechar{ℤ}{\ensuremath{\mathbb{Z}}}
\newunicodechar{ℂ}{\ensuremath{\mathbb{C}}}
\newunicodechar{ℕ}{\ensuremath{\mathbb{N}}}
\newunicodechar{ℚ}{\ensuremath{\mathbb{Q}}}
\newunicodechar{𝔻}{\ensuremath{\mathbb{D}}}
\newunicodechar{α}{\ensuremath{\alpha}}
\newunicodechar{β}{\ensuremath{\beta}}
\newunicodechar{γ}{\ensuremath{\gamma}}
\newunicodechar{δ}{\ensuremath{\delta}}
\newunicodechar{ε}{\ensuremath{\varepsilon}}
\newunicodechar{θ}{\ensuremath{\theta}}
\newunicodechar{λ}{\ensuremath{\lambda}}
\newunicodechar{μ}{\ensuremath{\mu}}
\newunicodechar{σ}{\ensuremath{\sigma}}
\newunicodechar{τ}{\ensuremath{\tau}}
\newunicodechar{φ}{\ensuremath{\varphi}}
\newunicodechar{ψ}{\ensuremath{\psi}}
\newunicodechar{ω}{\ensuremath{\omega}}
\newunicodechar{Σ}{\ensuremath{\Sigma}}
% majuscules produites par \MakeUppercase dans les titres (α-aminés -> Α)
\newunicodechar{Α}{\ensuremath{\alpha}}
\newunicodechar{Β}{\ensuremath{\beta}}
\newunicodechar{Γ}{\ensuremath{\Gamma}}
\newunicodechar{Δ}{\ensuremath{\Delta}}
\newunicodechar{Ε}{\ensuremath{\varepsilon}}
\newunicodechar{Θ}{\ensuremath{\Theta}}
\newunicodechar{Λ}{\ensuremath{\Lambda}}
\newunicodechar{Μ}{\ensuremath{\mu}}
\newunicodechar{Τ}{\ensuremath{\tau}}
\newunicodechar{Φ}{\ensuremath{\Phi}}
\newunicodechar{Ψ}{\ensuremath{\Psi}}
\newunicodechar{Ω}{\ensuremath{\Omega}}
\newunicodechar{↦}{\ensuremath{\mapsto}}
\newunicodechar{∈}{\ensuremath{\in}}
\newunicodechar{∉}{\ensuremath{\notin}}
\newunicodechar{∧}{\ensuremath{\wedge}}
\newunicodechar{⇔}{\ensuremath{\Leftrightarrow}}
\newunicodechar{⟺}{\ensuremath{\Longleftrightarrow}}
\newunicodechar{⇒}{\ensuremath{\Rightarrow}}
\newunicodechar{⟹}{\ensuremath{\Longrightarrow}}
\newunicodechar{∘}{\ensuremath{\circ}}
\newunicodechar{∪}{\ensuremath{\cup}}
\newunicodechar{∩}{\ensuremath{\cap}}
\newunicodechar{≡}{\ensuremath{\equiv}}
\newunicodechar{⊂}{\ensuremath{\subset}}
\newunicodechar{⊥}{\ensuremath{\perp}}
\newunicodechar{∃}{\ensuremath{\exists}}
\newunicodechar{∀}{\ensuremath{\forall}}
\newunicodechar{∛}{\ensuremath{\sqrt[3]{\,}}}
\newunicodechar{∜}{\ensuremath{\sqrt[4]{\,}}}
\newunicodechar{⃗}{\ensuremath{{}^{\to}}}
\newunicodechar{ⁿ}{\ifmmode^{n}\else\textsuperscript{\ensuremath{n}}\fi}
\newunicodechar{ˣ}{\ifmmode^{x}\else\textsuperscript{\ensuremath{x}}\fi}
\newunicodechar{ᵇ}{\ifmmode^{b}\else\textsuperscript{\ensuremath{b}}\fi}
\newunicodechar{ʳ}{\ifmmode^{r}\else\textsuperscript{\ensuremath{r}}\fi}
\newunicodechar{ᵘ}{\ifmmode^{u}\else\textsuperscript{\ensuremath{u}}\fi}
\newunicodechar{ʸ}{\ifmmode^{y}\else\textsuperscript{\ensuremath{y}}\fi}
\newunicodechar{ᵃ}{\ifmmode^{a}\else\textsuperscript{\ensuremath{a}}\fi}
\newunicodechar{ᵉ}{\ifmmode^{e}\else\textsuperscript{\ensuremath{e}}\fi}
\newunicodechar{ᵏ}{\ifmmode^{k}\else\textsuperscript{\ensuremath{k}}\fi}
\newunicodechar{ᵖ}{\ifmmode^{p}\else\textsuperscript{\ensuremath{p}}\fi}
\newunicodechar{ᴺ}{\ifmmode^{N}\else\textsuperscript{\ensuremath{N}}\fi}
\newunicodechar{ᶜ}{\ifmmode^{c}\else\textsuperscript{\ensuremath{c}}\fi}
\newunicodechar{ʰ}{\ifmmode^{h}\else\textsuperscript{\ensuremath{h}}\fi}
\newunicodechar{ᵠ}{\ifmmode^{q}\else\textsuperscript{\ensuremath{q}}\fi}
\newunicodechar{ₙ}{\ifmmode_{n}\else\textsubscript{\ensuremath{n}}\fi}
\newunicodechar{ₐ}{\ifmmode_{a}\else\textsubscript{\ensuremath{a}}\fi}
\newunicodechar{ᵢ}{\ifmmode_{i}\else\textsubscript{\ensuremath{i}}\fi}
\newunicodechar{ᵣ}{\ifmmode_{r}\else\textsubscript{\ensuremath{r}}\fi}
\newunicodechar{ₖ}{\ifmmode_{k}\else\textsubscript{\ensuremath{k}}\fi}
\newunicodechar{ᵦ}{\ifmmode_{\beta}\else\textsubscript{\ensuremath{\beta}}\fi}
\newunicodechar{ₓ}{\ifmmode_{x}\else\textsubscript{\ensuremath{x}}\fi}
\newfontfamily\popdisplay{ArchivoBlack-Regular.ttf}[Path=../../../fonts/]
\newfontfamily\poplogo{SpaceGrotesk-Bold.ttf}[Path=../../../fonts/]
\newfontfamily\popsemi{PlusJakartaSans-SemiBold.ttf}[Path=../../../fonts/]
\newfontfamily\popxbold{PlusJakartaSans-ExtraBold.ttf}[Path=../../../fonts/]
\newfontfamily\popxboldls{PlusJakartaSans-ExtraBold.ttf}[Path=../../../fonts/,LetterSpace=3.0]

\setlist[enumerate]{leftmargin=1.7em,itemsep=5pt,topsep=4pt}
\setlist[itemize]{leftmargin=1.7em,itemsep=4pt,topsep=4pt,label={\color{poporange}\textbullet}}
\setlength{\parindent}{0pt}
\setlength{\parskip}{5pt}
\renewcommand{\arraystretch}{1.25}

% pgfplots : style Pop par défaut
\pgfplotsset{
  every axis/.append style={axis line style={popink,line width=1pt},tick style={popink},
    grid style={popline,line width=0.5pt},label style={font=\small},tick label style={font=\footnotesize}},
  popcurve/.style={poporange,line width=1.8pt,no markers,smooth},
}
% Même style disponible dans un \draw TikZ simple (hors pgfplots)
\tikzset{popcurve/.style={poporange,line width=1.8pt}}
\tikzset{popgrid/.style={popline,very thin}}
\colorlet{popcurve}{poporange} % le nom du style est parfois utilisé comme couleur

% ============================================================
% Pill / squiggle
% ============================================================
\newcommand{\poppill}[3]{%
  \tikz[baseline=-0.55ex]\node[draw=popink,line width=1.2pt,rounded corners=2.6pt,%
    fill=#2,inner xsep=6.5pt,inner ysep=3pt]{%
    \popxboldls\fontsize{7}{7}\selectfont\color{#3}\MakeUppercase{#1}};%
}
\newcommand{\popsquiggle}[1]{%
  \tikz[baseline=-0.4ex]{
    \draw[#1,line width=2.6pt,line cap=round]
      plot [smooth] coordinates {(0,0.09) (0.34,0.01) (0.68,0.21) (1.02,0.01) (1.36,0.10)};
  }%
}
% Mot-clé surligné (marqueur orange sous le texte)
\newcommand{\cle}[1]{{\popxbold\color{popdark}#1}}

% ============================================================
% En-tête / pied
% ============================================================
\pagestyle{fancy}
\fancyhf{}
\renewcommand{\headrulewidth}{0pt}
\renewcommand{\footrulewidth}{0pt}
\lhead{\raisebox{-0.15ex}{\poplogo\fontsize{13}{13}\selectfont\color{popink}OnBuch\color{poporange}+}}
\rhead{\poppill{\DOCMATIERE\ \textbullet\ \DOCNIVEAU\ \textbullet\ Leçon \DOCLECON}{white}{popink}}
\lfoot{\popxbold\fontsize{7.6}{7.6}\selectfont\color{popmuted}\MakeUppercase{Cours signé OnBuch+}\ \textbullet\ {\popsemi\DOCTITRE}}
\rfoot{\tikz[baseline=-0.6ex]\node[rounded corners=8.5pt,fill=popink,minimum height=17pt,minimum width=17pt,inner xsep=5pt,inner sep=0]{\popxbold\fontsize{8}{8}\selectfont\color{popcream}\thepage\,/\,\pageref*{LastPage}};}

% ============================================================
% Titres
% ============================================================
\newcommand{\chaptitre}[2]{%
  \begin{tcolorbox}[enhanced,colback=poporange,colframe=popink,boxrule=2.6pt,arc=11pt,
      drop shadow=popink,left=15pt,right=15pt,top=12pt,bottom=11pt,before skip=3pt,after skip=18pt]
    \poppill{#2}{popink}{popcream}\par\vspace{9pt}
    {\raggedright\hyphenpenalty=10000\exhyphenpenalty=10000\popdisplay\fontsize{22}{24}\selectfont\color{popcream}\MakeUppercase{#1}\par}
  \end{tcolorbox}
}
% Section numérotée : pastille orange avec le numéro + titre Archivo
\newcounter{popsec}
\newcommand{\coursec}[1]{%
  \stepcounter{popsec}\setcounter{popsub}{0}%
  \par\vspace{16pt}\noindent
  \begin{minipage}{\linewidth}
  \tikz[baseline=-0.7ex]\node[draw=popink,line width=1.6pt,fill=poporange,rounded corners=4pt,minimum size=22pt,inner sep=0]{\popdisplay\fontsize{12}{12}\selectfont\color{popcream}\thepopsec};%
  \hspace{8pt}{\popdisplay\fontsize{15}{17}\selectfont\color{popink}\MakeUppercase{#1}}\par
  \vspace{4pt}\noindent{\color{poporange}\rule{36pt}{3.6pt}}
  \end{minipage}\par\nopagebreak\vspace{8pt}
}
\newcounter{popsub}
\newcommand{\courssub}[1]{%
  \stepcounter{popsub}%
  \par\vspace{9pt}\noindent{\popxbold\fontsize{11.5}{13}\selectfont\color{popink}{\color{poporange}\thepopsec.\thepopsub}\hspace{6pt}#1}\par\nopagebreak\vspace{4pt}
}

% ============================================================
% Boîtes pédagogiques — même recette : fond blanc, bordure épaisse colorée,
% ombre dure encre, badge plein. Titre optionnel [#1].
% ============================================================
\tcbset{popbase/.style={breakable,enhanced,colback=white,boxrule=2.3pt,arc=9pt,
  shadow={3pt}{-3pt}{0pt}{popink},left=14pt,right=14pt,top=11pt,bottom=11pt,before skip=11pt,after skip=15pt,
  fontupper=\popsemi\fontsize{10.6}{14.5}\selectfont\color{popink},
  fontlower=\popsemi\fontsize{10.6}{14.5}\selectfont\color{popink}}}
% Coche dessinée (le glyphe ✓ n'existe pas dans les polices du document)
\newcommand{\popcheck}[1]{\tikz[baseline=-0.5ex]\draw[#1,line width=1.6pt,line cap=round,line join=round](-0.13,0.0)--(-0.04,-0.09)--(0.13,0.1);}
\providecommand{\coche}{\popcheck{popgreen}}
\providecommand{\resultat}[1]{\par\smallskip\noindent\fcolorbox{popgreen}{popgreenL}{\parbox{\dimexpr\linewidth-2\fboxsep-2\fboxrule\relax}{\textbf{Résultat.} #1}}\par\smallskip}
\newcommand{\popboxhead}[3]{\poppill{#1}{#2}{white}\ifx\relax#3\relax\else\hspace{6pt}{\popxbold\fontsize{10.6}{12}\selectfont\color{popink}#3}\fi\par\vspace{7pt}}

\newtcolorbox{aretenir}[1][]{popbase,colframe=popgreen,before upper={\popboxhead{À retenir}{popgreen}{#1}}}
% Texte en chinois (document de lecture, dialogue, extrait) : cadre
% d'étude. Un titre optionnel (source, auteur) s'affiche dans l'en-tête.
\newtcolorbox{textebox}[1][]{popbase,colframe=popblue,before upper={\popboxhead{文本 · Texte}{popblue}{#1}},after upper={}}
% Marques ✓ / ✗ (forme correcte / fautive) souvent utilisées par les modèles
\providecommand{\cross}{\textcolor{poppink}{\ensuremath{\times}}}
\providecommand{\xmark}{\cross}\providecommand{\crossmark}{\cross}\providecommand{\cmark}{\ensuremath{\checkmark}}
% Ligne de réponse / justification à compléter (inventée par les modèles)
\providecommand{\justification}[1]{\par\noindent\textit{Justification~:} \dotfill\par}
% \textipa (package tipa non chargé) : la police ne couvre pas l'API -> simple passage
\providecommand{\textipa}[1]{#1}
% Soulignement (ulem non chargé) : \uline, \ul
\providecommand{\uline}[1]{\underline{#1}}\providecommand{\ul}[1]{\underline{#1}}
% \glossaire{\item ...} inventée par les modèles : liste de glossaire
\providecommand{\glossaire}[1]{\begin{itemize}#1\end{itemize}}
% \alert (beamer) inventée par les modèles : texte mis en évidence
\providecommand{\alert}[1]{\textbf{\textcolor{poppink}{#1}}}
% variantes ASCII d'ordinaux écrites en macro
\providecommand{\iere}{\textsuperscript{re}}\providecommand{\ieme}{\textsuperscript{e}}\providecommand{\ere}{\textsuperscript{re}}\providecommand{\eme}{\textsuperscript{e}}\providecommand{\er}{\textsuperscript{er}}\providecommand{\ie}{\textsuperscript{e}}
% Ordinaux français écrits en macro par les modèles : 1\ière, 3\ième
\newcommand{\ière}{\textsuperscript{re}}\newcommand{\ième}{\textsuperscript{e}}
% \exemple (inventée par les modèles) : étiquette « Exemple : »
\providecommand{\exemple}{\textbf{Exemple~:}\ }
% \square utilisable aussi hors mode math (cases à cocher)
\AtBeginDocument{\let\squareMath\square\renewcommand{\square}{\ensuremath{\squareMath}}}
% Mot ou phrase en chinois à l'intérieur d'un paragraphe français
\newcommand{\zh}[1]{\ifmmode\text{#1}\else{#1}\fi}
\let\eng\zh \let\esp\zh \let\all\zh % alias tolérés
% flèches d'intonation inventées par les modèles
\providecommand{\textsearrow}{}\renewcommand{\textsearrow}{\ensuremath{\searrow}}\providecommand{\textnearrow}{}\renewcommand{\textnearrow}{\ensuremath{\nearrow}}
% \fr{..} : mot français (inventé par les modèles)
\providecommand{\fr}[1]{\textit{#1}}
% zhbox : encadré de texte chinois inventé par les modèles
\newenvironment{zhbox}{\begin{quote}}{\end{quote}}
% \courssubsub : titre de niveau 3 inventé par les modèles
\providecommand{\courssubsub}[1]{\par\medskip\noindent\textbf{#1}\par\smallskip}
% \zhbf : texte chinois en gras inventé par les modèles
\providecommand{\zhbf}[1]{\textbf{#1}}
% \vs : « versus » inventé par les modèles
\providecommand{\vs}{\textit{vs}\ }
% \blank : case à compléter inventée par les modèles
\providecommand{\blank}[1][]{\underline{\hspace{1.6em}}}
\newtcolorbox{exemplebox}[1][]{popbase,colframe=poporange,before upper={\popboxhead{Exemple}{poporange}{#1}}}
\newtcolorbox{definition}[1][]{popbase,colframe=popblue,before upper={\popboxhead{Définition}{popblue}{#1}}}
\newtcolorbox{propriete}[1][]{popbase,colframe=poppurple,before upper={\popboxhead{Propriété}{poppurple}{#1}}}
\newtcolorbox{methode}[1][]{popbase,colframe=popgold,before upper={\popboxhead{Méthode}{popgold}{#1}}}
\newtcolorbox{attention}[1][]{popbase,colframe=poppink,before upper={\popboxhead{Attention}{poppink}{#1}}}
\newtcolorbox{experience}[1][]{popbase,colframe=popblue,colback=popblueL,before upper={\popboxhead{Expérience}{popblue}{#1}}}
% « Le savais-tu ? » : carte violette pleine (contexte, culture, Cameroun)
\newtcolorbox{savaistu}[1][]{popbase,colback=poppurple,colframe=popink,
  fontupper=\popsemi\fontsize{10.4}{14.2}\selectfont\color{white},
  before upper={\poppill{Le savais-tu\,?}{white}{poppurple}\ifx\relax#1\relax\else\hspace{6pt}{\popxbold\color{white}#1}\fi\par\vspace{7pt}}}
% Exercice résolu : énoncé en haut, \tcblower puis la solution sur fond crème
\newcounter{popexo}\newcounter{popexr}
\newtcolorbox{exoresolu}[1][]{popbase,colframe=poporange,colbacklower=popcream,
  segmentation style={popink,line width=1.2pt,dashed},
  before upper={\stepcounter{popexr}\popboxhead{Exercice résolu \thepopexr}{poporange}{#1}},
  before lower={\poppill{Solution}{popink}{popcream}\par\vspace{6pt}}}
% Exercice (non résolu en place) — la correction est dans la partie Corrigés
\newtcolorbox{exercice}[1][]{popbase,colframe=popink,
  before upper={\stepcounter{popexo}\popboxhead{Exercice \thepopexo}{popink}{#1}}}
% Corrigé
\newtcolorbox{corrige}[1][]{popbase,colframe=popgreen,colback=popgreenL,
  before upper={\popboxhead{Corrigé}{popgreen}{#1}}}
% Objectifs / prérequis / bilan
\newtcolorbox{objectifs}{popbase,colframe=popink,colback=poporangeL,
  before upper={\popboxhead{Objectifs de la leçon}{popink}{}}}
\newtcolorbox{prerequis}{popbase,colframe=popink,colback=popblueL,
  before upper={\popboxhead{Prérequis}{popblue}{}}}
\newtcolorbox{bilan}{popbase,colframe=popink,colback=white,boxrule=2.8pt,
  before upper={\popboxhead{Fiche bilan}{poporange}{L'essentiel à connaître pour l'examen}}}
% Figure encadrée avec légende
\newtcolorbox{popfigure}[1][]{enhanced,breakable=false,colback=white,colframe=popline,boxrule=1.4pt,arc=8pt,
  left=10pt,right=10pt,top=10pt,bottom=8pt,before skip=10pt,after skip=12pt,halign=center,
  lower separated=false,halign lower=center,
  fontlower=\popsemi\fontsize{9}{11.5}\selectfont\color{popmuted},
  #1}
\newcommand{\legende}[1]{\par\vspace{6pt}{\popsemi\fontsize{9}{11.5}\selectfont\color{popmuted}#1}}

% ============================================================
% Signature OnBuch+
% ============================================================
\newcommand{\poplogoinline}[1]{{\poplogo\fontsize{#1}{#1}\selectfont\color{popink}OnBuch\color{poporange}+}}
\newcommand{\signatureonbuch}{%
  \begin{tcolorbox}[enhanced,colback=popink,colframe=popink,boxrule=2.2pt,arc=10pt,
      drop shadow=poporange,left=14pt,right=14pt,top=10pt,bottom=10pt,before skip=0pt,after skip=0pt]
    \tikz[baseline=-0.7ex]\node[circle,fill=popgreen,draw=popcream,line width=1.3pt,minimum size=22pt,inner sep=0]{\popcheck{white}};%
    \hspace{9pt}\begin{minipage}[c]{0.62\linewidth}
      {\popxbold\fontsize{12}{13}\selectfont\color{popcream}Signé\ \textbullet\ {\poplogo OnBuch\color{poporange}+}}\par\vspace{2pt}
      {\raggedright\popsemi\fontsize{8}{10.5}\selectfont\color{popline}\MakeUppercase{Équipe pédagogique OnBuch+}\\\MakeUppercase{Conforme au programme officiel MINESEC}\par}
    \end{minipage}\hfill
    {\poplogo\fontsize{22}{22}\selectfont\color{popcream}OnBuch\color{poporange}+}
  \end{tcolorbox}%
}

% ============================================================
% Page de garde
% #1 titre · #2 matière · #3 niveau · #4 module · #5 n° leçon · #6 durée ·
% #7 description / objectifs (texte ou itemize)
% ============================================================
\newcommand{\pagegarde}[7]{%
  \thispagestyle{empty}
  \newgeometry{a4paper,margin=1.7cm,top=1.6cm,bottom=1.4cm}%
  \begin{tikzpicture}[remember picture,overlay]
    \fill[poporange!18] ([xshift=-3.2cm,yshift=-3.6cm]current page.north east) circle (4.6cm);
    \fill[poppurple!14] ([xshift=2.4cm,yshift=7.6cm]current page.south west) circle (3.8cm);
  \end{tikzpicture}%
  {\poplogo\fontsize{30}{30}\selectfont\color{popink}OnBuch\color{poporange}+}\hfill
  \raisebox{0.6ex}{\poppill{Cours signé \textbullet\ Premium}{popink}{popcream}}\par
  \vspace{1pt}\popsquiggle{poppurple}\par
  \vspace{8pt}
  \begin{tcolorbox}[enhanced,colback=poporange,colframe=popink,boxrule=2.8pt,arc=13pt,
      drop shadow=popink,left=18pt,right=18pt,top=12pt,bottom=13pt,after skip=10pt]
    \poppill{#2 \textbullet\ #3}{popink}{popcream}\hspace{5pt}\poppill{#4}{white}{popink}\par\vspace{8pt}
    {\popdisplay\fontsize{14}{14}\selectfont\color{popink}LEÇON #5}\par\vspace{4pt}
    {\raggedright\hyphenpenalty=10000\exhyphenpenalty=10000\popdisplay\fontsize{25}{27}\selectfont\color{popcream}\MakeUppercase{#1}\par}
    \vspace{8pt}
    {\popxbold\fontsize{9.5}{9.5}\selectfont\color{popink}\MakeUppercase{Durée conseillée : #6}}
  \end{tcolorbox}
  \begin{tcolorbox}[enhanced,colback=white,colframe=popink,boxrule=2.2pt,arc=10pt,
      drop shadow=popink,left=16pt,right=16pt,top=10pt,bottom=10pt,after skip=8pt]
    \poppill{Dans cette leçon}{poppurple}{white}\par\vspace{5pt}
    \popsemi\fontsize{9.3}{12.6}\selectfont\color{popink}#7
  \end{tcolorbox}
  \noindent\begin{minipage}[t]{0.487\linewidth}
    \begin{tcolorbox}[enhanced,colback=popgreen,colframe=popink,boxrule=2.2pt,arc=10pt,
        drop shadow=popink,left=11pt,right=11pt,top=10pt,bottom=10pt,height=3.1cm,valign=center]
      \tcbox[colback=white,colframe=popink,boxrule=1.3pt,arc=5pt,boxsep=0pt,nobeforeafter,
        left=3pt,right=3pt,top=3pt,bottom=3pt,valign=center]{\qrcode[height=1.9cm]{https://play.google.com/store/apps/details?id=com.luvvix.onbuch&hl=fr}}%
      \hspace{8pt}\begin{minipage}[c]{0.5\linewidth}
        {\popdisplay\fontsize{11}{12.5}\selectfont\color{white}\MakeUppercase{Télécharge OnBuch}}\par\vspace{4pt}
        {\raggedright\popsemi\fontsize{8.4}{11}\selectfont\color{white}Gratuit \textbullet\ Google Play\\Tuteur IA, annales, concours, résultats\par}
      \end{minipage}
    \end{tcolorbox}
  \end{minipage}\hfill
  \begin{minipage}[t]{0.487\linewidth}
    \begin{tcolorbox}[enhanced,colback=poppurple,colframe=popink,boxrule=2.2pt,arc=10pt,
        drop shadow=popink,left=11pt,right=11pt,top=10pt,bottom=10pt,height=3.1cm,valign=center]
      \tikz[baseline=-0.7ex]\node[circle,fill=white,draw=popink,line width=1.3pt,minimum size=1.6cm,inner sep=0]{\popdisplay\fontsize{24}{24}\selectfont\color{poppurple}+};%
      \hspace{8pt}\begin{minipage}[c]{0.6\linewidth}
        {\popdisplay\fontsize{11}{12.5}\selectfont\color{white}\MakeUppercase{Passe à OnBuch+}}\par\vspace{4pt}
        {\raggedright\popsemi\fontsize{8.4}{11}\selectfont\color{white}Tous les cours signés, Tuteur IA illimité, fiches PDF — onglet OnBuch+ de l'app\par}
      \end{minipage}
    \end{tcolorbox}
  \end{minipage}
  \vspace{8pt}
  \popcontenu
  \vfill
  \signatureonbuch
  \restoregeometry
  \newpage
}

% Rangée « Ce cours contient » (page de garde)
\newcommand{\poptile}[3]{% #1 couleur #2 gros texte #3 légende
  \begin{tcolorbox}[enhanced,colback=#1,colframe=popink,boxrule=2pt,arc=9pt,shadow={2.5pt}{-2.5pt}{0pt}{popink},
      left=6pt,right=6pt,top=9pt,bottom=9pt,halign=center,valign=center,height=1.75cm,width=\linewidth,nobeforeafter]
    {\popdisplay\fontsize{12.5}{13}\selectfont\color{white}\MakeUppercase{#2}}\par\vspace{3pt}
    {\centering\popsemi\fontsize{7.8}{9.5}\selectfont\color{white}#3\par}
  \end{tcolorbox}}
\newcommand{\popcontenu}{%
  \par\noindent{\popxboldls\fontsize{7.5}{7.5}\selectfont\color{popmuted}CE COURS SIGNÉ CONTIENT}\par\vspace{7pt}
  \noindent\begin{minipage}[t]{0.235\linewidth}\poptile{popblue}{Cours}{complet, illustré, pas à pas}\end{minipage}\hfill
  \begin{minipage}[t]{0.235\linewidth}\poptile{poporange}{Méthodes}{et exercices résolus}\end{minipage}\hfill
  \begin{minipage}[t]{0.235\linewidth}\poptile{poppink}{Exercices}{type examen + corrigés}\end{minipage}\hfill
  \begin{minipage}[t]{0.235\linewidth}\poptile{popgreen}{Fiche bilan}{l'essentiel à retenir}\end{minipage}\par}

% Sommaire
\newtcolorbox{sommaire}{enhanced,colback=white,colframe=popink,boxrule=2.2pt,arc=10pt,
  drop shadow=popink,left=18pt,right=18pt,top=13pt,bottom=13pt,before skip=6pt,after skip=20pt,
  before upper={\poppill{Sommaire}{poppurple}{white}\par\vspace{9pt}\popsemi\fontsize{10.6}{14.5}\selectfont\color{popink}}}

% Signature de fin de cours — poussée tout en bas de la dernière page
\newcommand{\findecours}{%
  \par\vfill\nopagebreak
  \begin{center}\popsquiggle{poporange}\end{center}\vspace{4pt}
  \signatureonbuch
}
\AtBeginDocument{\def\llbracket{\mathopen{[\mkern-3.5mu[}}\def\rrbracket{\mathclose{]\mkern-3.5mu]}}} % intervalles d'entiers (glyphe ⟦ absent de la police)
% Glyphes absents de Fira Math : redéfinis à partir de glyphes disponibles
\AtBeginDocument{%
  \def\setminus{\mathbin{\backslash}}\let\smallsetminus\setminus
  \def\vdots{\mathord{\vbox{\baselineskip3.2pt\lineskiplimit0pt\kern4pt\hbox{.}\hbox{.}\hbox{.}}}}%
  \def\ddots{\mathinner{\mkern1mu\raise6.5pt\hbox{.}\mkern2mu\raise3.6pt\hbox{.}\mkern2mu\raise0.7pt\hbox{.}\mkern1mu}}%
  \def\triangle{\mathord{\scalebox{0.9}{$\symup{\Delta}$}}}%
  \def\nabla{\mathord{\raisebox{\depth}{\scalebox{1}[-1]{$\symup{\Delta}$}}}}%
  \def\checkmark{\popcheck{popdark}}%
  \def\star{\mathbin{\ast}}\def\bigstar{\mathord{\ast}}%
  \def\diamond{\mathbin{\scalebox{0.6}{$\Diamond$}}}%
  \def\bigcup{\mathop{\mathchoice{\raisebox{-0.3em}{\scalebox{1.7}{$\cup$}}}{\scalebox{1.25}{$\cup$}}{\cup}{\cup}}\displaylimits}%
  \def\bigcap{\mathop{\mathchoice{\raisebox{-0.3em}{\scalebox{1.7}{$\cap$}}}{\scalebox{1.25}{$\cap$}}{\cap}{\cap}}\displaylimits}%
  \def\lesseqqgtr{\mathrel{\lessgtr}}\def\bigoplus{\mathop{\oplus}\displaylimits}%
}
\providecommand{\ssi}{si et seulement si} % abréviation fréquente
\AtBeginDocument{\providecommand{\wideparen}{\overparen}} % arc de cercle

%% FIGFIX-BEGIN v5 — correctifs globaux des figures (docs/onbuch-generation/tools/figfix)
\usepackage{adjustbox}\usepackage{etoolbox}\tcbuselibrary{raster}
% toute figure TikZ/pgfplots plus large que la ligne est réduite à la largeur disponible
\BeforeBeginEnvironment{tikzpicture}{\begin{adjustbox}{max width=\linewidth}}
\AfterEndEnvironment{tikzpicture}{\end{adjustbox}}
% légendes pgfplots lisibles par-dessus les courbes
\pgfplotsset{every axis/.append style={legend style={fill=white,fill opacity=0.92,text opacity=1,draw=black!15,inner sep=3pt}}}
% étiquettes d'arêtes (syntaxe edge["..."]) avec fond blanc
\tikzset{every edge quotes/.append style={fill=white,inner sep=1.2pt}}
%% FIGFIX-END
\begin{document}

\pagegarde{Leçon 32 — Grammaire : compl. d’état}{Chinois}{1ère A}{Module 4 : Activités économiques et monde du travail}{ZHA32}{2 h}{Cette leçon de grammaire présente le complément d'état (V + 得 + adjectif/adverbe), structure essentielle pour évaluer la manière dont une action est accomplie. Les élèves apprendront à construire, transformer et employer cette structure en évitant les pièges typiques des francophones.
\vspace{4pt}
\begin{multicols}{2}\small
\begin{itemize}
\item À la fin de cette leçon, je sais reconnaître un complément d'état dans une phrase chinoise
\item À la fin de cette leçon, je sais construire la structure Sujet + Verbe + 得 + adjectif
\item À la fin de cette leçon, je sais gérer la répétition du verbe quand celui-ci a un objet (V + O + V + 得 + adj)
\item À la fin de cette leçon, je sais former la négation avec 不 devant l'adjectif (V + 得 + 不 + adj)
\item À la fin de cette leçon, je sais poser une question avec 得 (V + 得 + adj + 吗 / V + 得 + 怎么样)
\item À la fin de cette leçon, je sais distinguer 得 (complément d'état) de 的 et 地
\end{itemize}\end{multicols}}

\begin{sommaire}
\begin{multicols}{2}
\begin{enumerate}
\item Observation et découverte
\item Schéma de la structure
\item Négation, question et emplois
\item Comparaison avec le français et pièges
\item Exercices d'application et de transformation
\item Activité d'intégration
\item Méthodes et exercices résolus
\item Exercices
\item Corrigés des exercices
\item Fiche bilan
\end{enumerate}
\end{multicols}
\end{sommaire}

\begin{objectifs}
\begin{itemize}
\item À la fin de cette leçon, je sais reconnaître un complément d'état dans une phrase chinoise
\item À la fin de cette leçon, je sais construire la structure Sujet + Verbe + 得 + adjectif
\item À la fin de cette leçon, je sais gérer la répétition du verbe quand celui-ci a un objet (V + O + V + 得 + adj)
\item À la fin de cette leçon, je sais former la négation avec 不 devant l'adjectif (V + 得 + 不 + adj)
\item À la fin de cette leçon, je sais poser une question avec 得 (V + 得 + adj + 吗 / V + 得 + 怎么样)
\item À la fin de cette leçon, je sais distinguer 得 (complément d'état) de 的 et 地
\item À la fin de cette leçon, je sais traduire correctement du français vers le chinois des phrases évaluant une action
\item À la fin de cette leçon, je sais utiliser le complément d'état dans un court texte descriptif sur mes activités
\end{itemize}
\end{objectifs}

\begin{prerequis}
\begin{itemize}
\item Les adjectifs qualificatifs et leur emploi avec 很 (leçon de Seconde)
\item La structure de base Sujet + Verbe + Objet
\item Les adverbes de degré : 很, 非常, 太, 不太
\item La particule 的 pour la détermination et 地 pour les compléments circonstanciels de manière
\item Le vocabulaire des activités quotidiennes et scolaires (Modules 1-3)
\end{itemize}
\end{prerequis}

\newpage

\coursec{Observation et découverte}

\courssub{Situation d'accroche : le match du lycée}

Après le match de football du lycée de Yaoundé, les élèves commentent dans la cour de récréation : « Il a bien joué ! », « Elle court vite ! », « Tu parles chinois couramment ! ». En français, c'est facile : on ajoute simplement un \emph{adverbe} après le verbe (jouer \emph{bien}, courir \emph{vite}, parler \emph{couramment}).

Mais en chinois, comment dit-on « Il parle chinois très bien » ? Regardez cette phrase :

\zh{他说汉语说得很好。} \emph{Tā shuō Hànyǔ shuō de hěn hǎo.} « Il parle (le) chinois très bien. »

Surprise : le verbe \zh{说} (parler) apparaît \textbf{deux fois} ! Pourquoi ? Et que fait ce petit caractère \zh{得} au milieu de la phrase ?

\textbf{Problématique :} comment le chinois exprime-t-il la \emph{manière} dont une action est réalisée (bien, vite, clairement, couramment) ? C'est toute la logique du \cle{complément d'état} (\zh{程度补语} \emph{chéngdù bǔyǔ}), introduit par la particule \zh{得}, que nous allons découvrir aujourd'hui.

\courssub{Premier contact : un dialogue au stade}

Lisons d'abord un court dialogue. Deux élèves, Étienne et son ami chinois Li Ming, discutent après le match.

\begin{textebox}[Au stade de Yaoundé]
\zh{艾天：李明，你看，他打篮球打得很好，但是足球踢得不太好。}\\
\emph{Àitiān : Lǐ Míng, nǐ kàn, tā dǎ lánqiú dǎ de hěn hǎo, dànshì zúqiú tī de bú tài hǎo.}\\
« Étienne : Li Ming, regarde, il joue très bien au basket, mais il ne joue pas très bien au football. »\\[4pt]
\zh{李明：是啊！可是他跑得很快。你汉语说得怎么样？}\\
\emph{Lǐ Míng : Shì a ! Kěshì tā pǎo de hěn kuài. Nǐ Hànyǔ shuō de zěnmeyàng?}\\
« Li Ming : C'est vrai ! Mais il court vite. Et toi, comment parles-tu chinois ? »\\[4pt]
\zh{艾天：我说得还可以，可是汉字写得不太好。}\\
\emph{Àitiān : Wǒ shuō de hái kěyǐ, kěshì Hànzì xiě de bú tài hǎo.}\\
« Étienne : Je parle passablement, mais je n'écris pas très bien les caractères. »
\end{textebox}

\textbf{Questions de compréhension :}
\begin{enumerate}
\item Que fait bien le joueur ? Que fait-il moins bien ?
\item Combien de fois la particule \zh{得} apparaît-elle dans ce dialogue ?
\item Quel mot suit \zh{得} à chaque fois : un nom, un verbe ou un adjectif ?
\end{enumerate}

\courssub{Lecture guidée des cinq exemples}

Observons maintenant attentivement les cinq phrases modèles de la leçon. Lisez chaque phrase à voix haute, en suivant le pinyin, puis la traduction.

\begin{center}
\begin{tabularx}{\linewidth}{|X|X|X|}
\hline
\rowcolor{popblueL} \textbf{Phrase chinoise} & \textbf{Pinyin} & \textbf{Traduction française} \\
\hline
他说汉语\textcolor{poporange}{\textbf{得}}很好。 & Tā shuō Hànyǔ shuō \textcolor{poporange}{\textbf{de}} hěn hǎo. & Il parle (le) chinois très bien. \\
\hline
她跑\textcolor{poporange}{\textbf{得}}很快。 & Tā pǎo \textcolor{poporange}{\textbf{de}} hěn kuài. & Elle court vite. \\
\hline
你汉字写\textcolor{poporange}{\textbf{得}}怎么样？ & Nǐ Hànzì xiě \textcolor{poporange}{\textbf{de}} zěnmeyàng? & Comment écris-tu les caractères ? \\
\hline
他唱歌唱\textcolor{poporange}{\textbf{得}}不太好。 & Tā chàng gē chàng \textcolor{poporange}{\textbf{de}} bú tài hǎo. & Il ne chante pas très bien. \\
\hline
老师说\textcolor{poporange}{\textbf{得}}清楚吗？ & Lǎoshī shuō \textcolor{poporange}{\textbf{de}} qīngchu ma? & Le professeur explique-t-il clairement ? \\
\hline
\end{tabularx}
\end{center}

\textbf{Repérage guidé.} Pour chaque phrase, répondez aux questions suivantes :

\begin{enumerate}
\item \textbf{Où est \zh{得} ?} Dans les cinq phrases, \zh{得} se trouve toujours \emph{immédiatement après un verbe} : \zh{说得}, \zh{跑得}, \zh{写得}, \zh{唱得}.
\item \textbf{Qu'y a-t-il après \zh{得} ?} Toujours un \emph{adjectif} (souvent précédé de \zh{很}, \zh{不}, \zh{太}) : \zh{很好} (très bien), \zh{很快} (vite), \zh{怎么样} (comment), \zh{不太好} (pas très bien), \zh{清楚} (clair).
\item \textbf{Que fait cette partie après \zh{得} ?} Elle \emph{évalue} l'action : elle dit \emph{comment} l'action est réalisée.
\end{enumerate}

\begin{exemplebox}[Hypothèse des élèves]
Formulons ensemble une hypothèse : la particule \zh{得} sert à introduire, après le verbe, un commentaire qui \textbf{évalue la manière} dont l'action est faite. Le français utilise un adverbe (« bien », « vite ») ; le chinois utilise \zh{得} + adjectif. Nous vérifierons cette hypothèse dans la suite de la leçon.
\end{exemplebox}

\courssub{De l'observation à la notion}

\begin{definition}[Le complément d'état]
Le \cle{complément d'état} (\zh{程度补语} \emph{chéngdù bǔyǔ}) est un complément placé \textbf{après le verbe}, introduit par la particule \zh{得}, qui \textbf{évalue la manière ou le degré de réalisation d'une action} : bien, mal, vite, clairement, couramment\ldots
\end{definition}

\begin{methode}[Comment repérer un complément d'état ?]
\begin{enumerate}
\item Cherche la particule \zh{得} dans la phrase.
\item Vérifie qu'elle est placée \textbf{juste après un verbe} (说, 跑, 写, 唱, 踢\ldots).
\item Vérifie qu'elle est suivie d'un \textbf{adjectif} ou d'une expression adjectivale : \zh{好} (bien), \zh{快} (vite), \zh{清楚} (clair), \zh{漂亮} (beau), \zh{怎么样} (comment), \zh{不太好} (pas très bien).
\item Si ces trois conditions sont réunies, tu as un complément d'état : la phrase évalue la manière de faire l'action.
\end{enumerate}
\end{methode}

\begin{attention}[Ne pas confondre les trois « de » !]
Le chinois possède trois particules homophones prononcées \emph{de} : \zh{的} (possession, qualification : \zh{我的书} « mon livre »), \zh{得} (complément d'état, après le verbe) et \zh{地} (complément circonstanciel, avant le verbe, niveau HSK supérieur). Ici, c'est bien \zh{得} (avec le radical 彳, « pas, marcher ») qu'il faut écrire. Erreur fréquente des francophones : écrire \zh{的} à la place de \zh{得}, car les deux se prononcent pareil. Retenez : \textbf{après un verbe, c'est 得}.
\end{attention}

\courssub{Carte mentale : les fonctions de 得}

\begin{popfigure}
\begin{tikzpicture}[mindmap, grow cyclic, every node/.style={concept, align=center, font=\small}, root concept/.append style={concept color=poporange, font=\large\bfseries}, level 1/.append style={level distance=3.2cm, sibling angle=120}]
\node[root concept, align=center]{得 de\\ complément\\ d'état}
  child[concept color=popblue] { node[align=center]{manière\\ 他说得很好。\\ \emph{il parle bien}} }
  child[concept color=popgreen] { node[align=center]{évaluation\\ 她写得怎么样？\\ \emph{comment écrit-elle ?}} }
  child[concept color=poppurple] { node[align=center]{degré\\ 他跑得很快。\\ \emph{il court très vite}} };
\end{tikzpicture}
\legende{Carte mentale : la particule 得 introduit un commentaire sur la manière, l'évaluation ou le degré de l'action.}
\end{popfigure}

\courssub{Comparaison avec le français}

\begin{center}
\begin{tabularx}{\linewidth}{|X|X|X|}
\hline
\rowcolor{popblueL} \textbf{Français} & \textbf{Chinois} & \textbf{Observation} \\
\hline
Il parle bien. & 他说得很好。Tā shuō de hěn hǎo. & L'adverbe « bien » devient 得 + 好. \\
\hline
Elle court vite. & 她跑得很快。Tā pǎo de hěn kuài. & « vite » = 得 + 快 (adjectif). \\
\hline
Il chante mal. & 他唱得不好。Tā chàng de bù hǎo. & « mal » = 得 + 不好 (négation après 得). \\
\hline
Le professeur explique clairement. & 老师说得很清楚。Lǎoshī shuō de hěn qīngchu. & « clairement » = 得 + 清楚. \\
\hline
\end{tabularx}
\end{center}

\begin{aretenir}[Ce qu'il faut retenir de l'observation]
\begin{itemize}
\item Le français utilise un \textbf{adverbe} ; le chinois utilise \zh{得} + \textbf{adjectif}.
\item \zh{得} se place \textbf{immédiatement après le verbe}.
\item L'adjectif après \zh{得} est presque toujours accompagné d'un modificateur : \zh{很} (très), \zh{不} (ne\ldots pas), \zh{太} (trop), ou un mot interrogatif comme \zh{怎么样}.
\item Quand le verbe a un objet (说汉语, 唱歌), on répète le verbe : \zh{他说汉语说得很好。} — ce point sera étudié en détail dans la section suivante.
\end{itemize}
\end{aretenir}

\begin{exoresolu}[Repérage du complément d'état]
\textbf{Énoncé.} Dans les phrases suivantes, souligne la particule \zh{得}, encadre l'adjectif qui la suit, puis traduis la phrase.\\
(a) 她跳舞跳得很漂亮。 (b) 我做作业做得很快。 (c) 他开车开得很好。
\tcblower
\textbf{Solution détaillée.}\\
(a) \zh{她跳舞跳\textbf{得}很[漂亮]。} \emph{Tā tiào wǔ tiào de hěn piàoliang.} « Elle danse très joliment. » — \zh{得} suit le verbe répété \zh{跳} ; l'adjectif est \zh{漂亮} (beau, joli).\\
(b) \zh{我做作业做\textbf{得}很[快]。} \emph{Wǒ zuò zuòyè zuò de hěn kuài.} « Je fais mes devoirs vite. » — \zh{得} suit le verbe répété \zh{做} ; l'adjectif est \zh{快} (rapide).\\
(c) \zh{他开车开\textbf{得}很[好]。} \emph{Tā kāi chē kāi de hěn hǎo.} « Il conduit bien. » — \zh{得} suit le verbe répété \zh{开} ; l'adjectif est \zh{好} (bien).\\
\textbf{Conclusion :} dans les trois cas, le schéma est identique : Sujet + Verbe (+ Objet) + Verbe + 得 + (很) + Adjectif.
\end{exoresolu}

\begin{exercice}[À toi de jouer]
Observe les phrases suivantes et dis si elles contiennent un complément d'état. Justifie ta réponse.\\
1. 我的老师很好。 2. 他说法语说得很流利。 3. 我们踢足球踢得不太好。 4. 她跑得怎么样？
\end{exercice}

\begin{savaistu}[Le compliment magique en Chine]
En Chine, féliciter quelqu'un avec \zh{你说得真好！} \emph{Nǐ shuō de zhēn hǎo!} « Tu parles vraiment bien ! » est un compliment très courant envers les étrangers qui apprennent le chinois. Les Chinois encouragent volontiers les apprenants, même pour quelques mots ! Si un jour tu visites Pékin ou Shanghai, tu entendras souvent cette phrase — et tu pourras répondre modestement : \zh{哪里哪里，我说得不太好。} \emph{Nǎli nǎli, wǒ shuō de bú tài hǎo.} « Mais non, mais non, je ne parle pas très bien. » La modestie est une valeur très appréciée dans la culture chinoise.
\end{savaistu}

\coursec{Observation et découverte}

\begin{textebox}[Dialogue : Après les cours au lycée de Yaoundé]
\zh{— 小张，你昨晚复习得怎么样？} \emph{— Xiǎo Zhāng, nǐ zuówǎn fùxí de zěnmeyàng?} \\
\zh{— 我复习得很认真，但是中文语法记得不太清楚。} \emph{— Wǒ fùxí de hěn rènzhēn, dànshì Zhōngwén yǔfǎ jì de bú tài qīngchu.} \\
\zh{— 别担心。王老师讲得很清楚，你听得仔细就行。} \emph{— Bié dānxīn. Wáng Lǎoshī jiǎng de hěn qīngchu, nǐ tīng de zǐxì jiù xíng.} \\
\zh{— 对了，你踢球踢得好吗？周六我们去踢足球。} \emph{— Duìle, nǐ tī qiú tī de hǎo ma? Zhōuliù wǒmen qù tī zúqiú.} \\
\zh{— 我踢得一般，但我跑得很快！} \emph{— Wǒ tī de yībān, dàn wǒ pǎo de hěn kuài!}
\tcblower
\textbf{Traduction :}
« — Petit Zhang, comment as-tu révisé hier soir ?
— J'ai révisé très sérieusement, mais je n'ai pas retenu la grammaire chinoise très clairement.
— Ne t'inquiète pas. Le prof Wang explique très clairement, tu n'as qu'à écouter attentivement.
— Au fait, est-ce que tu joues bien au foot ? Samedi on va jouer au foot.
— Je joue moyennement, mais je cours très vite ! »
\end{textebox}

\begin{methode}[Comment repérer le complément d'état]
\begin{enumerate}
\item Identifiez le verbe d'action principal (\zh{复习}, \zh{讲}, \zh{听}, \zh{踢}, \zh{跑}).
\item Cherchez la particule \zh{得} (de) immédiatement après le verbe (ou après le verbe répété).
\item Observez ce qui suit \zh{得} : un adverbe de degré (\zh{很}, \zh{不太}) + un adjectif qualitatif (\zh{认真}, \zh{清楚}, \zh{仔细}, \zh{好}, \zh{快}).
\item Ce bloc \zh{得 + adv. + adj.} évalue la \textbf{manière} ou le \textbf{degré} de réalisation de l'action. C'est le \cle{complément d'état} (程度补语).
\end{enumerate}
\end{methode}

\begin{aretenir}[Vocabulaire clé du dialogue]
\begin{center}
\begin{tabularx}{\linewidth}{|X|X|X|X|}
\hline
\rowcolor{popblueL} Caractères & Pinyin & Nature & Traduction \\
\hline
复习 & fùxí & verbe & réviser \\
\hline
认真 & rènzhēn & adjectif & sérieux, appliqué, soigneux \\
\hline
语法 & yǔfǎ & nom & grammaire \\
\hline
记 & jì & verbe & mémoriser, retenir \\
\hline
清楚 & qīngchu & adjectif & clair, clairement \\
\hline
担心 & dānxīn & verbe & s'inquiéter \\
\hline
仔细 & zǐxì & adjectif & attentif, minutieux \\
\hline
踢球 / 踢足球 & tī qiú / tī zúqiú & verbe + objet & jouer au football \\
\hline
一般 & yībān & adjectif & moyen, ordinaire \\
\hline
\end{tabularx}
\end{center}
\end{aretenir}

\coursec{Schéma de la structure}

Dans le dialogue, nous avons observé des phrases comme \zh{我复习得很认真} et \zh{王老师讲得很清楚}. Le petit mot \zh{得} (de, ton neutre) joue le rôle de \cle{marque du complément d'état} : il relie le verbe à une évaluation qualitative (adjectif) de l'action. En linguistique chinoise, on parle de \zh{程度补语} (complément de degré), terme que le programme officiel nomme « complément d'état ». Il ne faut pas le confondre avec le \zh{的} (attributif) ni le \zh{地} (adverbialisant), bien que les trois se prononcent \emph{de}.

\courssub{Cas 1 : Le verbe n'a pas d'objet (verbe intransitif ou employé absolument)}

\begin{propriete}[Structure de base : S + V + 得 + (Adv. degré) + Adj.]
Lorsque le verbe n'est pas suivi d'un objet, le complément d'état s'insère directement après le verbe grâce à \zh{得}.

\textbf{Schéma :} \textbf{Sujet + Verbe + 得 + (Adverbe de degré) + Adjectif}

L'adjectif décrit la \emph{manière} dont l'action est effectuée. En français, cela correspond souvent à un adverbe en \emph{-ment} (vite, clairement, sérieusement) ou à une locution (de manière appliquée).
\end{propriete}

\begin{exemplebox}[Verbe sans objet — Exemples camerounais et quotidiens]
\zh{他跑得很快。} \emph{Tā pǎo de hěn kuài.} « Il court très vite. » \\
\zh{她起得很早。} \emph{Tā qǐ de hěn zǎo.} « Elle se lève très tôt. » \\
\zh{老师讲得真清楚。} \emph{Lǎoshī jiǎng de zhēn qīngchu.} « Le professeur explique vraiment clairement. » \\
\zh{弟弟睡得不太好。} \emph{Dìdi shuì de bú tài hǎo.} « Le petit frère ne dort pas très bien. » \\
\zh{我们在雅温得住得很舒服。} \emph{Wǒmen zài Yǎwēndé zhù de hěn shūfu.} « Nous habitons très confortablement à Yaoundé. » \\
\zh{这道题我做得对。} \emph{Zhè dào tí wǒ zuò de duì.} « Cet exercice, je l'ai fait correctement. » (Adj. monosyllabique \zh{对} accepté car sens accompli/résultat).
\end{exemplebox}

\begin{attention}[Piège n°1 : L'ordre des mots — L'adjectif ne précède jamais le verbe]
En français, on dit « Il \emph{vite} court » (inversion poétique) ou « Il court \emph{vite} ». En chinois, la position après \zh{得} est \textbf{fixe et obligatoire}.
\begin{itemize}
\item \textbf{Correct :} \zh{他跑得很快。}
\item \textbf{Faux :} \zh{他很快跑得。} \quad \zh{他得很快跑。} \quad \zh{很快他跑得。}
\end{itemize}
Le complément d'état est un \emph{complément post-verbal} : il suit toujours le noyau verbal (ou sa répétition).
\end{attention}

\courssub{Cas 2 : Le verbe a un objet — La règle du verbe doublé}

En chinois, un verbe ne peut pas gouverner simultanément un \textbf{objet} et un \textbf{complément d'état} (contrainte de « unique complément »).
\begin{exemplebox}[Le conflit Objet vs Complément]
\zh{他说汉语。} \emph{Tā shuō Hànyǔ.} (Verbe + Objet : OK) \\
\zh{他说得很好。} \emph{Tā shuō de hěn hǎo.} (Verbe + Complément d'état : OK) \\
\zh{他说汉语得很好。} \emph{Tā shuō Hànyǔ de hěn hǎo.} \textbf{INCORRECT} (Verbe + Objet + Complément : INTERDIT)
\end{exemplebox}

\begin{propriete}[Structure au verbe doublé : S + V + O + V + 得 + (Adv.) + Adj.]
Pour exprimer l'objet \textbf{et} la manière, on \cle{répète le verbe} : la 1ʳᵉ occurrence prend l'objet, la 2ᵉ porte le \zh{得} + adjectif.
\end{propriete}

\begin{exemplebox}[Verbe doublé — Contexte scolaire et camerounais]
\zh{她跳舞跳得很漂亮。} \emph{Tā tiào wǔ tiào de hěn piàoliang.} « Elle danse très joliment. » \\
\zh{他写汉字写得很认真。} \emph{Tā xiě Hànzì xiě de hěn rènzhēn.} « Il écrit les caractères chinois très sérieusement. » \\
\zh{我们踢足球踢得不太快。} \emph{Wǒmen tī zúqiú tī de bú tài kuài.} « Nous jouons au foot pas très vite. » \\
\zh{小张复习功课复习得很仔细。} \emph{Xiǎo Zhāng fùxí gōngkè fùxí de hěn zǐxì.} « Petit Zhang révise ses leçons très soigneusement. » \\
\zh{妈妈做擀菜做得真好吃。} \emph{Māma zuò gǎncài zuò de zhēn hǎochī.} « Maman cuisine le \emph{ndolé} (gǎncài) vraiment délicieusement. »
\end{exemplebox}

\begin{methode}[La technique du « verbe doublé » en 3 temps]
Pour ne plus jamais faire l'erreur \zh{*V O 得 Adj} :
\begin{enumerate}
\item \textbf{Isolez le bloc Verbe + Objet :} \zh{说汉语} / \zh{写汉字} / \zh{踢足球} / \zh{做擀菜}.
\item \textbf{Construisez le bloc Évaluation :} Verbe seul + \zh{得} + Adv. + Adj. $\rightarrow$ \zh{说得很好} / \zh{写得很认真} / \zh{踢得不太快} / \zh{做得真好吃}.
\item \textbf{Collez les deux blocs :} Sujet + [Bloc 1] + [Bloc 2]. $\rightarrow$ \zh{他说汉语说得很好。}
\end{enumerate}
\emph{Astuce :} À l'oral, marquez une micro-pause après l'objet : \zh{他说汉语 | 说得很好。}
\end{methode}

\courssub{Variante élégante : L'objet avancé (Topicalisation)}

Très courant à l'oral et à l'écrit journalistique, on déplace l'objet en tête de phrase (position de \cle{thème}). Le verbe, désormais sans objet à sa droite, peut recevoir directement \zh{得}.

\begin{propriete}[Structure à objet avancé : S + O + V + 得 + (Adv.) + Adj.]
L'objet devient le sujet de la prédication (ce dont on parle). Nuance : « Quant au chinois, il le parle bien. »
\end{propriete}

\begin{exemplebox}[Objet avancé — Comparaison]
\zh{他汉语说得很好。} \emph{Tā Hànyǔ shuō de hěn hǎo.} « Le chinois, il le parle très bien. » (Focus sur \zh{汉语}) \\
\zh{这道菜我做得不太会。} \emph{Zhè dào cài wǒ zuò de bú tài huì.} « Ce plat, je ne sais pas trop le cuisiner. » \\
\zh{喀麦隆足球我们踢得很认真。} \emph{Kāmàilóng zúqiú wǒmen tī de hěn rènzhēn.} « Le football camerounais, on le joue très sérieusement. »
\end{exemplebox}

\begin{attention}[Choix entre Verbe Doublé et Objet Avancé]
\begin{itemize}
\item \textbf{Verbe doublé (S V O V 得...)} : Neutre, narratif, réponse à « Que fait-il ? ». \emph{Standard à l'écrit scolaire.}
\item \textbf{Objet avancé (S O V 得...)} : Thématise l'objet (déjà connu, mis en contraste). \emph{Très naturel à l'oral.}
\end{itemize}
Les deux sont corrects. À l'examen, les deux sont acceptés sauf consigne contraire (« Utilisez la structure à verbe doublé »).
\end{attention}

\courssub{Le rôle obligatoire de l'adverbe de degré}

\begin{propriete}[Règle du modificateur obligatoire]
Dans un complément d'état \textbf{énonciatif neutre}, l'adjectif qui suit \zh{得} est \textbf{quasiment toujours précédé d'un adverbe de degré} (\zh{很}, \zh{非常}, \zh{真}, \zh{太}, \zh{比较}, \zh{不太}, etc.).
\end{propriete}

\begin{center}
\begin{tabularx}{\linewidth}{|X|X|X|X|}
\hline
\rowcolor{popblueL} Caractères & Pinyin & Nature & Traduction / Nuance \\
\hline
很 & hěn & adv. degré & \textbf{Marque standard} (sens « très » affaibli, simple liaison) \\
\hline
非常 & fēicháng & adv. degré & extrêmement, exceptionnellement (fort) \\
\hline
真 & zhēn & adv. degré & vraiment, sincèrement (subjectif, affectif) \\
\hline
太 & tài & adv. degré & trop, tellement (souvent exclamatif, \zh{了} possible en fin de phrase) \\
\hline
比较 & bǐjiào & adv. degré & relativement, assez (modéré) \\
\hline
不太 & bú tài & adv. négatif & pas très, pas vraiment (négation standard) \\
\hline
好 & hǎo & adjectif & bien / correctement \\
\hline
快 / 慢 & kuài / màn & adjectif & vite / lentement \\
\hline
清楚 / 明确 & qīngchu / míngquè & adjectif & clairement / explicitement \\
\hline
漂亮 & piàoliang & adjectif & joliment, magnifiquement \\
\hline
流利 & liúlì & adjectif & couramment (langue) \\
\hline
认真 / 仔细 & rènzhēn / zǐxì & adjectif & sérieusement / attentivement \\
\hline
\end{tabularx}
\end{center}

\begin{attention}[Piège n°2 : Adjectif monosyllabique nu et \zh{很} seul]
\begin{itemize}
\item \textbf{Interdit :} \zh{他跑得快。} (Sauf contexte comparatif explicite : « C'est \emph{lui} qui court vite » / « Il court \emph{vite} [pas lentement] »). En énoncé neutre : \zh{他跑得\textbf{很}快。}
\item \textbf{Interdit :} \zh{他说得很。} (\zh{很} ne peut pas finir la phrase, il attend un adjectif).
\item \textbf{Correct :} \zh{得 + 很/真/不太 + Adj.}
\end{itemize}
\emph{Exception :} Certains adjectifs monosyllabiques figés dans des expressions idiomatiques ou sens résultatif : \zh{吃得饱} (manger à satiété), \zh{看得见} (voir distinctement), \zh{做得对} (faire correctement).
\end{attention}

\courssub{Schémas visuels des structures}

\begin{popfigure}
\begin{tikzpicture}[font=\small,
 bloc/.style={rectangle, rounded corners=3pt, draw=popblue, fill=popblueL, minimum height=0.9cm, align=center, inner sep=4pt},
 bloc2/.style={rectangle, rounded corners=3pt, draw=poporange, fill=poporangeL, minimum height=0.9cm, align=center, inner sep=4pt},
 fleche/.style={-{Stealth[length=2.5mm]}, thick, popdark}]
\node[bloc, align=center] (s) at (0,0){Sujet\\他};
\node[bloc, align=center] (v) at (2.6,0){Verbe\\跑};
\node[bloc2, align=center] (de) at (5.2,0){得\\de};
\node[bloc2, align=center] (adj) at (8.4,0){Adv. degré + Adj.\\很快};
\draw[fleche] (s) -- (v);
\draw[fleche] (v) -- (de);
\draw[fleche] (de) -- (adj);
\node[bloc, align=center] (s2) at (0,-2){Sujet\\他};
\node[bloc, align=center] (v2) at (2.2,-2){Verbe 1\\说};
\node[bloc, align=center] (o2) at (4.2,-2){Objet\\汉语};
\node[bloc, align=center] (v3) at (6.2,-2){Verbe 2 (répété)\\说};
\node[bloc2, align=center] (de2) at (8.2,-2){得\\de};
\node[bloc2, align=center] (adj2) at (10.6,-2){Adv. + Adj.\\很好};
\draw[fleche] (s2) -- (v2);
\draw[fleche] (v2) -- (o2);
\draw[fleche] (o2) -- (v3);
\draw[fleche] (v3) -- (de2);
\draw[fleche] (de2) -- (adj2);
\draw[fleche, poppurple, dashed, bend left=35] (v2.north) to node[above, poppurple, font=\footnotesize]{Répétition obligatoire} (v3.north);
\end{tikzpicture}
\legende{En haut : Cas 1 (Verbe intransitif). En bas : Cas 2 (Verbe transitif $\rightarrow$ Verbe doublé).}
\end{popfigure}

\begin{popfigure}
\begin{tikzpicture}[font=\small,
 decision/.style={diamond, draw=poppurple, fill=poppurpleL, aspect=2.2, align=center, inner sep=2pt},
 action/.style={rectangle, rounded corners=3pt, draw=popgreen, fill=popgreenL, align=center, inner sep=5pt},
 fleche/.style={-{Stealth[length=2.5mm]}, thick, popdark}]
\node[decision, align=center] (q) at (0,0){Le verbe a-t-il\\un objet \textbf{exprimé} ?};
\node[action, align=center] (non) at (-4.5,-2.8){NON\\S + V + 得 + Adv. + Adj.\\ \zh{他跑得很快。}};
\node[action, align=center] (oui) at (4.5,-2.8){OUI\\Choix 1 : Verbe doublé\\S + V + O + V + 得 + Adv. + Adj.\\ \zh{他说汉语说得很好。}\\
Choix 2 : Objet avancé\\S + O + V + 得 + Adv. + Adj.\\ \zh{他汉语说得很好。}};
\draw[fleche] (q) -- (non) node[midway, above left, font=\footnotesize]{Non};
\draw[fleche] (q) -- (oui) node[midway, above right, font=\footnotesize]{Oui};
\end{tikzpicture}
\legende{Organigramme de décision pour construire le complément d'état.}
\end{popfigure}

\coursec{Négation, question et emplois particuliers}

\courssub{La négation : \zh{不} / \zh{没} devant \zh{得}}

\begin{propriete}[Négation du complément d'état]
Le négateur (\zh{不} pour état habituel/futur, \zh{没} pour passé/accompli) se place \textbf{immédiatement avant \zh{得}}, jamais devant le verbe principal.
\textbf{Schéma :} S + V (+ O + V) + \textbf{不/没} + 得 + Adv. + Adj.
\end{propriete}

\begin{exemplebox}[Négation]
\zh{他跑得\textbf{不}快。} \emph{Tā pǎo de \textbf{bú} kuài.} « Il ne court pas vite. » (Habitude) \\
\zh{他说汉语说得\textbf{不}太好。} \emph{Tā shuō Hànyǔ shuō de \textbf{bú} tài hǎo.} « Il ne parle pas très bien chinois. » \\
\zh{昨天我睡得\textbf{没}好。} \emph{Zuótiān wǒ shuì de \textbf{méi} hǎo.} « Hier, je n'ai pas bien dormi. » (Passé accompli $\rightarrow$ \zh{没}) \\
\zh{这道题他做得\textbf{不}对。} \emph{Zhè dào tí tā zuò de \textbf{bú} duì.} « Il a fait cet exercice incorrectement. »
\end{exemplebox}

\begin{attention}[Erreur fréquente : *\zh{不跑得快} / *\zh{没说得好}]
Le négateur ne nie pas l'action (il \emph{a} couru, il \emph{a} parlé), il nie la \emph{manière}. Donc \zh{不/没} se colle à \zh{得}.
\begin{itemize}
\item \zh{他不跑。} = Il ne court pas (action niée).
\item \zh{他跑得不快。} = Il court, mais pas vite (manière niée).
\end{itemize}
\end{attention}

\courssub{L'interrogation : \zh{怎么样} / \zh{得 + Adj + 吗} / \zh{得 + 不 + Adj}}

\begin{propriete}[Trois façons de demander « Comment ? »]
\begin{enumerate}
\item \textbf{S + V (+ O + V) + 得 + 怎么样 ?} (Ouvert, standard).
\item \textbf{S + V (+ O + V) + 得 + Adj + 吗 ?} (Fermé, confirmation).
\item \textbf{S + V (+ O + V) + 得 + 不 + Adj ?} (Choix binaire, attendue positive).
\end{enumerate}
\end{propriete}

\begin{exemplebox}[Questions sur la manière]
\zh{你复习得怎么样？} \emph{Nǐ fùxí de zěnmeyàng?} « Comment as-tu révisé ? » \\
\zh{他法语说得流利吗？} \emph{Tā Fǎyǔ shuō de liúlì ma?} « Est-ce qu'il parle couramment français ? » (Objet avancé) \\
\zh{这道菜你做得不做得好吃？} \emph{Zhè dào cài nǐ zuò de bù zuò de hǎochī?} « Ce plat, est-ce que tu le cuisines bien ou pas ? » (Forme A-pas-A sur le complément, rare mais possible) \\
\zh{踢球踢得累不累？} \emph{Tī qiú tī de lèi bu lèi?} « Jouer au foot, est-ce que ça fatigue ? » (Verbe doublé + A-pas-A sur l'adj).
\end{exemplebox}

\courssub{Emploi aspectuel : \zh{得 + Adj + 了} (Changement d'état)}

\begin{propriete}[Complément d'état + \zh{了} = Nouvel état / Changement]
L'ajout de la particule \zh{了} en fin de phrase indique que le degré atteint est \textbf{nouveau}, \textbf{surpris} ou \textbf{abouti}.
\end{propriete}

\begin{exemplebox}[Avec \zh{了}]
\zh{现在他汉语说得很好\textbf{了}。} \emph{Xiànzài tā Hànyǔ shuō de hěn hǎo \textbf{le}.} « Maintenant, il parle très bien chinois (alors qu'avant non). » \\
\zh{你跑得真快\textbf{了}！} \emph{Nǐ pǎo de zhēn kuài \textbf{le}!} « Tu cours vraiment vite maintenant ! (Progrès constaté). »
\end{exemplebox}

\coursec{Comparaison avec le français et pièges majeurs}

\begin{center}
\begin{tabularx}{\linewidth}{|X|X|X|}
\hline
\rowcolor{popblueL} Structure & Exemple en caractères + pinyin & Traduction \\
\hline
S + V + 得 + Adv. + Adj. & 他跑得很快。Tā pǎo de hěn kuài. & Il court très vite. \\
\hline
S + V + O + V + 得 + Adv. + Adj. & 他说汉语说得很好。Tā shuō Hànyǔ shuō de hěn hǎo. & Il parle très bien le chinois. \\
\hline
S + O + V + 得 + Adv. + Adj. & 他汉语说得很好。Tā Hànyǔ shuō de hěn hǎo. & Le chinois, il le parle très bien. \\
\hline
Négation & 他跑得不快。Tā pǎo de bú kuài. & Il ne court pas vite. \\
\hline
Question & 你睡得怎么样？Nǐ shuì de zěnmeyàng? & Comment as-tu dormi ? \\
\hline
\end{tabularx}
\end{center}

\begin{center}
\begin{tabularx}{\linewidth}{|X|X|}
\hline
\rowcolor{popblueL} Français (Piège) & Chinois (Réalité) \\
\hline
Un seul verbe porte adverbe + objet : « Il parle \emph{bien le chinois} ». & \textbf{Deux verbes} (ou objet avancé) : \zh{说汉语说得好} / \zh{汉语说得好}. \\
\hline
Adverbes en \emph{-ment} : bien, vite, clairement. & \textbf{Adjectifs} + modificateur : \zh{很好}, \zh{很快}, \zh{很清楚}. \\
\hline
Négation : « Il ne parle \emph{pas} bien. » & Négation sur \zh{得} : \zh{说得\textbf{不}好} (pas \zh{不说得好}). \\
\hline
« Il \emph{bien} parle » (inversion) impossible. & Ordre figé : V + \zh{得} + Adj. Jamais Adj + V. \\
\hline
\end{tabularx}
\end{center}

\begin{attention}[Les 3 erreurs fatales des francophones (Check-list avant de rendre la copie)]
\begin{enumerate}
\item \textbf{La fusion interdite :} \zh{*他说汉语得很好。} $\rightarrow$ \textbf{Corrigez :} \zh{他说汉语说得很好。} ou \zh{他汉语说得很好。}
\item \textbf{La négation déplacée :} \zh{*他不说得好。} / \zh{*他说不得好。} $\rightarrow$ \textbf{Corrigez :} \zh{他说得不好。} (\zh{不} colle à \zh{得}).
\item \textbf{L'adjectif nu :} \zh{*我写得快。} (Énoncé neutre) $\rightarrow$ \textbf{Corrigez :} \zh{我写得\textbf{很}快。} (Sauf comparatif : \zh{我写得快，不写得慢。}).
\end{enumerate}
\end{attention}

\begin{savaistu}[Le \zh{得} des trois \emph{de} : \zh{的・得・地}]
En chinois, trois particules homophones (ton neutre) structurent la phrase :
\begin{itemize}
\item \zh{\textbf{的}} (radical \zh{白}) : \textbf{Attributif} (N + \zh{的} + N). \emph{Ex:} \zh{我的书} (mon livre), \zh{红色的花} (fleur rouge).
\item \zh{\textbf{得}} (radical \zh{彳} double personne) : \textbf{Complément d'état/degré} (V + \zh{得} + Adj). \emph{Ex:} \zh{跑得快}. Mémo : « Il faut \textbf{marcher (彳)} pour \textbf{atteindre} le résultat. »
\item \zh{\textbf{地}} (radical \zh{土}) : \textbf{Adverbialisant} (Adj + \zh{地} + V). \emph{Ex:} \zh{快快地跑} (courir vite). Mémo : « Cela se passe sur le \textbf{sol (土)} / terrain de l'action. »
\end{itemize}
\emph{Attention :} À l'oral, tout est \emph{de}. À l'écrit (examen, HSK), la distinction est \textbf{obligatoire} et notée.
\end{savaistu}

\coursec{Exercices d'application et de transformation}

\begin{exoresolu}[Transformation : Verbe + Objet $\rightarrow$ Complément d'état (Verbe doublé)]
Énoncé : Transformez les phrases « Sujet + Verbe + Objet » en phrases avec complément d'état en utilisant l'adjectif indiqué. Utilisez la structure à verbe doublé.
\begin{enumerate}
\item \zh{他 学 中文} (adjectif : \zh{认真} rènzhēn)
\item \zh{我们 吃 擀菜} (adjectif : \zh{香} xiāng — parfumé/délicieux)
\item \zh{妹妹 唱 歌} (adjectif : \zh{好听} hǎotīng — agréable à entendre)
\end{enumerate}
\tcblower
\textbf{Solution détaillée :}
\begin{enumerate}
\item \textbf{Étape 1 :} Verbe + Objet = \zh{学中文}. \textbf{Étape 2 :} Verbe répété + \zh{得} + \zh{很} + Adj = \zh{学得很认真}. \textbf{Résultat :} \zh{他学中文学得很认真。} \emph{Tā xué Zhōngwén xué de hěn rènzhēn.} « Il étudie le chinois très sérieusement. »
\item \textbf{Étape 1 :} \zh{吃擀菜}. \textbf{Étape 2 :} \zh{吃得很香}. \textbf{Résultat :} \zh{我们吃擀菜吃得很香。} \emph{Wǒmen chī gǎncài chī de hěn xiāng.} « Nous mangeons le ndolé avec grand appétit / c'est très bon. »
\item \textbf{Étape 1 :} \zh{唱歌}. \textbf{Étape 2 :} \zh{唱得很好听}. \textbf{Résultat :} \zh{妹妹唱歌唱得很好听。} \emph{Mèimei chàng gē chàng de hěn hǎotīng.} « Ma petite sœur chante très agréablement. »
\end{enumerate}
\end{exoresolu}

\begin{exoresolu}[Correction d'erreurs (Diagnostic)]
Énoncé : Chaque phrase contient une erreur typique. Identifiez-la, expliquez-la et corrigez.
\begin{enumerate}
\item \zh{她做饭做得真好吃了。} (Contexte : Elle cuisine d'habitude bien.)
\item \zh{他们不踢球踢得好。}
\item \zh{老师讲课讲得清楚不清楚？}
\item \zh{我法语学得很。}
\end{enumerate}
\tcblower
\textbf{Corrigé :}
\begin{enumerate}
\item \textbf{Erreur :} \zh{了} en fin de phrase marque un changement d'état / nouvel état. Si c'est une habitude, \zh{了} est en trop. \textbf{Correction :} \zh{她做饭做得真好吃。} (Ou \zh{做得真好吃了} si on veut dire « Maintenant elle cuisine bien [alors qu'avant non] »).
\item \textbf{Erreur :} Négation \zh{不} placée devant le verbe principal (\zh{踢球}) au lieu de devant \zh{得}. \textbf{Correction :} \zh{他们踢球踢得不好。} ou \zh{他们踢球踢得不太好。}
\item \textbf{Erreur :} Forme interrogative \zh{A不A} (\zh{清楚不清楚}) mal placée. On ne met pas \zh{不} dans l'adjectif du complément d'état pour une question standard. \textbf{Correction :} \zh{老师讲课讲得怎么样？} ou \zh{老师讲课讲得清楚吗？}
\item \textbf{Erreur :} \zh{很} utilisé seul sans adjectif. \textbf{Correction :} \zh{我法语学得很好。} / \zh{我法语学得很认真。}
\end{enumerate}
\end{exoresolu}

\begin{exercice}[Entraînement guidé : À vous de jouer !]
\textbf{A. Complétez avec la structure correcte (Verbe doublé ou Objet avancé) + \zh{得} + Adv. + Adj.}
\begin{enumerate}
\item 我 \_\_\_\_\_\_ 汉字 \_\_\_\_\_\_ \_\_\_\_\_\_ 很认真。 (écrire / 写)
\item 他们 \_\_\_\_\_\_ 足球 \_\_\_\_\_\_ \_\_\_\_\_\_ 不太好。 (jouer au foot / 踢足球)
\item 这道菜 \_\_\_\_\_\_ 我 \_\_\_\_\_\_ \_\_\_\_\_\_ 真好吃。 (cuisiner / 做) \emph{Indice : Objet avancé.}
\item 小李 \_\_\_\_\_\_ 法语 \_\_\_\_\_\_ \_\_\_\_\_\_ 非常流利。 (parler / 说)
\end{enumerate}

\textbf{B. Mettez à la forme négative (habitude) :}
\begin{enumerate}
\item \zh{他跑得很快。} $\rightarrow$ \_\_\_\_\_\_
\item \zh{她唱歌唱得很好听。} $\rightarrow$ \_\_\_\_\_\_
\item \zh{我们住得很舒服。} $\rightarrow$ \_\_\_\_\_\_
\end{enumerate}

\textbf{C. Traduisez en chinois (caractères + pinyin) :}
\begin{enumerate}
\item « Mon professeur de chinois explique très clairement. »
\item « Est-ce que tu dors bien la nuit ? » (nuit = \zh{晚上} wǎnshang)
\item « Le football, nous y jouons assez sérieusement. » (Objet avancé + \zh{比较})
\item « Hier, il n'a pas bien mangé. » (Passé $\rightarrow$ \zh{没})
\end{enumerate}
\end{exercice}

\begin{corrige}[Exercice : Entraînement guidé]
\textbf{A. Complétion :}
\begin{enumerate}
\item \zh{我写汉字写得很认真。} (Wǒ xiě Hànzì xiě de hěn rènzhēn.) — Verbe doublé.
\item \zh{他们踢足球踢得不太好。} (Tāmen tī zúqiú tī de bú tài hǎo.) — Verbe doublé.
\item \zh{这道菜我做得真好吃。} (Zhè dào cài wǒ zuò de zhēn hǎochī.) — Objet avancé.
\item \zh{小李说法语说得非常流利。} (Xiǎo Lǐ shuō Fǎyǔ shuō de fēicháng liúlì.) — Verbe doublé.
\end{enumerate}

\textbf{B. Négation (habitude \zh{不}) :}
\begin{enumerate}
\item \zh{他跑得不快。} (Tā pǎo de bú kuài.)
\item \zh{她唱歌唱得不好听。} (Tā chàng gē chàng de bù hǎotīng.) — \zh{好听} est un adjectif composé, \zh{不} devant l'adjectif complet.
\item \zh{我们住得不舒服。} (Wǒmen zhù de bù shūfu.)
\end{enumerate}

\textbf{C. Traduction :}
\begin{enumerate}
\item \zh{我的中文老师讲得很清楚。} (Wǒ de Zhōngwén lǎoshī jiǎng de hěn qīngchu.) — Verbe \zh{讲} sans objet exprimé ici.
\item \zh{你晚上睡得好吗？} (Nǐ wǎnshang shuì de hǎo ma?) — \zh{睡} intransitif ici. \zh{好} suffit comme adjectif court dans la question \zh{...好吗?}.
\item \zh{足球我们踢得比较认真。} (Zúqiú wǒmen tī de bǐjiào rènzhēn.) — Objet avancé + \zh{比较}.
\item \zh{他昨天吃得没好。} (Tā zuótiān chī de méi hǎo.) — Passé accompli $\rightarrow$ \zh{没}. Verbe \zh{吃} sans objet précis (manger en général). Si objet \zh{饭} : \zh{他昨天吃饭吃得没好。}
\end{enumerate}
\end{corrige}

\begin{aretenir}[Synthèse de la Leçon 32 — Complément d'état (程度补语)]
\begin{itemize}
\item \textbf{Fonction :} Évaluer le degré / la manière de réalisation d'une action (Verbe + \zh{得} + Adj).
\item \textbf{Structure 1 (Sans objet) :} S + V + \zh{得} + Adv. + Adj. \quad \zh{他跑得很快。}
\item \textbf{Structure 2 (Avec objet — Verbe doublé) :} S + V + O + V + \zh{得} + Adv. + Adj. \quad \zh{他说汉语说得很好。}
\item \textbf{Structure 3 (Avec objet — Objet avancé) :} S + O + V + \zh{得} + Adv. + Adj. \quad \zh{他汉语说得很好。}
\item \textbf{Modificateur obligatoire :} Adv. de degré (\zh{很}, \zh{真}, \zh{不太}, \zh{非常}...) devant l'adjectif. Pas d'adj. monosyllabique nu (sauf cas figés).
\item \textbf{Négation :} \zh{不/没} + \zh{得} + Adv. + Adj. \quad \zh{说得不好。}
\item \textbf{Question :} \zh{得 + 怎么样?} / \zh{得 + Adj + 吗?}
\item \textbf{Distinction graphique :} \zh{的} (N), \zh{\textbf{得}} (V), \zh{地} (Adv). \textbf{Ne les confondez pas à l'écrit !}
\end{itemize}
\end{aretenir}

\coursec{Négation, question et emplois}

Dans la section précédente, nous avons découvert le schéma de base du complément d'état : \textbf{Sujet + Verbe + 得 + adjectif}, qui permet d'évaluer la manière dont une action est accomplie (\zh{他说得很好} — il parle bien). Mais comment dire « il ne parle pas bien » ? Comment demander « parles-tu bien ? » ou « comment chantes-tu ? » ? C'est précisément l'objet de cette section : la \cle{négation}, les \cle{questions} et les \cle{emplois} du complément d'état.

\courssub{Observation : un dialogue en classe de chinois}

\begin{textebox}[À la fin du cours de chinois, au lycée de Yaoundé]
A : 你中文学得怎么样？\\
\emph{Nǐ Zhōngwén xué de zěnmeyàng ?}\\
« Comment se passe ton apprentissage du chinois ? »\\[0.4em]
B : 我学得还不错，就是汉字写得不太好。\\
\emph{Wǒ xué de hái búcuò, jiùshì Hànzì xiě de bú tài hǎo.}\\
« J'apprends plutôt bien, seulement je n'écris pas très bien les caractères. »\\[0.4em]
A : 你听得懂老师的话吗？\\
\emph{Nǐ tīng de dǒng lǎoshī de huà ma ?}\\
« Comprends-tu ce que dit le professeur quand tu écoutes ? »\\[0.4em]
B : 听得懂，可是说得不太流利。\\
\emph{Tīng de dǒng, kěshì shuō de bú tài liúlì.}\\
« Je comprends en écoutant, mais je ne parle pas très couramment. »
\end{textebox}

\textbf{Observons.} Repérez dans le dialogue les formes suivantes :
\begin{itemize}
\item \zh{学得怎么样} : question avec \zh{怎么样} (comment) ;
\item \zh{写得不太好} : négation avec \zh{不} placé \textbf{devant l'adjectif} ;
\item \zh{听得懂吗} : question avec \zh{吗} (structure de complément de possibilité, voir encadré ci-dessous) ;
\item \zh{说得不太流利} : encore une négation, \zh{不} + adjectif.
\end{itemize}

Dans tous les cas, la structure de base \zh{V + 得 + adjectif} reste intacte : c'est autour d'elle que viennent se greffer la négation et les marques de question. Analysons chaque cas.

\begin{attention}[Distinction importante : Complément d'état vs Complément de possibilité]
Le dialogue contient \zh{听得懂} (entendre et comprendre). \zh{懂} (comprendre) est un \textbf{verbe}, pas un adjectif. La structure \zh{Verbe + 得/不 + Verbe} (ex. \zh{听得懂}, \zh{看不见}, \zh{吃得完}) est le \textbf{complément de possibilité} (HSK 3), qui exprime la \emph{capacité} d'atteindre un résultat.
Le \textbf{complément d'état} (cette leçon) suit la structure \zh{Verbe + 得 + Adjectif} (ex. \zh{写得好}, \zh{跑得快}) et évalue la \emph{manière}.
La négation et l'interrogation fonctionnent de façon très similaire pour les deux, mais ne les confondez pas : \cle{État = V + 得 + Adj} \quad vs \quad \cle{Possibilité = V + 得/不 + V}.
\end{attention}

\courssub{La négation : V + 得 + 不 + adjectif}

\begin{propriete}[Règle de la négation du complément d'état]
La négation du complément d'état se construit ainsi :
\begin{center}
\textbf{Sujet + Verbe + 得 + 不 + adjectif}
\end{center}
Le mot négatif \zh{不} se place \textbf{devant l'adjectif}, jamais devant le verbe. C'est l'\emph{évaluation} (bien, mal, vite, clairement) qui est niée, pas l'action elle-même.
\end{propriete}

\begin{exemplebox}[Exemples de négation]
\zh{他写得不好。} \emph{Tā xiě de bù hǎo.} « Il n'écrit pas bien. »\\[0.3em]
\zh{他睡得不好。} \emph{Tā shuì de bù hǎo.} « Il dort mal. »\\[0.3em]
\zh{我跑得不太快。} \emph{Wǒ pǎo de bú tài kuài.} « Je ne cours pas très vite. »\\[0.3em]
\zh{她汉语说得不好。} \emph{Tā Hànyǔ shuō de bù hǎo.} « Elle ne parle pas bien chinois. »\\[0.3em]
\zh{弟弟今天吃得不多。} \emph{Dìdi jīntiān chī de bù duō.} « Mon petit frère n'a pas beaucoup mangé aujourd'hui. »
\end{exemplebox}

\begin{attention}[Piège majeur : la place de 不]
Un francophone veut traduire mot à mot « il ne parle pas bien » par *\zh{他不说话得很好} ou *\zh{他不写得好}. \textbf{C'est faux !} En chinois, \zh{不} ne peut pas se placer devant le verbe dans cette structure : \zh{不} modifie l'\textbf{adjectif}, pas le verbe. On dit : \zh{他说得不好} (littéralement : « il parle de-manière-pas-bonne »). Retenez : la négation porte sur l'\textbf{évaluation}, pas sur l'action.
\end{attention}

\begin{popfigure}
\begin{center}
\begin{tikzpicture}[font=\small]
% Schéma correct
\node[draw=popgreen, fill=popgreenL, rounded corners, align=center, line width=1pt] (ok) at (0,0) {Sujet + \textbf{V} + 得 + \textbf{不} + adjectif};
\node[below=0.15cm of ok, text=popgreen] {\zh{他写得不好。} — correct};
% Schéma faux
\node[draw=poppink, fill=poppinkL, rounded corners, align=center, line width=1pt] (ko) at (0,-1.8) {Sujet + \textbf{不} + V + 得 + adjectif};
\node[below=0.15cm of ko, text=poppink] {*\zh{他不写得好。} — FAUX};
% Croix sur le schéma faux
\draw[poppink, line width=2pt] ($(ko.north west)+(-0.15,0.15)$) -- ($(ko.south east)+(0.15,-0.15)$);
\draw[poppink, line width=2pt] ($(ko.north east)+(0.15,0.15)$) -- ($(ko.south west)+(-0.15,-0.15)$);
% Flèche 不 vers adjectif
\draw[->, popblue, line width=1.2pt] (0,0.6) .. controls (1.5,1.3) .. (2.6,0.35) node[midway, above, text=popblue, align=center] {不 nie\\ l'adjectif};
\end{tikzpicture}
\legende{Position de 不 : devant l'adjectif (correct), jamais devant le verbe (faux).}
\end{center}
\end{popfigure}

\courssub{Les trois types de questions}

\textbf{1. La question fermée avec 吗 (oui / non).}

On ajoute simplement \zh{吗} à la fin de la phrase affirmative, comme pour toute phrase déclarative :

\begin{exemplebox}[Questions avec 吗]
\zh{你说得清楚吗？} \emph{Nǐ shuō de qīngchu ma ?} « Parles-tu clairement ? »\\[0.3em]
\zh{他汉语说得好吗？} \emph{Tā Hànyǔ shuō de hǎo ma ?} « Parle-t-il bien chinois ? »\\[0.3em]
\zh{老师说得快吗？} \emph{Lǎoshī shuō de kuài ma ?} « Le professeur parle-t-il vite ? »
\end{exemplebox}

\textbf{2. La question ouverte avec 怎么样 (comment ?).}

C'est la question la plus naturelle pour demander une \textbf{évaluation}. Deux constructions sont possibles : \zh{V + 得 + 怎么样} ou \zh{Nom + V + 得 + 怎么样} quand on précise l'objet.

\begin{exemplebox}[Questions avec 怎么样]
\zh{你汉语说得怎么样？} \emph{Nǐ Hànyǔ shuō de zěnmeyàng ?} « Comment parles-tu chinois ? »\\[0.3em]
\zh{他歌唱得怎么样？} \emph{Tā gē chàng de zěnmeyàng ?} « Comment chante-t-il ? »\\[0.3em]
\zh{你最近睡得怎么样？} \emph{Nǐ zuìjìn shuì de zěnmeyàng ?} « Comment dors-tu ces jours-ci ? »
\end{exemplebox}

\textbf{3. La question alternative : V + 得 + adj + 不 + adj.}

C'est l'équivalent chinois de « vite ou pas vite ? » : on juxtapose l'adjectif et sa négation. L'adjectif est généralement monosyllabique ou bisyllabique.

\begin{exemplebox}[Questions alternatives]
\zh{他跑得快不快？} \emph{Tā pǎo de kuài bù kuài ?} « Court-il vite ? »\\[0.3em]
\zh{你写得好不好？} \emph{Nǐ xiě de hǎo bù hǎo ?} « Écris-tu bien ? »\\[0.3em]
\zh{她说得清楚不清楚？} \emph{Tā shuō de qīngchu bù qīngchu ?} « Parle-t-elle clairement ? »
\end{exemplebox}

\begin{aretenir}[Tableau synoptique des cinq formes]
\begin{center}
\begin{tabularx}{\linewidth}{|l|X|X|}
\hline
\rowcolor{popblueL} \textbf{Forme} & \textbf{Exemple (caractères + pinyin)} & \textbf{Traduction} \\
\hline
Affirmative & \zh{他说得好。} \emph{Tā shuō de hǎo.} & Il parle bien. \\
\hline
Négative & \zh{他说得不好。} \emph{Tā shuō de bù hǎo.} & Il ne parle pas bien. \\
\hline
Question 吗 & \zh{他说得好吗？} \emph{Tā shuō de hǎo ma ?} & Parle-t-il bien ? \\
\hline
Question 怎么样 & \zh{他说得怎么样？} \emph{Tā shuō de zěnmeyàng ?} & Comment parle-t-il ? \\
\hline
Question alternative & \zh{他说得好不好？} \emph{Tā shuō de hǎo bù hǎo ?} & Parle-t-il bien (ou non) ? \\
\hline
\end{tabularx}
\end{center}
\end{aretenir}

\courssub{Les emplois du complément d'état}

\begin{propriete}[Emploi principal : habitude ou capacité générale]
Le complément d'état sert à \textbf{évaluer une habitude, une capacité générale ou une manière habituelle d'agir} : \zh{他唱得好} (il chante bien, en général), \zh{她字写得漂亮} (elle écrit joliment). Il ne décrit \textbf{pas} une action ponctuelle et unique avec un aspect accompli : pour dire « hier, il a bien couru (cette fois-là) », le chinois utilise un complément de temps en tête de phrase \zh{昨天他跑得很快}, mais le verbe ne prend pas \zh{了}.
\end{propriete}

\begin{attention}[Incompatibilité avec 了 (aspect accompli)]
Le complément d'état avec \zh{得} évalue une manière générale ou constante ; le verbe ne peut donc pas porter \zh{了} (marque de l'accompli). On dit \zh{他跑得很快} (il court vite), jamais *\zh{他跑了得很快}. Pour raconter un événement passé précis, on place le complément de temps au début : \zh{昨天他跑得很快} (Hier, il a couru vite).
\end{attention}

\begin{propriete}[Emploi étendu : insister avec 得很 / 极了]
Pour renforcer l'évaluation (utile à l'examen), on peut ajouter \zh{得很} après l'adjectif ou utiliser la structure \zh{adjectif + 极了} :
\begin{itemize}
\item \zh{好得很} \emph{hǎo de hěn} : vraiment excellent / très bon ;
\item \zh{他说得好极了。} \emph{Tā shuō de hǎo jí le.} « Il parle extrêmement bien. »
\item \zh{今天热得很。} \emph{Jīntiān rè de hěn.} « Il fait vraiment très chaud aujourd'hui. »
\end{itemize}
\end{propriete}

\begin{savaistu}[La modestie chinoise : 还可以]
\zh{怎么样} (comment ?) est la question la plus fréquente pour demander une évaluation : du travail, d'un repas, d'un match de football. La réponse typique des Chinois est \zh{还可以} \emph{hái kěyǐ} — « ça va, passable » — même quand ils excellent ! Cette modestie est une valeur culturelle importante en Chine : se vanter est mal vu. Au Cameroun, on répondra plus volontiers avec enthousiasme ; en Chine, un élève excellent dira souvent simplement \zh{还行} \emph{hái xíng} (« c'est acceptable »).
\end{savaistu}

\begin{definition}[Vocabulaire clé de la section]
\begin{center}
\begin{tabularx}{\linewidth}{|X|X|X|X|}
\hline
\rowcolor{popblueL} \textbf{Caractères} & \textbf{Pinyin} & \textbf{Nature} & \textbf{Traduction} \\
\hline
怎么样 & zěnmeyàng & pron. interrogatif & comment, quel genre \\
\hline
不太 & bú tài & adv. & pas très (négation atténuée) \\
\hline
流利 & liúlì & adj. & fluent, couramment \\
\hline
清楚 & qīngchu & adj. & clair, distinct \\
\hline
极了 & jí le & structure & extrêmement, au plus haut point \\
\hline
得很 & de hěn & structure & vraiment très (renforcement) \\
\hline
还可以 & hái kěyǐ & expr. & ça va, passable, pas mal \\
\hline
听得懂 & tīng de dǒng & compl. possibilité & comprendre en écoutant \\
\hline
\end{tabularx}
\end{center}
\end{definition}

\courssub{Exercice résolu et entraînement}

\begin{exoresolu}[Transformation de phrases]
\textbf{Énoncé.} Transformez chaque phrase affirmative en (a) phrase négative et (b) question avec 吗.\\
1. \zh{他写得好。} 2. \zh{她唱得不错。} 3. \zh{你跑得很快。}
\tcblower
\textbf{Solution détaillée.}\\
1a. \zh{他写得不好。} \emph{Tā xiě de bù hǎo.} « Il n'écrit pas bien. » — \zh{不} se place devant l'adjectif \zh{好}.\\
1b. \zh{他写得好吗？} \emph{Tā xiě de hǎo ma ?} « Écrit-il bien ? »\\
2a. \zh{她唱得不太好。} \emph{Tā chàng de bú tài hǎo.} « Elle ne chante pas très bien. » (on évite \zh{不不错}, on remplace par \zh{不太好}).\\
2b. \zh{她唱得不错吗？} \emph{Tā chàng de búcuò ma ?} « Chante-t-elle bien ? »\\
3a. \zh{你跑得不快。} \emph{Nǐ pǎo de bù kuài.} « Tu ne cours pas vite. »\\
3b. \zh{你跑得快吗？} \emph{Nǐ pǎo de kuài ma ?} « Cours-tu vite ? »
\end{exoresolu}

\begin{exercice}[À vous de jouer]
1. Traduisez en chinois : « Il ne dort pas bien. » / « Comprends-tu en écoutant ? » / « Comment écrit-elle les caractères ? » / « Court-il vite ? » (question alternative).\\
2. Corrigez la phrase fautive : *\zh{他不写得好。}\\
3. Répondez avec modestie (comme en Chine !) à la question : \zh{你汉语说得怎么样？}
\end{exercice}

\begin{corrige}[Exercice « À vous de jouer »]
1. \zh{他睡得不好。} \emph{Tā shuì de bù hǎo.} / \zh{你听得懂吗？} \emph{Nǐ tīng de dǒng ma ?} / \zh{她汉字写得怎么样？} \emph{Tā Hànzì xiě de zěnmeyàng ?} / \zh{他跑得快不快？} \emph{Tā pǎo de kuài bù kuài ?}\\
2. \zh{他写得不好。} — \zh{不} doit précéder l'adjectif \zh{好}, pas le verbe \zh{写}.\\
3. Réponse modeste attendue : \zh{还可以。} \emph{Hái kěyǐ.} « Ça va, passable. » ou \zh{我说得不太好。} \emph{Wǒ shuō de bú tài hǎo.} « Je ne parle pas très bien. »
\end{corrige}

\begin{aretenir}[L'essentiel de la section]
\begin{itemize}
\item Négation : \textbf{V + 得 + 不 + adjectif} — \zh{不} nie l'adjectif, jamais le verbe.
\item Question fermée : phrase + \zh{吗} ; question ouverte : \zh{V + 得 + 怎么样} ; question alternative : \zh{V + 得 + adj + 不 + adj}.
\item Le complément d'état évalue une \textbf{habitude ou capacité générale}, pas une action ponctuelle accomplie ; il est incompatible avec \zh{了} après le verbe.
\item Pour insister : \zh{得很} ou \zh{极了} (\zh{好得很} / \zh{好极了} — vraiment excellent).
\item Distinction : \textbf{État} = V + 得 + \textbf{Adj} \quad vs \quad \textbf{Possibilité} = V + 得/不 + \textbf{Verbe}.
\end{itemize}
\end{aretenir}

\coursec{Leçon 32 — Grammaire : Complément d'état}

\courssub{Observation et découverte}

\begin{textebox}[Dialogue : Évaluation de stage à Douala]
\textbf{Contexte :} \emph{M. Ngono, tuteur de stage à la SONARA (Limbe), fait le point avec son stagiaire, Paul, étudiant en 1ère A au Lycée de Buea.}\\[6pt]
\zh{Ngono: 保罗，你这个月表现怎么样？} \emph{Bǎoluó, nǐ zhège yuè biǎoxiàn zěnmeyàng?}\\
\zh{Paul:  我觉得我学得很快，但是汉语说得不太好。} \emph{Wǒ juéde wǒ xué de hěn kuài, dànshì Hànyǔ shuō de bù tài hǎo.}\\
\zh{Ngono:  没关系。你听得很仔细，字也写得工整。继续努力！} \emph{Méi guānxi. Nǐ tīng de hěn zǐxì, zì yě xiě de gōngzhěng. Jìxù nǔlì!}\\
\zh{Paul:  谢谢老师。我会认真地听，认真地写。} \emph{Xièxie lǎoshī. Wǒ huì rènzhēn de tīng, rènzhēn de xiě.}\\[6pt]
\textbf{Traduction :}\\
\textbf{Ngono :} Paul, comment s'est passé ton mois ? (Quelle a été ta performance ?)\\
\textbf{Paul :} Je trouve que j'apprends vite, mais je ne parle pas très bien le chinois.\\
\textbf{Ngono :} Ce n'est pas grave. Tu écoutes très attentivement, tu écris aussi les caractères très soigneusement. Continue tes efforts !\\
\textbf{Paul :} Merci monsieur. J'écouterai et j'écrirai sérieusement.
\end{textebox}

\begin{aretenir}[Mots nouveaux du dialogue]
\begin{center}
\begin{tabularx}{\linewidth}{|X|X|X|X|}
\hline
\rowcolor{popblueL} Caractères & Pinyin & Nature & Traduction \\
\hline
表现 & biǎoxiàn & n./v. & performance, comportement ; se comporter, performer \\
\hline
觉得 & juéde & v. & penser, trouver, avoir l'impression \\
\hline
学 & xué & v. & apprendre, étudier \\
\hline
快 & kuài & adj. & vite, rapide \\
\hline
关系 & guānxi & n. & relation ; \zh{没关系} : ça ne fait rien, pas de problème \\
\hline
仔细 & zǐxì & adj. & attentif, soigneux, minutieux \\
\hline
字 & zì & n. & caractère, mot, lettre \\
\hline
工整 & gōngzhěng & adj. & soigné, régulier (écriture, rangement) \\
\hline
继续 & jìxù & v. & continuer \\
\hline
努力 & nǔlì & v./adj. & faire des efforts, travailler dur ; assidu \\
\hline
认真 & rènzhēn & adj. & sérieux, consciencieux, soigneux \\
\hline
\end{tabularx}
\end{center}
\end{aretenir}

\begin{experience}[Activité de classe : Repérage guidé]
\textbf{Consigne :} En binôme, relevez dans le dialogue toutes les structures où un verbe est suivi de \zh{得} + adjectif. Classez-les en deux colonnes : \textbf{Verbe sans objet} (ex: \zh{学得快}) et \textbf{Verbe avec objet} (ex: \zh{汉语说得好}). Présentez vos trouvailles au tableau.
\end{experience}

\courssub{Schéma de la structure}

\begin{propriete}[Structure de base du complément d'état]
Le complément d'état ( \zh{状态补语} ) évalue le \textbf{résultat}, le \textbf{degré} ou la \textbf{manière} dont une action est réalisée. Il répond à la question « Comment ? » (résultat/degré).
\begin{center}
\textbf{Schéma fondamental :}\quad Sujet + Verbe + \cle{得} + Adjectif (+ \zh{很/真/特别/不太} etc.)
\end{center}
\begin{itemize}
\item L'adjectif fonctionne comme noyau du complément (souvent précédé de \zh{很} pour rythmer la phrase).
\item La particule \zh{得} (ton neutre) est \textbf{obligatoire} : elle marque la frontière entre l'action et son évaluation.
\end{itemize}
\end{propriete}

\begin{exemplebox}[Verbe sans objet — Structure simple]
\begin{tabularx}{\linewidth}{@{}lX@{}}
\zh{他跑得快。} & \emph{Tā pǎo de kuài.} — Il court vite. \\
\zh{我睡得早。} & \emph{Wǒ shuì de zǎo.} — Je me couche tôt / Je dors tôt. \\
\zh{这本书写得好。} & \emph{Zhè běn shū xiě de hǎo.} — Ce livre est bien écrit. (Sujet patient, verbe passif) \\
\end{tabularx}
\end{exemplebox}

\begin{propriete}[Cas du verbe transitif : l'objet et la répétition verbale]
Quand le verbe a un objet (\zh{宾语}), celui-ci \textbf{bloque} l'accès au complément d'état. Deux solutions grammaticales s'offrent à vous :
\begin{enumerate}
\item \textbf{Répétition du verbe} (structure standard, neutre) :\\
\hspace{1cm} Sujet + Verbe + Objet + \textbf{Verbe} + \zh{得} + Adjectif
\item \textbf{Antéposition de l'objet} (topicalisation, mise en relief de l'objet) :\\
\hspace{1cm} Sujet + \textbf{Objet} + Verbe + \zh{得} + Adjectif
\end{enumerate}
\end{propriete}

\begin{exemplebox}[Verbe avec objet — Les deux structures correctes]
\begin{tabularx}{\linewidth}{@{}lX@{}}
\textbf{Répétition :} \zh{他说汉语说得流利。} & \emph{Tā shuō Hànyǔ shuō de liúlì.} — Il parle couramment le chinois. \\
\textbf{Antéposition :} \zh{他汉语说得流利。} & \emph{Tā Hànyǔ shuō de liúlì.} — Le chinois, il le parle couramment. \\
\textbf{Répétition :} \zh{我们吃饭吃得很饱。} & \emph{Wǒmen chī fàn chī de hěn bǎo.} — Nous avons mangé à satiété. \\
\textbf{Antéposition :} \zh{我们饭吃得很饱。} & \emph{Wǒmen fàn chī de hěn bǎo.} — Le riz, nous en avons mangé à satiété. \\
\end{tabularx}
\end{exemplebox}

\begin{attention}[Interdits absolus — Erreurs de structure]
Ne placez \textbf{jamais} l'objet après le complément d'état, et n'oubliez \textbf{jamais} \zh{得}.
\begin{itemize}
\item *\zh{他说得流利汉语} \hspace{2cm} (Objet après le complément → \textbf{Incorrect})
\item *\zh{他说汉语流利} \hspace{2.2cm} (Manque \zh{得} → \textbf{Incorrect})
\item *\zh{他汉语说流利} \hspace{2.3cm} (Manque \zh{得} → \textbf{Incorrect})
\end{itemize}
\end{attention}

\begin{center}
\begin{tabularx}{\linewidth}{|X|X|X|}
\hline
\rowcolor{popblueL} Structure & Exemple (Caractères + Pinyin) & Traduction \\
\hline
V + 得 + Adj & \zh{跑得快} \emph{pǎo de kuài} & Courir vite \\
\hline
V + O + V + 得 + Adj & \zh{说汉语说得好} \emph{shuō Hànyǔ shuō de hǎo} & Bien parler chinois \\
\hline
O + V + 得 + Adj & \zh{汉语说得好} \emph{Hànyǔ shuō de hǎo} & Le chinois, bien le parler \\
\hline
S + V + 得 + 很 + Adj & \zh{他写得很工整} \emph{Tā xiě de hěn gōngzhěng} & Il écrit très soigneusement \\
\hline
\end{tabularx}
\end{center}

\courssub{Négation, question et emplois}

\begin{propriete}[Négation : \zh{不} ou \zh{没} ?]
Le complément d'état évalue une action \textbf{réalisée} ou \textbf{habituelle}. On utilise \textbf{\zh{不}} (pas \zh{没}) inséré \textbf{entre \zh{得} et l'adjectif}.
\begin{center}
\textbf{Sujet + Verbe (+ Objet + Verbe) + 得 + \cle{不} + Adjectif}
\end{center}
\end{propriete}

\begin{exemplebox}[Exemples de négation]
\zh{他跑得不快。} \emph{Tā pǎo de bù kuài.} — Il ne court pas vite.\\
\zh{她说汉语说得不好。} \emph{Tā shuō Hànyǔ shuō de bù hǎo.} — Elle ne parle pas bien le chinois.\\
\zh{这道题我做得不对。} \emph{Zhè dào tí wǒ zuò de bù duì.} — Je n'ai pas fait cet exercice correctement.
\end{exemplebox}
\begin{attention}[Piège : \zh{没} + Verbe + \zh{得}... n'existe pas pour nier l'état.]
*\zh{他没跑得快} est incorrect. Pour nier le passé accompli, on nie le verbe principal : \zh{他没跑} (Il n'a pas couru), mais on ne peut pas dire « il n'a pas couru vite » avec \zh{没}.
\end{attention}


\begin{propriete}[Formes interrogatives]
Trois façons principales de questionner le complément d'état :
\begin{enumerate}
\item \textbf{Particule \zh{吗} en fin de phrase} (question fermée) :\\
\zh{他跑得快吗？} \emph{Tā pǎo de kuài ma?} — Court-il vite ?
\item \textbf{Forme A-\zh{不}-A sur l'adjectif} (question alternative, très courante) :\\
\zh{他跑得快不快？} \emph{Tā pǎo de kuài bu kuài?} — Court-il vite ou pas ?\\
\zh{汉语说得好不好？} \emph{Hànyǔ shuō de hǎo bu hǎo?} — Le chinois, est-il bien parlé ?
\item \textbf{Mot interrogatif \zh{怎么样}} (question ouverte) :\\
\zh{他跑得怎么样？} \emph{Tā pǎo de zěnmeyàng?} — Comment court-il ? / Quelle est sa vitesse ?\\
\zh{字写得怎么样？} \emph{Zì xiě de zěnmeyàng?} — Comment sont écrits les caractères ?
\end{enumerate}
\end{propriete}

\begin{methode}[Comment répondre ?]
\begin{itemize}
\item À \zh{...吗？} / \zh{...不...?} : \zh{很快} / \zh{不快} / \zh{一般} (moyen).
\item À \zh{...怎么样？} : Réponse complète avec \zh{得} : \zh{跑得很快} / \zh{写得不太工整}.
\end{itemize}
\end{methode}

\begin{propriete}[Emplois principaux du complément d'état]
\begin{enumerate}
\item \textbf{Évaluation de la compétence / performance} (très fréquent en milieu scolaire/pro) :\\
\zh{你中文说得真标准！} \emph{Nǐ Zhōngwén shuō de zhēn biāozhǔn!} — Ton chinois est vraiment standard !
\item \textbf{Description du résultat d'une action concrète} :\\
\zh{屋子打扫得干干净净。} \emph{Wūzi dǎsǎo de gāngānjìngjìng.} — La pièce a été nettoyée à fond. (Reduplication d'adjectif pour l'intensité).
\item \textbf{Jugement sur un comportement / attitude} :\\
\zh{他考虑得很周到。} \emph{Tā kǎolǜ de hěn zhōudào.} — Il a pensé à tout / Il a réfléchi très soigneusement.
\end{enumerate}
\end{propriete}

\courssub{Comparaison avec le français et pièges}

Nous savons désormais construire le complément d'état à l'affirmative (\zh{说得好}), à la négative (\zh{说得不好}) et à la forme interrogative (\zh{说得好不好？}). Mais pour un élève francophone, le vrai danger n'est pas la structure elle-même : c'est le \cle{français} qui vient parasiter le chinois. Cette section met face à face les deux langues, démonte les trois particules homophones \zh{的}、\zh{地}、\zh{得}, et vous donne un test infaillible pour ne plus jamais vous tromper.

\courssub{Le français dit « il court vite », le chinois dit \zh{他跑得快}}

Commençons par observer deux phrases qui expriment exactement la même idée.

\begin{exemplebox}[Français et chinois face à face]
\textbf{Français :} Il court \emph{vite}. (verbe + adverbe, dans cet ordre)\\
\textbf{Chinois :} \zh{他跑得快。} \emph{Tā pǎo de kuài.} « Il court vite. » (verbe + 得 + adjectif)\\[4pt]
\textbf{Français :} Elle parle \emph{bien} le chinois.\\
\textbf{Chinois :} \zh{她说汉语说得很好。} \emph{Tā shuō Hànyǔ shuō de hěn hǎo.} « Elle parle bien le chinois. »
\end{exemplebox}

Deux différences fondamentales sautent aux yeux :

\begin{enumerate}
\item En français, l'évaluation est un \cle{adverbe} (\emph{vite, bien, lentement}) placé directement après le verbe, sans aucun mot-outil. En chinois, l'évaluation est un \cle{adjectif} (\zh{快}、\zh{好}、\zh{慢}) qui ne peut \textbf{jamais} toucher directement le verbe : la particule \zh{得} est obligatoire entre les deux.
\item En français, on peut enchaîner « parler + bien + le chinois ». En chinois, c'est impossible : le verbe \zh{说} ne peut pas porter à la fois son objet \zh{汉语} et son complément d'état. Il faut donc \textbf{répéter le verbe} : \zh{说汉语说得好}.
\end{enumerate}

\begin{propriete}[La règle d'or de la transposition]
Quand le français dit « verbe + adverbe (+ objet) », le chinois exige :
\begin{center}
\textbf{Sujet + Verbe + (Objet) + Verbe + 得 + Adjectif}
\end{center}
Si le verbe n'a pas d'objet, on ne le répète pas : \zh{他跑得快。} \emph{Tā pǎo de kuài.} « Il court vite. »\\
Si le verbe a un objet, deux options correctes :
\begin{itemize}
\item répéter le verbe : \zh{他说汉语说得很好。} \emph{Tā shuō Hànyǔ shuō de hěn hǎo.}
\item avancer l'objet : \zh{他汉语说得很好。} \emph{Tā Hànyǔ shuō de hěn hǎo.} « Le chinois, il le parle très bien. »
\end{itemize}
Jamais : *\zh{他说得很好汉语} ni *\zh{他说汉语很好} dans ce sens (évaluation de l'action de parler).
\end{propriete}

\begin{popfigure}
\begin{tikzpicture}[font=\small, node distance=0.6cm]
\node[draw=popblue, fill=popblueL, rounded corners, align=center, text width=4.6cm] (fr) {\textbf{Français}\\ Il parle bien le chinois.};
\node[draw=popgreen, fill=popgreenL, rounded corners, align=center, text width=5.2cm, right=3.2cm of fr] (zh) {\textbf{Chinois}\\ 他说汉语说得很好。\\ \emph{Tā shuō Hànyǔ shuō de hěn hǎo.}};
\draw[-{Stealth[length=3mm]}, very thick, poporange] (fr) -- node[above, align=center, text width=2.6cm]{\scriptsize 1. répéter le verbe\\ \scriptsize 2. insérer 得\\ \scriptsize 3. adverbe $\rightarrow$ adjectif} (zh);
\node[draw=poppurple, fill=poppurpleL, rounded corners, align=center, text width=2.1cm, below=0.9cm of fr] (v) {\scriptsize verbe\\ 说 shuō};
\node[draw=poppurple, fill=poppurpleL, rounded corners, align=center, text width=2.1cm, right=0.4cm of v] (o) {\scriptsize objet\\ 汉语 Hànyǔ};
\node[draw=poporange, fill=poporangeL, rounded corners, align=center, text width=2.6cm, right=0.4cm of o] (c) {\scriptsize verbe + 得 + adj.\\ 说得很好};
\draw[-{Stealth[length=2mm]}, popmuted] (fr.south) -- (v.north -| fr.south);
\draw[-{Stealth[length=2mm]}, popmuted] (v.east) -- (o.west);
\draw[-{Stealth[length=2mm]}, popmuted] (o.east) -- (c.west);
\end{tikzpicture}
\legende{Le pont de transformation : du français « Il parle bien le chinois » au chinois \zh{他说汉语说得很好}.}
\end{popfigure}

\courssub{Les trois particules « de » : 的、地、得}

Voici le point le plus délicat de l'écriture chinoise pour un francophone. À l'oral, \zh{的}、\zh{地}、\zh{得} se prononcent toutes \emph{de} (ton neutre). À l'écrit, ce sont trois caractères différents, avec trois fonctions et trois positions différentes.

\begin{propriete}[Les trois « de » du chinois]
\begin{itemize}
\item \zh{的} + \textbf{nom} : marque de la détermination (possession, qualificatif). \zh{我的书} \emph{wǒ de shū} « mon livre » ; \zh{喀麦隆的足球} \emph{Kāmàilóng de zúqiú} « le football du Cameroun ».
\item \zh{地} + \textbf{verbe} : marque du complément circonstanciel de manière, placé \textbf{avant} le verbe. \zh{他慢慢地走。} \emph{Tā mànmàn de zǒu.} « Il marche lentement. »
\item verbe + \zh{得} + \textbf{adjectif} : complément d'état, placé \textbf{après} le verbe, qui évalue l'action. \zh{他走得慢。} \emph{Tā zǒu de màn.} « Il marche lentement » (constat, évaluation de sa façon de marcher).
\end{itemize}
\end{propriete}

\begin{center}
\begin{tabularx}{\linewidth}{|X|X|X|X|}
\hline
\rowcolor{popblueL} Particule & Position & Fonction & Exemple traduit \\
\hline
的 de & avant un nom & détermination, possession & 我的老师 wǒ de lǎoshī « mon professeur » \\
\hline
地 de & avant un verbe & manière (comment ?) & 她高兴地说。Tā gāoxìng de shuō. « Elle dit joyeusement. » \\
\hline
得 de & après un verbe & état, évaluation (résultat) & 她说得很高兴。Tā shuō de hěn gāoxìng. « Elle parle avec joie. » \\
\hline
\end{tabularx}
\end{center}

\begin{exemplebox}[地 ou 得 ? Le sens change avec la position]
\zh{她高兴地说。} \emph{Tā gāoxìng de shuō.} « Elle parle joyeusement. » — \zh{高兴地} décrit l'humeur \textbf{pendant} l'action (manière, avant le verbe).\\[4pt]
\zh{她说得很高兴。} \emph{Tā shuō de hěn gāoxìng.} « Elle parle avec joie / Elle est contente de parler. » — On \textbf{évalue} sa façon de parler ou son état d'esprit résultant (complément d'état, après le verbe).\\[4pt]
Autre paire : \zh{他认真地写字。} \emph{Tā rènzhēn de xiě zì.} « Il écrit avec sérieux » (attitude, manière) / \zh{他字写得很好。} \emph{Tā zì xiě de hěn hǎo.} « Il écrit bien les caractères » (évaluation du résultat).
\end{exemplebox}

\begin{methode}[Le test de position : trois questions, zéro erreur]
Avant d'écrire « de », posez-vous la question : \textbf{où suis-je dans la phrase ?}
\begin{enumerate}
\item Ce qui suit est-il un \textbf{nom} ? $\rightarrow$ écrivez \zh{的} (我的书, 老师的车).
\item Ce qui suit est-il un \textbf{verbe}, et le « de » décrit-il la \textbf{manière} de faire l'action ? $\rightarrow$ écrivez \zh{地} (慢慢地走, 高兴地唱).
\item Suis-je \textbf{juste après un verbe}, et ce qui suit \textbf{évalue} l'action (degré, résultat) ? $\rightarrow$ écrivez \zh{得} (跑得快, 说得流利).
\end{enumerate}
Résumé : \textbf{nom $\rightarrow$ 的 ; manière avant le verbe $\rightarrow$ 地 ; évaluation après le verbe $\rightarrow$ 得.}
\end{methode}

\begin{savaistu}[Pourquoi trois caractères pour un seul son ?]
À l'oral, \zh{的}、\zh{地}、\zh{得} se prononcent toutes \emph{de}, au ton neutre : un Chinois qui parle ne fait aucune différence entre elles, et personne ne la remarque. C'est uniquement à l'\textbf{écrit} que la distinction existe, héritée de l'écriture moderne normalisée au XX\textsuperscript{e} siècle. Conséquence pratique : dans une dictée, seule la \textbf{position dans la phrase} vous permet de choisir le bon caractère. C'est exactement comme en français « a / à », « et / est », « son / sont » : même son, orthographes différentes, et c'est la grammaire qui tranche.
\end{savaistu}

\courssub{Piège n° 1 : le caractère 得 a plusieurs prononciations}

\begin{attention}[得 se lit « de »… mais pas toujours !]
Le caractère \zh{得} est polyphonique :
\begin{itemize}
\item \zh{得} \emph{de} (ton neutre) : particule du complément d'état. \zh{跑得快} \emph{pǎo de kuài} « courir vite ».
\item \zh{得} \emph{dé} (2\textsuperscript{e} ton) : verbe « obtenir ». \zh{得到} \emph{dédào} « obtenir » ; \zh{他得了一百分。} \emph{Tā dé le yì bǎi fēn.} « Il a obtenu cent points. »
\item \zh{得} \emph{děi} (3\textsuperscript{e} ton) : auxiliaire « devoir, falloir ». \zh{我得走了。} \emph{Wǒ děi zǒu le.} « Je dois partir. »
\end{itemize}
Même caractère, trois lectures, trois fonctions. Dans le complément d'état, c'est toujours \emph{de}, ton neutre.
\end{attention}

\courssub{Piège n° 2 : *\zh{他跑很快} — l'adjectif seul ne suffit pas}

En français, « il court vite » met l'adverbe directement après le verbe. Le francophone traduit alors mot à mot : *\zh{他跑很快}. C'est \textbf{incorrect} dans le sens « il court vite » : après un verbe, l'évaluation exige \zh{得}.

\begin{attention}[很 + adjectif ne remplace pas 得 + adjectif]
\begin{itemize}
\item *\zh{他跑很快。} — faute typique du francophone (calque de « il court vite »).
\item \zh{他跑得很快。} \emph{Tā pǎo de hěn kuài.} « Il court très vite. » — correct.
\item \zh{他跑得很快吗？} \emph{Tā pǎo de hěn kuài ma ?} « Court-il vite ? »
\end{itemize}
Attention au faux ami : \zh{他很快。} \emph{Tā hěn kuài.} signifie « Il est rapide » (description de la personne, verbe \zh{是} sous-entendu), pas « il court vite » (évaluation de l'action de courir).
\end{attention}

\courssub{Piège n° 3 : placer le complément avant le verbe}

Le français place parfois la manière avant le verbe (« il \emph{bien} parle » est fautif, mais « il \emph{vite} courut » existe en littérature, et surtout l'élève pense « bien » avant de penser « parler »). En chinois, le complément d'état est \textbf{toujours} après le verbe, sans exception.

\begin{attention}[Ordre des mots : le complément d'état ne remonte jamais]
\begin{itemize}
\item *\zh{他得快跑。} — incorrect (et incompréhensible : on croirait \zh{得} « děi », devoir !).
\item *\zh{她很好说汉语。} — incorrect.
\item Correct : \zh{他跑得快。} / \zh{她说汉语说得很好。}
\end{itemize}
Si vous voulez exprimer la manière \textbf{avant} le verbe, c'est possible, mais alors il faut changer d'outil : \zh{地}. \zh{他高兴地唱。} \emph{Tā gāoxìng de chàng.} « Il chante joyeusement » (manière) ≠ \zh{他唱得很高兴。} « Il chante avec joie » (évaluation). Les deux existent, mais ce ne sont pas les mêmes structures.
\end{attention}

\begin{exoresolu}[Traduire et choisir la bonne particule]
Énoncé — Traduisez en chinois : 1) « Mon frère chante très bien. » 2) « Les élèves écoutent attentivement le professeur. » 3) « C'est le marché de Douala. » Indiquez pour chaque phrase quelle particule « de » vous utilisez et pourquoi.
\tcblower
Solution détaillée :
\begin{enumerate}
\item « Mon frère chante très bien » : évaluation \textbf{après} le verbe $\rightarrow$ \zh{得}. Le verbe \zh{唱} n'a pas d'objet, pas de répétition : \zh{我哥哥唱得很好。} \emph{Wǒ gēge chàng de hěn hǎo.}
\item « Les élèves écoutent attentivement » : manière \textbf{avant} le verbe $\rightarrow$ \zh{地}. \zh{学生们认真地听老师讲课。} \emph{Xuéshengmen rènzhēn de tīng lǎoshī jiǎng kè.}
\item « Le marché de Douala » : détermination \textbf{avant un nom} $\rightarrow$ \zh{的}. \zh{这是杜阿拉的市场。} \emph{Zhè shì Dù'ālā de shìchǎng.}
\end{enumerate}
Méthode : dans chaque cas, on a d'abord repéré la \textbf{position} (avant nom / avant verbe / après verbe), puis choisi le caractère. Position d'abord, caractère ensuite.
\end{exoresolu}

\begin{aretenir}[L'essentiel de la section]
\begin{itemize}
\item Le français évalue avec un adverbe après le verbe ; le chinois évalue avec \textbf{verbe + 得 + adjectif}, et répète le verbe s'il y a un objet : \zh{他说汉语说得好}.
\item Trois particules « de » : \zh{的} avant un nom, \zh{地} avant un verbe (manière), \zh{得} après un verbe (état). Même son, trois écritures : seule la position permet de choisir.
\item \zh{得} se lit aussi \emph{dé} (obtenir) et \emph{děi} (devoir) : ne pas confondre.
\item Interdits absolus : *\zh{他跑很快} (il manque 得), *\zh{得快跑} (le complément d'état ne précède jamais le verbe).
\end{itemize}
\end{aretenir}

\courssub{Exercices d'application et de transformation}

\begin{exercice}[Entraînement structuré]
\textbf{Exercice 1 — Complétez avec la bonne particule (\zh{的}、\zh{地}、\zh{得}).}
\begin{enumerate}
\item 这道菜做\zh{\_\_\_\_}真好吃！
\item 请大家安静\zh{\_\_\_\_}听讲。
\item 我哥哥跑\zh{\_\_\_\_}特别快。
\item 这是我\zh{\_\_\_\_}新书包。
\item 她中文说\zh{\_\_\_\_}很流利。
\item 同学们认真\zh{\_\_\_\_}复习课文。
\end{enumerate}

\textbf{Exercice 2 — Transformez les phrases en utilisant la structure « Verbe + Objet + Verbe + 得 + Adj ».}
\textit{Modèle : 他说中文说得好。 $\rightarrow$ 他说中文说得很好。 / 中文他说得很好。}
\begin{enumerate}
\item 我写汉字写得工整。
\item 他们踢足球踢得很好。
\item 妹妹弹钢琴弹得不好。
\item 我们学中文学得很快。
\end{enumerate}

\textbf{Exercice 3 — Mettez à la forme négative (avec \zh{不}) et interrogative (avec \zh{好不好} ou \zh{怎么样}).}
\begin{enumerate}
\item 他开车开得小心。 $\rightarrow$ Nég. : \_\_\_\_\_\_\_\_\_\_ ; Interr. : \_\_\_\_\_\_\_\_\_\_
\item 这件衣服洗得干净。 $\rightarrow$ Nég. : \_\_\_\_\_\_\_\_\_\_ ; Interr. : \_\_\_\_\_\_\_\_\_\_
\end{enumerate}

\textbf{Exercice 4 — Traduisez en chinois (caractères + pinyin).}
\begin{enumerate}
\item Le professeur explique très clairement. (expliquer = \zh{讲课} / \zh{解释} ; clairement = \zh{清楚})
\item Nous avons mal dormi hier soir. (dormir = \zh{睡觉} ; mal = \zh{不好})
\item Les travailleurs de l'usine travaillent dur. (usine = \zh{工厂} ; travailleur = \zh{工人} ; dur = \zh{努力})
\item Ce plat de ndolé est délicieux ! (plat = \zh{菜} ; ndolé = \zh{恩多莱} ; délicieux = \zh{好吃})
\end{enumerate}

\textbf{Exercice 5 — Correction d'erreurs : Chaque phrase contient une faute (structure, particule, ordre). Corrigez-la.}
\begin{enumerate}
\item *他汉语说很好。
\item *她跑得很慢地。
\item *学生们认真得听课。
\item *我得了一本书很好看。 (Indice : deux verbes \zh{得} différents !)
\item *我们吃得很饱饭。
\end{enumerate}
\end{exercice}

\begin{corrige}[Exercices de la Leçon 32]
\textbf{Exercice 1 :}
1. \zh{得} (évaluation après verbe \zh{做}) — \emph{Zhè dào cài zuò \textbf{de} zhēn hǎochī!}
2. \zh{地} (manière avant verbe \zh{听}) — \emph{Qǐng dàjiā ānjìng \textbf{de} tīng jiǎng.}
3. \zh{得} (évaluation après verbe \zh{跑}) — \emph{Wǒ gēge pǎo \textbf{de} tèbié kuài.}
4. \zh{的} (détermination avant nom \zh{书包}) — \emph{Zhè shì wǒ \textbf{de} xīn shūbāo.}
5. \zh{得} (évaluation, structure V+O+V+de+Adj) — \emph{Tā Zhōngwén shuō \textbf{de} hěn liúlì.}
6. \zh{地} (manière avant verbe \zh{复习}) — \emph{Tóngxuémen rènzhēn \textbf{de} fùxí kèwén.}

\textbf{Exercice 2 :} (Les deux formes sont acceptées)
1. \zh{我汉字写得工整。} / \zh{我写汉字写得很工整。}
2. \zh{他们足球踢得很好。} / \zh{他们踢足球踢得很好。}
3. \zh{妹妹钢琴弹得不好。} / \zh{妹妹弹钢琴弹得不好。}
4. \zh{我们中文学得很快。} / \zh{我们学中文学得很快。}

\textbf{Exercice 3 :}
1. Nég. : \zh{他开车开得不小心。} Interr. : \zh{他开车开得小心不小心？} / \zh{他开车开得怎么样？}
2. Nég. : \zh{这件衣服洗得不干净。} Interr. : \zh{这件衣服洗得干净不干净？} / \zh{这件衣服洗得怎么样？}

\textbf{Exercice 4 :}
1. \zh{老师讲课讲得很清楚。} \emph{Lǎoshī jiǎngkè jiǎng de hěn qīngchu.} (ou \zh{老师解释得很清楚。})
2. \zh{我们昨晚睡得不好。} \emph{Wǒmen zuówǎn shuì de bù hǎo.}
3. \zh{工厂的工人工作很努力。} \emph{Gōngchǎng de gōngrén gōngzuò hěn nǔlì.} (Attention : \zh{努力} ici est adverbe/adjectif prédicatif sans \zh{得} car \zh{工作} est nom+v, mais \zh{工人们工作得很努力} est aussi correct). \textit{Note :} \zh{努力} peut fonctionner comme adverbe pur (\zh{地}) ou adjectif dans compl. d'état (\zh{得}). Les deux passent : \zh{工人们努力地工作} / \zh{工人们工作得很努力}.
4. \zh{这道恩多莱菜真好吃！} \emph{Zhè dào Ēnduōlái cài zhēn hǎochī!} (Pas de \zh{得} ici, \zh{好吃} est adjectif prédicatif direct).

\textbf{Exercice 5 :}
1. *\zh{他汉语说很好} $\rightarrow$ \zh{他汉语说得很好} (manque \zh{得}) ou \zh{他说汉语说得很好}.
2. *\zh{她跑得很慢地} $\rightarrow$ \zh{她跑得很慢} (supprimer \zh{地}, \zh{得} déjà présent) ou \zh{她很慢地跑} (changer structure).
3. *\zh{学生们认真得听课} $\rightarrow$ \zh{学生们认真地听课} (\zh{得} $\rightarrow$ \zh{地} : manière avant verbe).
4. *\zh{我得了一本书很好看} $\rightarrow$ \zh{我得到了一本很好看的书。} (\zh{得} \emph{dé} = obtenir ; \zh{的} pour \zh{很好看的书}). Ne pas confondre les trois \zh{de} !
5. *\zh{我们吃得很饱饭} $\rightarrow$ \zh{我们饭吃得很饱} (antéposition) ou \zh{我们吃饭吃得很饱} (répétition). L'objet \zh{饭} ne peut pas rester après \zh{饱}.
\end{corrige}

\coursec{Exercices d'application et de transformation}

Dans la section précédente, nous avons comparé le complément d'état avec le français et repéré les pièges qui guettent les apprenants francophones : l'oubli de \zh{得}, la répétition du verbe, la place de l'objet. Il est maintenant temps de passer de la théorie à la pratique. Cette dernière section vous propose un parcours progressif en cinq étapes : d'abord \emph{reconnaître} les mots-outils, puis \emph{manipuler} l'ordre des mots, ensuite \emph{transformer} les phrases (négation, question), puis \emph{traduire} dans les deux sens, et enfin \emph{produire} librement. C'est exactement le chemin que suivent les épreuves d'examen.

\begin{popfigure}
\begin{tikzpicture}[node distance=0.35cm]
\node[draw=popblue, fill=popblueL, rounded corners, minimum width=2.5cm, minimum height=0.9cm, align=center, font=\small\bfseries] (e1) {1. Reconnaissance\\\zh{得 / 的 / 地}};
\node[draw=popgreen, fill=popgreenL, rounded corners, minimum width=2.5cm, minimum height=0.9cm, align=center, font=\small\bfseries, right=of e1] (e2) {2. Manipulation\\ordre des mots};
\node[draw=poporange, fill=poporangeL, rounded corners, minimum width=2.5cm, minimum height=0.9cm, align=center, font=\small\bfseries, right=of e2] (e3) {3. Transformation\\négation, question};
\node[draw=poppurple, fill=poppurpleL, rounded corners, minimum width=2.5cm, minimum height=0.9cm, align=center, font=\small\bfseries, right=of e3] (e4) {4. Traduction\\thème et version};
\node[draw=poppink, fill=poppinkL, rounded corners, minimum width=2.5cm, minimum height=0.9cm, align=center, font=\small\bfseries, right=of e4] (e5) {5. Production\\phrases libres};
\draw[-{Stealth}, thick, popdark] (e1) -- (e2);
\draw[-{Stealth}, thick, popdark] (e2) -- (e3);
\draw[-{Stealth}, thick, popdark] (e3) -- (e4);
\draw[-{Stealth}, thick, popdark] (e4) -- (e5);
\end{tikzpicture}
\legende{Frise de progression des exercices : de la reconnaissance à la production libre.}
\end{popfigure}

\begin{center}
\begin{tabularx}{\linewidth}{|X|X|X|X|}
\hline
\rowcolor{popblueL} \textbf{Caractères} & \textbf{Pinyin} & \textbf{Nature} & \textbf{Traduction} \\
\hline
越来越 & yuè lái yuè & adv. & de plus en plus \\
\hline
非常 & fēi cháng & adv. & extrêmement, très \\
\hline
不错 & bù cuò & adj. & pas mal, bien \\
\hline
清楚 & qīng chu & adj. & clair, distinct \\
\hline
流利 & liú lì & adj. & courant (langue), fluide \\
\hline
篮球 & lán qiú & n. & basket-ball \\
\hline
法语 & Fǎ yǔ & n. & français (langue) \\
\hline
踢 & tī & v. & jouer (au foot), donner un coup de pied \\
\hline
跑步 & pǎo bù & v. (VO) & courir, faire du jogging \\
\hline
跳舞 & tiào wǔ & v. (VO) & danser \\
\hline
汉字 & Hàn zì & n. & caractère chinois \\
\hline
\end{tabularx}
\end{center}

\courssub{Exercice 1 : choisir le bon mot-outil — 得, 的 ou 地}

Avant de vous lancer, rappelez-vous la grande distinction étudiée dans cette leçon : \cle{得} introduit le complément d'état \emph{après} un verbe ; \cle{的} relie un déterminant à un nom ; \cle{地} relie un adverbe à un verbe. Ces trois caractères se prononcent tous \emph{de} (ton neutre), mais ils ne s'échangent jamais.

\begin{exoresolu}[Exercice 1 — Compléter avec 得, 的 ou 地]
Complétez : a) 他唱\_\_很好。b) 我\_\_老师。c) 她高兴\_\_说。
\tcblower
\textbf{Solutions :}\\
a) \zh{他唱得很好。}\emph{Tā chàng de hěn hǎo.} « Il chante très bien. » Après un verbe (\zh{唱}), on introduit l'évaluation avec \cle{得}.\\
b) \zh{我的老师。}\emph{Wǒ de lǎoshī.} « Mon professeur. » Entre un possesseur et un nom, on utilise \cle{的}.\\
c) \zh{她高兴地说。}\emph{Tā gāoxìng de shuō.} « Elle parle joyeusement. » Entre un adverbe de manière et un verbe, on utilise \cle{地}.
\end{exoresolu}

\begin{attention}[Le trio 得 / 的 / 地]
Les trois se prononcent \emph{de} (ton neutre), mais : \zh{得} suit le verbe (complément d'état), \zh{的} précède le nom (possession, qualification), \zh{地} précède le verbe (adverbe de manière). Écrire \zh{他唱的很好} est une faute de caractère, même si la prononciation est correcte.
\end{attention}

\courssub{Exercice 2 : remettre les mots dans l'ordre}

\begin{exoresolu}[Exercice 2 — Phrase désordonnée]
Remettez dans l'ordre : 说 / 汉语 / 得 / 他 / 很好 / 说
\tcblower
\textbf{Solution :} \zh{他汉语说得很好。} ou \zh{他说汉语说得很好。}\emph{Tā Hànyǔ shuō de hěn hǎo. / Tā shuō Hànyǔ shuō de hěn hǎo.} « Il parle très bien chinois. »\\
\textbf{Méthode :} le verbe \zh{说} apparaît deux fois : c'est le signal de la répétition du verbe. On place d'abord Sujet + verbe + objet (\zh{他说汉语}), puis on répète le verbe, puis \zh{得} + adverbe de degré + adjectif (\zh{说得很好}). Variante élégante : anticiper l'objet \zh{汉语} juste après le sujet (topicalisation).
\end{exoresolu}

\courssub{Exercices 3 et 4 : transformations — négation et question}

\begin{methode}[Transformer une phrase à complément d'état]
\begin{enumerate}
\item \textbf{Négation :} remplacer l'adverbe de degré (souvent \zh{很}) par \zh{不} devant l'adjectif : \zh{写得很快} → \zh{写得不快}. On ne dit jamais *\zh{不得快}.
\item \textbf{Question d'évaluation :} remplacer tout le groupe « degré + adjectif » par \zh{怎么样} : \zh{跳得很好} → \zh{跳得怎么样？}
\item Vérifier que \zh{得} reste en place dans les deux cas.
\end{enumerate}
\end{methode}

\begin{exercice}[Exercice 3 — Passage à la négative]
Transformez : a) 他写得很快。b) 她唱得很好。c) 我哥哥足球踢得不错。
\end{exercice}

\begin{corrige}[Exercice 3]
a) \zh{他写得不快。}\emph{Tā xiě de bù kuài.} « Il n'écrit pas vite. »\\
b) \zh{她唱得不好。}\emph{Tā chàng de bù hǎo.} « Elle ne chante pas bien. »\\
c) \zh{我哥哥足球踢得不好。}\emph{Wǒ gēge zúqiú tī de bù hǎo.} « Mon grand frère ne joue pas bien au football. »
\end{corrige}

\begin{exercice}[Exercice 4 — Question avec 怎么样]
Transformez en question : a) 她跳舞跳得很好。b) 他汉语说得很流利。
\end{exercice}

\begin{corrige}[Exercice 4]
a) \zh{她跳舞跳得怎么样？}\emph{Tā tiàowǔ tiào de zěnmeyàng?} « Comment danse-t-elle ? »\\
b) \zh{他汉语说得怎么样？}\emph{Tā Hànyǔ shuō de zěnmeyàng?} « Comment parle-t-il chinois ? »
\end{corrige}

\begin{attention}[怎么样 exige 得]
Dans les questions d'évaluation de la manière, n'oubliez pas \zh{得} : on dit \zh{你学得怎么样？}\emph{Nǐ xué de zěnmeyàng?} « Comment apprends-tu ? », et non *\zh{你学习怎么样}, qui interroge sur la situation générale (« comment vont tes études ? ») et non sur la manière de faire.
\end{attention}

\courssub{Exercices 5 et 6 : traduction — thème et version}

\begin{methode}[Traduire « il + verbe + adverbe » du français]
Recette en trois étapes :
\begin{enumerate}
\item Écrire le verbe + l'objet éventuel : « jouer au football » → \zh{踢足球}.
\item Répéter le verbe (obligatoire si l'objet est exprimé) : \zh{踢足球踢}.
\item Ajouter \zh{得} + adverbe de degré + adjectif : \zh{踢得非常好}.
\end{enumerate}
Résultat : \zh{他踢足球踢得非常好。} « Il joue très bien au football. »
\end{methode}

\begin{exemplebox}[Modèle de traduction réussie]
« Mon frère joue très bien au football. » → \zh{我哥哥踢足球踢得非常好。}\emph{Wǒ gēge tī zúqiú tī de fēicháng hǎo.} Remarquez : le verbe \zh{踢} est répété, \zh{得} introduit l'évaluation, et l'adverbe de degré \zh{非常} précède l'adjectif \zh{好}.
\end{exemplebox}

\begin{exercice}[Exercice 5 — Thème (français → chinois)]
Traduisez : a) « Elle chante très bien. » b) « Il ne court pas vite. » c) « Comment écris-tu les caractères ? »
\end{exercice}

\begin{corrige}[Exercice 5]
a) \zh{她唱歌唱得很好。}\emph{Tā chàng gē chàng de hěn hǎo.}\\
b) \zh{他跑步跑得不快。}\emph{Tā pǎobù pǎo de bù kuài.} (ou \zh{他跑得不快。})\\
c) \zh{你写汉字写得怎么样？}\emph{Nǐ xiě Hànzì xiě de zěnmeyàng?} (ou \zh{你汉字写得怎么样？})
\end{corrige}

\begin{exercice}[Exercice 6 — Version (chinois → français)]
Traduisez : a) 老师说得很清楚。b) 你英语学得怎么样？
\end{exercice}

\begin{corrige}[Exercice 6]
a) \zh{老师说得很清楚。}\emph{Lǎoshī shuō de hěn qīngchu.} « Le professeur explique (parle) très clairement. »\\
b) \zh{你英语学得怎么样？}\emph{Nǐ Yīngyǔ xué de zěnmeyàng?} « Comment apprends-tu l'anglais ? / Quel est ton niveau en anglais ? »
\end{corrige}

\begin{textebox}[Phrase utile pour l'examen]
我们班的同学汉语说得越来越好。\\
\emph{Wǒmen bān de tóngxué Hànyǔ shuō de yuè lái yuè hǎo.}\\
« Les élèves de notre classe parlent chinois de mieux en mieux. »\\
Retenez la structure \cle{越来越} + adjectif : « de plus en plus… ». Elle se combine parfaitement avec le complément d'état : \zh{说得越来越好} « parler de mieux en mieux ».
\end{textebox}

\courssub{Exercice 7 : production libre guidée}

\begin{experience}[À vous de parler !]
Rédigez trois phrases sur vos propres compétences (sport, langues, musique) avec le complément d'état, puis présentez-les à la classe. Modèles possibles :\\
\zh{我打篮球打得很好。}\emph{Wǒ dǎ lánqiú dǎ de hěn hǎo.} « Je joue bien au basket. »\\
\zh{我法语说得不错，英语说得不太好。}\emph{Wǒ Fǎyǔ shuō de bùcuò, Yīngyǔ shuō de bù tài hǎo.} « Je parle pas mal français, pas très bien anglais. »\\
\zh{我唱歌唱得越来越好。}\emph{Wǒ chàng gē chàng de yuè lái yuè hǎo.} « Je chante de mieux en mieux. »
\end{experience}

\begin{popfigure}
\begin{tikzpicture}
\node[draw=popblue, fill=popblueL, rounded corners, minimum width=3.2cm, minimum height=0.8cm, align=center, font=\small\bfseries] (t) at (0,0) {Grille d'auto-correction};
\node[draw=popgreen, fill=popgreenL, rounded corners, align=center, font=\small, minimum width=4.2cm] (a) at (-4.3,-1.6) {Structure respectée ?\\verbe + 得 + adjectif};
\node[draw=poporange, fill=poporangeL, rounded corners, align=center, font=\small, minimum width=4.2cm] (b) at (0,-1.6) {Verbe répété ?\\si l'objet est exprimé};
\node[draw=poppurple, fill=poppurpleL, rounded corners, align=center, font=\small, minimum width=4.2cm] (c) at (4.3,-1.6) {Adverbe de degré ?\\\zh{很 / 非常 / 不太}};
\draw[-{Stealth}, thick, popdark] (t) -- (a);
\draw[-{Stealth}, thick, popdark] (t) -- (b);
\draw[-{Stealth}, thick, popdark] (t) -- (c);
\end{tikzpicture}
\legende{Grille d'auto-correction : trois questions à se poser avant de rendre sa phrase.}
\end{popfigure}

\begin{center}
\begin{tabularx}{\linewidth}{|X|X|X|}
\hline
\rowcolor{popblueL} \textbf{Forme} & \textbf{Exemple (caractères + pinyin)} & \textbf{Traduction} \\
\hline
Affirmative & \zh{他说汉语说得很好。}\emph{Tā shuō Hànyǔ shuō de hěn hǎo.} & Il parle très bien chinois. \\
\hline
Négative & \zh{他写得不快。}\emph{Tā xiě de bù kuài.} & Il n'écrit pas vite. \\
\hline
Question (\zh{怎么样}) & \zh{她跳舞跳得怎么样？}\emph{Tā tiàowǔ tiào de zěnmeyàng?} & Comment danse-t-elle ? \\
\hline
Avec \zh{越来越} & \zh{汉语说得越来越好。}\emph{Hànyǔ shuō de yuè lái yuè hǎo.} & Parler chinois de mieux en mieux. \\
\hline
\end{tabularx}
\end{center}

\begin{aretenir}[Bilan de la leçon]
Le complément d'état obéit à la recette : \textbf{verbe (+ objet) + verbe répété + 得 + degré + adjectif}. La négation se fait avec \zh{不} devant l'adjectif, la question d'évaluation avec \zh{怎么样} — sans jamais perdre \zh{得}. Distinguez soigneusement \zh{得} (après le verbe), \zh{的} (avant le nom) et \zh{地} (avant le verbe). Avant de rendre une phrase, passez-la à la grille : structure, verbe répété, adverbe de degré. Vous êtes désormais prêts pour l'évaluation du module 4 !
\end{aretenir}

\newpage

\coursec{Activité d'intégration}

\courssub{Objectifs de la leçon}

À la fin de cette leçon, tu seras capable de :
\begin{itemize}
\item identifier et construire le \cle{complément d'état} (结构助词 \zh{得}) pour décrire la manière ou le résultat d'une action ;
\item poser une question d'évaluation avec \zh{得怎么样？} et y répondre par l'affirmative, la négation (\zh{得不太好}, \zh{得不快}) ou la progression (\zh{得越来越} + adjectif) ;
\item respecter l'ordre des mots quand le verbe a un objet : répéter le verbe avant \zh{得} (ex. \zh{他说汉语说得很好}) ;
\item rédiger un paragraphe descriptif (5--6 phrases) présentant les compétences d'une personne avec au moins trois compléments d'état variés ;
\item interagir oralement pour comparer des performances (structure \zh{A 比 B V 得 Adj}).
\end{itemize}

\courssub{Situation déclenchante : la « Semaine des talents »}

Le club de chinois de ton lycée à Yaoundé organise une « Semaine des talents ». Mme Ngo Bell, votre professeure, lance le projet « Mur des talents » (\zh{才艺墙}) : chaque binôme interview un camarade, note ses compétences, puis affiche une présentation écrite. Voici le modèle d'interview affiché au tableau pour guider vos échanges.

\begin{textebox}[Modèle d'interview — Semaine des talents]
A：你会做什么？\\
\emph{Nǐ huì zuò shénme ?}\\
« Que sais-tu faire ? »\\
B：我会说汉语，也会做饭。\\
\emph{Wǒ huì shuō Hànyǔ, yě huì zuòfàn.}\\
« Je sais parler chinois et je sais aussi cuisiner. »\\
A：你汉语说得怎么样？\\
\emph{Nǐ Hànyǔ shuō de zěnmeyàng ?}\\
« Comment parles-tu chinois ? »\\
B：我说得还可以，但是我说得不太快。\\
\emph{Wǒ shuō de hái kěyǐ, dànshì wǒ shuō de bù tài kuài.}\\
« Je parle assez bien, mais je ne parle pas très vite. »\\
A：你做饭做得怎么样？\\
\emph{Nǐ zuòfàn zuò de zěnmeyàng ?}\\
« Comment cuisines-tu ? »\\
B：我做饭做得越来越好！\\
\emph{Wǒ zuòfàn zuò de yuè lái yuè hǎo !}\\
« Je cuisine de mieux en mieux ! »
\end{textebox}

\courssub{Observation et découverte}

\begin{enumerate}
\item \textbf{Souligne} dans le dialogue toutes les occurrences de la particule \zh{得}.
\item \textbf{Repère} la place de l'objet (\zh{汉语}, \zh{做饭}) par rapport au verbe et à \zh{得}.
\item \textbf{Note} les trois types de réponse : évaluation positive (\zh{还可以}), négation (\zh{不太快}), progression (\zh{越来越好}).
\end{enumerate}

\courssub{Schéma de la structure : le complément d'état}

\begin{propriete}[Structure affirmative]
\textbf{Sujet + Verbe (+ Objet) + Verbe + 得 + Adjectif}\\[4pt]
Lorsque le verbe est transitif (il a un objet), on \textbf{répète le verbe} avant \zh{得}. L'objet reste après le premier verbe.
\end{propriete}

\begin{center}
\begin{tabularx}{\linewidth}{|X|X|X|}
\hline
\rowcolor{popblueL} Structure & Exemple (caractères + pinyin) & Traduction \\
\hline
Verbe intransitif & \zh{他跑得很快。} (Tā pǎo de hěn kuài.) & Il court très vite. \\
\hline
Verbe transitif (objet) & \zh{他说汉语说得很流利。} (Tā shuō Hànyǔ shuō de hěn liúlì.) & Il parle chinois très couramment. \\
\hline
Verbe + objet composé & \zh{她做饭做得很好吃。} (Tā zuòfàn zuò de hěn hǎochī.) & Elle cuisine très bon (ses plats sont délicieux). \\
\hline
\end{tabularx}
\end{center}

\begin{propriete}[Négation, question, progression]
\begin{itemize}
\item \textbf{Négation} : \zh{得 + 不 + (太) + Adjectif} \quad → \zh{做得不太好} (pas très bien), \zh{跑得不快} (pas vite).
\item \textbf{Question ouverte} : \zh{得 + 怎么样？} \quad → \zh{你汉语说得怎么样？}
\item \textbf{Question fermée} : \zh{... 得 + Adjectif + 吗？} \quad → \zh{他唱歌唱得好听吗？}
\item \textbf{Progression} : \zh{得 + 越来越 + Adjectif} \quad → \zh{写得越来越清楚} (écrire de plus en plus clair). \textbf{Jamais} \zh{越来越很 + Adj}.
\end{itemize}
\end{propriete}

\courssub{Comparaison avec le français et pièges fréquents}

\begin{attention}[Erreurs typiques des francophones — à bannir]
\begin{itemize}
\item \textbf{Ordre objet / complément d'état} : \textcolor{popred}{✗} \zh{他说得很好汉语} \quad → \textcolor{popgreen}{✓} \zh{他说汉语说得很好}. \emph{(En français : « Il parle bien chinois » — l'adverbe suit le verbe ; en chinois, l'objet s'intercale et le verbe se répète.)}
\item \textbf{Oubli de \zh{得}} : \textcolor{popred}{✗} \zh{他跑很快} \quad → \textcolor{popgreen}{✓} \zh{他跑得很快}. \emph{(Sans \zh{得}, \zh{很快} devient un adverbe de manière placé avant le verbe : \zh{他很快跑完了} « Il a fini de courir vite », sens différent.)}
\item \textbf{Négation mal placée} : \textcolor{popred}{✗} \zh{不得好} \quad → \textcolor{popgreen}{✓} \zh{得不好} / \zh{得不太好}. \emph{(\zh{不} nie l'adjectif, pas la particule.)}
\item \textbf{\zh{越来越} + \zh{很}} : \textcolor{popred}{✗} \zh{越来越很好} \quad → \textcolor{popgreen}{✓} \zh{越来越好}. \emph{(\zh{越来越} inclut déjà l'idée de degré élevé.)}
\item \textbf{Confusion \zh{得} / \zh{的} / \zh{地}} : \zh{得} (complément d'état) ≠ \zh{的} (attributif, possession) ≠ \zh{地} (adverbialisant). \emph{Ex. : \zh{跑得快} (court vite) vs \zh{快跑} (cours vite ! / ordre) vs \zh{快的跑法} (méthode de course rapide).}
\end{itemize}
\end{attention}

\begin{aretenir}[Français vs Chinois : résumé contrastif]
\begin{center}
\begin{tabularx}{\linewidth}{|X|X|X|}
\hline
\rowcolor{popblueL} Français & Chinois (structure) & Remarque \\
\hline
Il parle bien chinois. & \zh{他说中文说得好。} & Répétition du verbe \zh{说} ; \zh{得} + adjectif. \\
\hline
Il ne parle pas très vite. & \zh{他说得不太快。} & \zh{不太} devant l'adjectif après \zh{得}. \\
\hline
Il parle de mieux en mieux. & \zh{他说得越来越好。} & \zh{越来越} + adjectif seul. \\
\hline
Comment parle-t-il chinois ? & \zh{他中文说得怎么样？} & Objet fronté (topicalisation) possible. \\
\hline
\end{tabularx}
\end{center}
\end{aretenir}

\courssub{Lexique de la leçon (HSK 2--3)}

\begin{center}
\begin{tabularx}{\linewidth}{|X|X|X|X|}
\hline
\rowcolor{popblueL} Caractères & Pinyin & Nature & Traduction \\
\hline
得 & de & particule & marque du complément d'état \\
\hline
怎么样 & zěnmeyàng & pronom interrog. & comment, comment est-ce que \\
\hline
越来越 & yuè lái yuè & adverbe & de plus en plus \\
\hline
唱歌 & chànggē & verbe & chanter \\
\hline
跳舞 & tiàowǔ & verbe & danser \\
\hline
踢足球 & tī zúqiú & loc. verbale & jouer au football \\
\hline
做饭 & zuòfàn & loc. verbale & cuisiner, faire la cuisine \\
\hline
写字 & xiězì & loc. verbale & écrire (des caractères) \\
\hline
跑步 & pǎobù & verbe & courir, faire du jogging \\
\hline
画画 & huàhuà & verbe & dessiner, peindre \\
\hline
快 & kuài & adjectif & rapide, vite \\
\hline
慢 & màn & adjectif & lent, lentement \\
\hline
清楚 & qīngchu & adjectif & clair, distinct \\
\hline
流利 & liúlì & adjectif & courant, fluide (langue) \\
\hline
好吃 & hǎochī & adjectif & bon (à manger), délicieux \\
\hline
还可以 & hái kěyǐ & locution & passable, assez bien, ça va \\
\hline
非常 & fēicháng & adverbe & extrêmement, très \\
\hline
才艺 & cáiyì & nom & talent, aptitude artistique \\
\hline
墙 & qiáng & nom & mur \\
\hline
球员 & qiúyuán & nom & joueur (de ballon) \\
\hline
有才华 & yǒu cáihuá & loc. adj. & doué, talentueux \\
\hline
\end{tabularx}
\end{center}

\courssub{Exercices d'application et de transformation}

\begin{exercice}[Mise en forme : ordre des mots]
Remets les mots dans l'ordre pour former une phrase correcte avec complément d'état.
\begin{enumerate}
\item 很 / 说 / 他 / 英语 / 说 / 得 / 好 \\
\item 不 / 跑 / 得 / 他 / 快 / 太 \\
\item 越来越 / 写 / 字 / 她 / 写 / 得 / 工整 (zhěngqí — soigné) \\
\item 怎么样 / 做 / 饭 / 你 / 做 / 得 \\
\item 很 / 跳 / 舞 / 她 / 跳 / 得 / 棒 (bàng — super)
\end{enumerate}
\end{exercice}

\begin{corrige}[Exercice 1]
\begin{enumerate}
\item \zh{他说英语说得很好。} (Tā shuō Yīngyǔ shuō de hěn hǎo.)
\item \zh{他跑得不太快。} (Tā pǎo de bù tài kuài.)
\item \zh{她写字写得越来越工整。} (Tā xiězì xiě de yuè lái yuè zhěngqí.)
\item \zh{你做饭做得怎么样？} (Nǐ zuòfàn zuò de zěnmeyàng?)
\item \zh{她跳舞跳得很棒。} (Tā tiàowǔ tiào de hěn bàng.)
\end{enumerate}
\end{corrige}

\begin{exercice}[Transformation : négation et progression]
Transforme chaque phrase affirmative en phrase négative (avec \zh{不太}) puis en phrase de progression (avec \zh{越来越}).
\begin{enumerate}
\item \zh{他唱歌唱得很好。}
\item \zh{我做饭做得很好吃。}
\item \zh{她跑步跑得很快。}
\end{enumerate}
\end{exercice}

\begin{corrige}[Exercice 2]
\begin{enumerate}
\item Nég. : \zh{他唱歌唱得不太好。} \quad Prog. : \zh{他唱歌唱得越来越好。}
\item Nég. : \zh{我做饭做得不太好吃。} \quad Prog. : \zh{我做饭做得越来越好吃。}
\item Nég. : \zh{她跑步跑得不太快。} \quad Prog. : \zh{她跑步跑得越来越快。}
\end{enumerate}
\end{corrige}

\begin{exercice}[Production guidée : interview écrite]
Complète l'interview ci-dessous en utilisant les indices entre parenthèses.
\begin{textebox}[Interview à compléter]
A：你会做什么？\\
B：我会\zh{踢足球} 和 \zh{画画}。\\
A：你足球踢得\zh{怎么样}？\\
B：我踢得\zh{非常好}，我是校队队员 (xiàoduì duìyuán)。\\
A：那你画画画得\zh{怎么样}？\\
B：我画得\zh{还可以}，但是我画得\zh{不太快}。\\
A：你觉得以后会进步吗？\\
B：当然！我画得\zh{越来越好}！
\end{textebox}
\end{exercice}

\begin{corrige}[Exercice 3]
Voir le texte ci-dessus (réponses en gras). Vérifie la répétition du verbe : \zh{踢足球踢得...}, \zh{画画画得...}.
\end{corrige}

\courssub{Production écrite : le « Mur des talents »}

\begin{exoresolu}[Tâche finale — Présentation d'un camarade]
Rédige un paragraphe de 5 à 6 phrases pour le « Mur des talents ». Contraintes obligatoires :
\begin{itemize}
\item au moins \textbf{trois} compléments d'état différents (verbes distincts) ;
\item au moins \textbf{une} négation (\zh{得不太...} ou \zh{得不...}) ;
\item au moins \textbf{une} phrase avec \zh{越来越} + adjectif ;
\item une phrase de conclusion avec \zh{我觉得} (je pense que) ou \zh{我认为}.
\end{itemize}
\tcblower
\textbf{Production modèle :}

我的同学叫阿里。他会做很多事情。他踢足球踢得非常好，是我们班最好的球员。他说汉语说得越来越流利，但是他说得不太快。他唱歌唱得还可以，可是他跳舞跳得不太好。他做饭做得很好吃，他做的菜大家都喜欢。我觉得阿里是一个很有才华的人！

\emph{Wǒ de tóngxué jiào Ālǐ. Tā huì zuò hěn duō shìqing. Tā tī zúqiú tī de fēicháng hǎo, shì wǒmen bān zuì hǎo de qiúyuán. Tā shuō Hànyǔ shuō de yuè lái yuè liúlì, dànshì tā shuō de bù tài kuài. Tā chànggē chàng de hái kěyǐ, kěshì tā tiàowǔ tiào de bù tài hǎo. Tā zuòfàn zuò de hěn hǎochī, tā zuò de cài dàjiā dōu xǐhuan. Wǒ juéde Ālǐ shì yí ge hěn yǒu cáihuá de rén !}

« Mon camarade s'appelle Ali. Il sait faire beaucoup de choses. Il joue au football extrêmement bien, c'est le meilleur joueur de notre classe. Il parle chinois de plus en plus couramment, mais il ne parle pas très vite. Il chante assez bien, mais il ne danse pas très bien. Il cuisine très bien (ses plats sont délicieux), tout le monde aime ses plats. Je trouve qu'Ali est quelqu'un de très doué ! »

\textbf{Analyse linguistique du modèle :}
\begin{itemize}
\item \zh{踢足球踢得非常好} : objet \zh{足球} après V₁, répétition de \zh{踢} avant \zh{得} + adv. de degré \zh{非常} + adj.
\item \zh{说得越来越流利} : progression \zh{越来越} + adj. \zh{流利} (sans \zh{很}).
\item \zh{说得不太快} / \zh{跳得不太好} : négation standard \zh{得不太} + adj.
\item \zh{做得很好吃} : adj. résultatif \zh{好吃} naturel pour \zh{做饭}.
\item \zh{他做的菜} : structure attributive \zh{的} (celui qui fait / ce qui est fait) — révision Module 3.
\item Conclusion \zh{我觉得...} : ancrage énonciatif personnel.
\end{itemize}
\end{exoresolu}

\courssub{Expression orale et interaction : le débat des talents}

\begin{experience}[Activité de classe : Podium des talents]
\begin{enumerate}
\item \textbf{Par groupes de 4} : lisez les 4 affiches du « Mur des talents » de votre groupe.
\item \textbf{Comparez} les performances à l'aide des structures comparatives et superlatives :
      \begin{itemize}
      \item \zh{A 比 B 唱得好。} (A chante mieux que B.)
      \item \zh{A 跑得最快。} (A court le plus vite.)
      \item \zh{B 没有 A 唱得好。} (B ne chante pas aussi bien que A.)
      \end{itemize}
\item \textbf{Établissez} votre « Podium » (1er, 2e, 3e) pour 3 catégories (ex. chant, foot, cuisine).
\item \textbf{Présentez} au tableau en 3 phrases minimum par groupe. Exemple :
      \zh{我们组足球第一名是阿里，他踢得最好。第二名是保罗，他比马克踢得好。第三名是马克。}
\end{enumerate}
\end{experience}

\courssub{Bilan et points-clés à retenir}

\begin{aretenir}[Récapitulatif : Complément d'état \zh{得}]
\begin{itemize}
\item \textbf{Formule de base} : Sujet + V (+ Objet) + V + \zh{得} + Adjectif.
\item \textbf{Règle d'or} : \textbf{Objet avant \zh{得}} → répétition du verbe. \zh{他汉语说得好} (topicalisation) ou \zh{他说汉语说得好} (standard).
\item \textbf{Négation} : \zh{得不 / 得不太} + Adj. \quad \textbf{Jamais} \zh{不得}.
\item \textbf{Question} : \zh{得怎么样？} (ouverte) / \zh{得...吗？} (fermée).
\item \textbf{Progression} : \zh{得越来越} + Adj. \quad \textbf{Jamais} \zh{越来越很}.
\item \textbf{Distinction} : \zh{得} (compl. d'état) ≠ \zh{地} (adverbe : \zh{快快地跑}) ≠ \zh{的} (attribut : \zh{快的车}).
\end{itemize}
\end{aretenir}

\begin{methode}[Pour réussir l'épreuve écrite (expression / traduction)]
\begin{enumerate}
\item \textbf{Identifie} le verbe principal et son objet éventuel.
\item \textbf{Décide} : si objet → répéter le verbe avant \zh{得} ; si pas d'objet → verbe + \zh{得} direct.
\item \textbf{Choisis} l'adjectif d'évaluation (bon/vite/clair/courant/délicieux...).
\item \textbf{Ajoute} le marqueur de degré (\zh{很}, \zh{非常}, \zh{不太}, \zh{越来越}) selon le sens visé.
\item \textbf{Relis} : vérifie l'absence de \zh{的}/\zh{地} à la place de \zh{得} et l'ordre Objet -- V -- \zh{得} -- Adj.
\end{enumerate}
\end{methode}

\newpage

\coursec{Méthodes et exercices résolus}

\courssub{Méthodes et exercices résolus}

\begin{methode}[Identifier et construire le complément d'état (verbe/adjectif + 得 + état)]
\emph{Objectif :} passer d'une phrase simple (Sujet + Verbe/Adjectif) à une phrase précisant \textbf{comment} l'action est réalisée ou \textbf{à quel degré} l'état se manifeste, en utilisant la particule \cle{得}.
\begin{enumerate}
    \item \textbf{Repérer le noyau :} trouver le verbe d'action (跑、写、吃、学) ou l'adjectif d'état (快、好、累、高兴) principal.
    \item \textbf{Placer \cle{得} :} le placer \textbf{immédiatement après} le verbe ou l'adjectif. Règle d'or : \textbf{jamais de complément d'objet entre le verbe et \cle{得}}.
    \item \textbf{Construire le complément :} ajouter après \cle{得} une structure adjectivale : adverbe de degré + adjectif (很、非常、真、特别 + adjectif), ou adjectif + particule d'intensité (极了、死了、得不得了).
    \item \textbf{Vérifier l'ordre :} Sujet + Verbe/Adjectif + \cle{得} + complément d'état. \textbf{Jamais :} Sujet + 得 + Verbe.
\end{enumerate}
\end{methode}

\begin{aretenir}[Structure type du complément d'état]
\textbf{Sujet + Verbe/Adjectif + 得 + (adverbe de degré) + adjectif (+ particule d'intensité)}\par
\zh{他跑得很快。} \emph{Tā pǎo de hěn kuài.} « Il court très vite. »\par
\zh{这道菜做得真好吃。} \emph{Zhè dào cài zuò de zhēn hǎochī.} « Ce plat est cuisiné vraiment bien (il est délicieux). »
\end{aretenir}

\begin{exoresolu}[Transformation : phrase simple $\rightarrow$ complément d'état]
\emph{Transformez les phrases suivantes en ajoutant le complément d'état indiqué entre parenthèses. Écrivez en caractères, pinyin et français.}
\begin{enumerate}
    \item 他说中文。 (很流利)
    \item 妹妹睡觉。 (很早)
    \item 今天热。 (极了)
    \item 我们学习汉语。 (很认真)
    \item 他跑步。 (不快 — \textbf{négation})
\end{enumerate}
\tcblower
\textbf{Solution détaillée :}
\begin{enumerate}
    \item \textbf{Analyse :} verbe transitif \zh{说} + objet \zh{中文}. L'objet ne peut pas rester entre le verbe et \cle{得} : on \textbf{répète le verbe}. \zh{他说中文说得很流利。} \emph{Tā shuō Zhōngwén shuō de hěn liúlì.} « Il parle chinois très couramment. » Variante par topicalisation : \zh{他中文说得很流利。}
    \item \textbf{Analyse :} \zh{睡觉} est un verbe séparable (verbe-objet). On répète le premier morphème \zh{睡}. \zh{妹妹睡觉睡得很早。} \emph{Mèimei shuìjiào shuì de hěn zǎo.} « Ma petite sœur dort (se couche) très tôt. » Forme courte acceptée : \zh{妹妹睡得很早。}
    \item \textbf{Analyse :} adjectif \zh{热}. Structure : adjectif + 得 + intensité. \zh{今天热极了。} \emph{Jīntiān rè jíle.} « Aujourd'hui, il fait extrêmement chaud. » \textit{Remarque : après un adjectif seul, on préfère souvent 极了 directement (热极了) ; 热得极了 est lourd mais correct.}
    \item \textbf{Analyse :} \zh{学习} + objet \zh{汉语} : répétition du verbe. \zh{我们学习汉语学得很认真。} \emph{Wǒmen xuéxí Hànyǔ xué de hěn rènzhēn.} « Nous étudions le chinois très sérieusement. »
    \item \textbf{Analyse :} négation du complément d'état : Verbe + 得 + \cle{不} + adjectif. \zh{他跑步跑得不快。} \emph{Tā pǎobù pǎo de bù kuài.} « Il ne court pas vite. » Forme courte : \zh{他跑得不快。}
\end{enumerate}
\textbf{Conclusion :} la clé est la place de \cle{得} (juste après le verbe/adjectif) et la gestion des verbes transitifs : répétition du verbe ou topicalisation de l'objet.
\end{exoresolu}

\begin{methode}[Négation, interrogation et emplois courants]
\emph{Objectif :} savoir nier, questionner et utiliser le complément d'état dans des contextes réels (compétences, météo, ressentis).
\begin{enumerate}
    \item \textbf{Négation :} Verbe + 得 + \cle{不 / 不太 / 不怎么} + adjectif. \zh{他跑得不太快。} « Il ne court pas très vite. » Pour un adjectif d'état simple, on nie directement sans \cle{得} : \zh{今天不太热。} « Il ne fait pas très chaud aujourd'hui. »
    \item \textbf{Interrogation :}
    \begin{itemize}
        \item avec \cle{吗} : \zh{他跑得快吗？} « Court-il vite ? »
        \item avec \cle{怎么样} : \zh{他跑得怎么样？} « Comment court-il ? »
        \item avec \cle{多} + adjectif mesurable : \zh{他跑得多快？} « À quelle vitesse court-il ? »
    \end{itemize}
    \item \textbf{Emplois courants :}
    \begin{itemize}
        \item compétences : \zh{法语说得很标准} « parler un français très standard » ; \zh{电脑修得真好} « réparer très bien les ordinateurs » ;
        \item vie quotidienne : \zh{睡得很香} « dormir profondément » ; \zh{吃得很饱} « manger à sa faim » ;
        \item ressenti : \zh{冷得受不了} « avoir un froid insupportable » ; \zh{累得动不了} « être fatigué au point de ne plus pouvoir bouger ».
    \end{itemize}
\end{enumerate}
\end{methode}

\begin{attention}[Ne pas confondre complément d'état et complément de possibilité]
\zh{他做得不好。} (complément d'\textbf{état} : « il ne le fait pas bien ») $\neq$ \zh{他做不好。} (complément de \textbf{possibilité} : « il n'est pas capable de bien le faire »). La présence de \cle{得} après le verbe change le sens : \zh{做得好} « peut bien le faire » (possibilité) se distingue de \zh{做得好} « le fait bien » (état) uniquement par le contexte et la négation (\zh{做不好} / \zh{做得不好}).
\end{attention}

\begin{exoresolu}[Négation, question et mise en situation (Yaoundé)]
\emph{Partie A : mettez les phrases à la forme négative (négation du degré ou de la manière).}
\begin{enumerate}
    \item 这个摊位的恩多莱做得很好吃。 (恩多莱 Ènduōlái : ndolé)
    \item 我在杜阿拉的交通堵塞中等了很久。 (堵塞 dǔsè ; 等 děng)
    \item 他跑得很快。
\end{enumerate}
\emph{Partie B : posez la question sur le complément d'état (avec \cle{怎么样} ou \cle{多}).}
\begin{enumerate}
    \setcounter{enumi}{3}
    \item 这位中国老师汉语教得很清楚。 (教 jiāo)
    \item 昨晚的足球赛踢得很精彩。 (踢 tī ; 精彩 jīngcǎi)
\end{enumerate}
\tcblower
\textbf{Solution détaillée :}
\textbf{Partie A (négation) :}
\begin{enumerate}
    \item \zh{这个摊位的恩多莱做得不太好吃。} \emph{Zhège tānwèi de Ènduōlái zuò de bù tài hǎochī.} « Le ndolé de ce stand n'est pas très bon. » Structure : verbe + 得 + 不太 + adjectif.
    \item \zh{我在杜阿拉的交通堵塞中等得不太久。} \emph{Wǒ zài Dùālā de jiāotōng dǔsè zhōng děng de bù tài jiǔ.} « Je n'ai pas attendu trop longtemps dans les embouteillages de Douala. » Verbe intransitif : pas de répétition.
    \item \zh{他跑得不快。} \emph{Tā pǎo de bù kuài.} « Il ne court pas vite. »
\end{enumerate}
\textbf{Partie B (interrogation) :}
\begin{enumerate}
    \setcounter{enumi}{3}
    \item \zh{这位中国老师汉语教得怎么样？} \emph{Zhè wèi Zhōngguó lǎoshī Hànyǔ jiāo de zěnmeyàng?} « Comment ce professeur chinois enseigne-t-il le chinois ? »
    \item \zh{昨晚的足球赛踢得怎么样？} \emph{Zuówǎn de zúqiúsài tī de zěnmeyàng?} « Comment le match de foot d'hier soir a-t-il été joué ? » \textit{Remarque : \cle{多} + adjectif convient surtout aux adjectifs mesurables (快、久、高、贵、远) ; avec 精彩, \cle{怎么样} est naturel.}
\end{enumerate}
\textbf{Conclusion :} la négation porte sur l'adjectif du complément (\cle{不/不太} placé après \cle{得}) ; l'interrogation se fait avec \cle{吗}, \cle{怎么样} ou \cle{多} + adjectif mesurable.
\end{exoresolu}

\begin{methode}[Différencier les trois « de » : 得, 的, 地]
\emph{Objectif :} éviter la confusion majeure qui coûte des points à l'examen. La prononciation est identique (\emph{de}, ton neutre), mais la fonction syntaxique diffère.
\begin{enumerate}
    \item \textbf{\cle{得} (complément d'état/degré) :} se place \textbf{après} un verbe ou un adjectif : Verbe/Adjectif + 得 + adjectif. \zh{他跑得快。}
    \item \textbf{\cle{的} (attribut/possession) :} se place \textbf{avant} un nom : adjectif/proposition + 的 + nom. \zh{跑得快的人} « la personne qui court vite ».
    \item \textbf{\cle{地} (complément circonstanciel de manière) :} se place \textbf{avant} un verbe : adjectif (souvent redoublé) + 地 + verbe. \zh{他慢慢地走。} « Il marche lentement. »
\end{enumerate}
\end{methode}

\begin{attention}[Piège classique : les trois « de »]
\begin{itemize}
    \item \zh{他跑得快。} « Il court vite » (complément d'état, après le verbe).
    \item \zh{跑得快的人很多。} « Les gens qui courent vite sont nombreux » (attribut avec \cle{的}, avant le nom).
    \item \zh{他快快地跑。} « Il court vite » (manière avec \cle{地}, avant le verbe ; style plutôt écrit).
    \item \textbf{Erreurs fatales :} \zh{他快的跑} (incorrect), \zh{跑地快} (incorrect), \zh{他说得中文很好} (incorrect : objet coincé entre le verbe et \cle{得}).
\end{itemize}
\end{attention}

\begin{exoresolu}[Correction d'erreurs « les trois de » et traduction]
\emph{Corrigez les phrases fautives (une erreur par phrase) et traduisez.}
\begin{enumerate}
    \item 这本书写地很有意思。
    \item 这个蛋糕吃的很好吃。
    \item 他中文说的很流利。
    \item 请你慢得走。
    \item 他说得中文很好。
\end{enumerate}
\tcblower
\textbf{Solution détaillée :}
\begin{enumerate}
    \item \textbf{Erreur :} \cle{地} après le verbe. \textbf{Correction :} \zh{这本书写得很有意思。} \emph{Zhè běn shū xiě de hěn yǒuyìsi.} « Ce livre est écrit de façon très intéressante. »
    \item \textbf{Erreur :} \cle{的} après le verbe. \textbf{Correction :} \zh{这个蛋糕吃得很好吃。} est lourd ; forme naturelle : \zh{这个蛋糕很好吃。} \emph{Zhège dàngāo hěn hǎochī.} « Ce gâteau est très bon. » Pour garder le complément d'état sur l'action : \zh{这个蛋糕做得很好吃。} « Ce gâteau est fait très bien (très réussi). »
    \item \textbf{Erreur :} \cle{的} après le verbe. \textbf{Correction :} \zh{他中文说得很流利。} \emph{Tā Zhōngwén shuō de hěn liúlì.} « Il parle chinois couramment. »
    \item \textbf{Erreur :} \cle{得} avant le verbe. \textbf{Correction :} \zh{请你慢慢地走。} \emph{Qǐng nǐ mànmàn de zǒu.} « Marche lentement, s'il te plaît. » Règle : adjectif redoublé + \cle{地} + verbe.
    \item \textbf{Erreur :} objet \zh{中文} placé entre le verbe et \cle{得}. \textbf{Correction :} \zh{他说中文说得很好。} ou \zh{他中文说得很好。} \emph{Tā shuō Zhōngwén shuō de hěn hǎo.} « Il parle très bien chinois. »
\end{enumerate}
\textbf{Conclusion :} réflexe de position — avant un nom $\rightarrow$ \cle{的} ; avant un verbe $\rightarrow$ \cle{地} ; après un verbe/adjectif $\rightarrow$ \cle{得}.
\end{exoresolu}

\begin{methode}[Intégrer le complément d'état dans une production écrite courte]
\emph{Objectif :} rédiger un paragraphe cohérent (5 à 6 phrases) pour l'épreuve d'expression écrite, en utilisant au moins 3 compléments d'état variés (action, ressenti, compétence).
\begin{enumerate}
    \item \textbf{Analyser le sujet :} par exemple « Décrivez une journée de marché, une expérience de stage, un match de foot. »
    \item \textbf{Sélectionner les verbes et adjectifs clés :} 卖 (vendre), 买 (acheter), 热 (chaud), 累 (fatigué), 开心 (content), 说 (parler), 做 (faire/cuisiner).
    \item \textbf{Construire les phrases « phares » avec \cle{得} :}
    \begin{itemize}
        \item action : \zh{摊主卖得很起劲。} « Le vendeur vend avec entrain. »
        \item degré : \zh{东西卖得真便宜。} « Les marchandises se vendent vraiment bon marché. »
        \item ressenti : \zh{热得受不了。} « Il fait une chaleur insupportable. » ; \zh{累得腰酸背痛。} « Être épuisé, avoir mal au dos et aux jambes. »
    \end{itemize}
    \item \textbf{Relier par des connecteurs :} \zh{因为……所以……}, \zh{虽然……但是……}, \zh{首先……然后……最后……}.
    \item \textbf{Relire :} vérifier \cle{得} (et non \cle{的}/\cle{地}), l'ordre des mots (objet topicalisé ou verbe répété) et les tons.
\end{enumerate}
\end{methode}

\begin{exoresolu}[Rédaction guidée : « Une journée de stage à Douala »]
\emph{Consigne :} écrivez un court texte (60 à 80 caractères) en chinois. Situation : vous êtes stagiaire dans une entreprise chinoise à Douala. Décrivez votre journée : le travail, le midi, le soir. Utilisez au moins 3 compléments d'état distincts.
\tcblower
\textbf{Proposition de rédaction (modèle) :}\par
\zh{今天我在公司实习。早上中文邮件写得很慢，但是老板看得很仔细。中午食堂的饭做得真好吃，我吃得很饱。下午会议开得很久，我听得有点累。下班回家，路上堵得厉害，但是我今天学得很开心。}\par
\emph{Jīntiān wǒ zài gōngsī shíxí. Zǎoshang Zhōngwén yóujiàn xiě de hěn màn, dànshì lǎobǎn kàn de hěn zǐxì. Zhōngwǔ shítáng de fàn zuò de zhēn hǎochī, wǒ chī de hěn bǎo. Xiàwǔ huìyì kāi de hěn jiǔ, wǒ tīng de yǒudiǎn lèi. Xiàbān huíjiā, lùshang dǔ de lìhai, dànshì wǒ jīntiān xué de hěn kāixīn.}\par
« Aujourd'hui, je fais un stage dans l'entreprise. Le matin, j'ai écrit les courriels en chinois très lentement, mais le patron les a lus très attentivement. À midi, la cantine a cuisiné vraiment bien, j'ai mangé à ma faim. L'après-midi, la réunion a duré très longtemps, j'ai écouté un peu fatigué. En rentrant du travail, la route était fortement bouchée, mais j'ai appris aujourd'hui avec beaucoup de joie. »\par
\textbf{Analyse grammaticale (points forts) :}
\begin{itemize}
    \item 6 compléments d'état corrects : \zh{写得很慢} (manière), \zh{看得很仔细} (manière), \zh{做得真好吃} (résultat), \zh{吃得很饱} (résultat), \zh{开得很久} (durée), \zh{听得有点累} (état simultané), \zh{堵得厉害} (degré), \zh{学得很开心} (état d'esprit).
    \item Vocabulaire professionnel : \zh{实习} (stage), \zh{邮件} (courriel), \zh{老板} (patron), \zh{食堂} (cantine), \zh{会议} (réunion), \zh{下班} (finir le travail), \zh{堵} (bouché).
    \item Connecteur : \zh{但是} (mais).
\end{itemize}
\textbf{Appréciation :} phrases variées, structures ciblées maîtrisées, vocabulaire contextuel (Cameroun/Chine), aucune faute sur \cle{得}.
\end{exoresolu}

\begin{methode}[Préparer l'oral : exposé et jeu de rôle « Évaluation de compétences »]
\emph{Objectif :} réussir l'épreuve orale (description, jeu de rôle, exposé) en utilisant le complément d'état pour \textbf{évaluer, qualifier, nuancer}.
\begin{enumerate}
    \item \textbf{Structure type de réponse :} Sujet + Verbe + 得 + appréciation. \zh{这个工作做得不错。} « Ce travail est bien fait. »
    \item \textbf{Jeu de rôle type :} l'examinateur demande : \zh{你觉得这个实习生表现怎么样？} « Comment trouves-tu la prestation de ce stagiaire ? » Réponse structurée : \zh{我觉得他中文说得很流利，工作做得很认真，但是来得有点晚。} « Je trouve qu'il parle chinois couramment, qu'il travaille sérieusement, mais qu'il arrive un peu tard. »
    \item \textbf{Exposé d'une minute :} choisir un thème (mon meilleur professeur, mon plat préféré, mon équipe de foot) et préparer 5 phrases avec \cle{得}.
\end{enumerate}
\end{methode}

\begin{center}
\begin{tabularx}{\linewidth}{|X|X|X|X|}
\hline
\rowcolor{popblueL} Caractères & Pinyin & Nature & Traduction \\
\hline
很好 / 很棒 & hěn hǎo / hěn bàng & appréciation positive & très bien / formidable \\
\hline
不太好 / 很差 & bù tài hǎo / hěn chà & appréciation négative & pas très bien / mauvais \\
\hline
很专业 & hěn zhuānyè & appréciation positive & très professionnel \\
\hline
很标准 & hěn biāozhǔn & appréciation positive & très standard (langue) \\
\hline
很流利 & hěn liúlì & appréciation positive & très courant (langue) \\
\hline
很认真 & hěn rènzhēn & appréciation positive & très sérieux \\
\hline
马马虎虎 & mǎmǎhūhū & appréciation faible & à la va-vite, bof \\
\hline
\end{tabularx}
\end{center}

\begin{exoresolu}[Simulation orale : jeu de rôle « Entretien d'embauche »]
\emph{Consigne :} l'examinateur (recruteur chinois à Yaoundé) demande : \zh{请介绍一下你的中文水平和工作能力。} « Présentez votre niveau de chinois et vos capacités de travail. » Répondez en 4 à 5 phrases en utilisant le complément d'état.
\tcblower
\textbf{Modèle de réponse :}\par
\zh{我学中文学了三年了。汉字写得比较标准，语法掌握得不错，口语说得还可以，但是听力理解得不太快。工作方面，电脑操作得很熟练，做事做得很认真，从来不马虎。希望能为贵公司贡献力量。}\par
\emph{Wǒ xué Zhōngwén xué le sān nián le. Hànzì xiě de bǐjiào biāozhǔn, yǔfǎ zhǎngwò de búcuò, kǒuyǔ shuō de hái kěyǐ, dànshì tīnglì lǐjiě de bù tài kuài. Gōngzuò fāngmiàn, diànnǎo cāozuò de hěn shúliàn, zuòshì zuò de hěn rènzhēn, cónglái bù mǎmǎhūhū. Xīwàng néng wèi guì gōngsī gòngxiàn lìliang.}\par
« J'apprends le chinois depuis trois ans. J'écris les caractères de façon assez standard, je maîtrise la grammaire plutôt bien, je parle correctement, mais je comprends l'oral pas très vite. Côté travail, j'utilise l'ordinateur très habilement, je travaille très sérieusement, jamais à la va-vite. J'espère pouvoir apporter ma contribution à votre entreprise. »\par
\textbf{Analyse des structures avec \cle{得} :}
\begin{enumerate}
    \item \zh{写得比较标准} : verbe + 得 + adverbe + adjectif (compétence écrite).
    \item \zh{掌握得不错} : verbe + 得 + adjectif (maîtrise).
    \item \zh{说得还可以} : verbe + 得 + adverbe + adjectif (nuance modérée).
    \item \zh{理解得不太快} : verbe + 得 + négation + adverbe + adjectif (négation correcte du degré).
    \item \zh{操作得很熟练} : verbe + 得 + adverbe + adjectif (compétence technique).
\end{enumerate}
\textbf{Conseil oral :} prononcer \cle{得} en ton neutre, léger et court, sans pause avant lui : \emph{xiě-de-biāozhǔn}.
\end{exoresolu}

\begin{methode}[Traduction thème (français $\rightarrow$ chinois) : pièges de transposition]
\emph{Objectif :} traduire des phrases françaises contenant des adverbes en \emph{-ment}, des tournures avec « de » ou des relatives, en structures chinoises appropriées.
\begin{enumerate}
    \item \textbf{Adverbe en \emph{-ment} (manière) $\rightarrow$ Verbe + 得 + adjectif :} « Il conduit prudemment. » $\rightarrow$ \zh{他开车开得很小心。} (verbe répété) ou, style écrit, \zh{他小心地开车。} À l'examen, privilégier Verbe + 得 + adjectif.
    \item \textbf{« Facile/difficile à + verbe » $\rightarrow$ Sujet patient + Verbe + 得 + adjectif :} « Ce livre est facile à lire. » $\rightarrow$ \zh{这本书看得很容易。} ou plus naturel : \zh{这本书很容易看。}
    \item \textbf{« Trop... pour » $\rightarrow$ adjectif + 得 + complément de résultat potentiel :} « Trop lourd pour être porté. » $\rightarrow$ \zh{重得提不起来。} (structure avancée, à reconnaître).
    \item \textbf{Qualité permanente $\rightarrow$ attribut avec \cle{的} :} « un gâteau très bon » $\rightarrow$ \zh{一个很好吃的蛋糕} (qualité intrinsèque, pas de \cle{得}).
\end{enumerate}
\end{methode}

\begin{attention}[Calque du français : « Il parle bien français »]
En français, l'adverbe suit le verbe ; en chinois, on ne peut pas dire \zh{他说法语好} dans ce sens. Formes correctes : \zh{他法语说得好。} (topicalisation), \zh{他说法语说得好。} (verbe répété), \zh{他流利地说法语。} (manière avec \cle{地}, style écrit).
\end{attention}

\begin{exoresolu}[Thème : traduction ciblée (contexte Cameroun)]
\emph{Traduisez en chinois (caractères + pinyin).}
\begin{enumerate}
    \item Le ndolé de ma mère est cuisiné très bien (délicieux).
    \item Les routes de Yaoundé sont très bouchées le matin.
    \item Ce joueur court très vite.
    \item J'ai dormi très tard hier soir.
    \item Ce problème est difficile à comprendre.
\end{enumerate}
\tcblower
\textbf{Solution détaillée :}
\begin{enumerate}
    \item \zh{我妈妈做恩多莱做得真好吃。} \emph{Wǒ māma zuò Ènduōlái zuò de zhēn hǎochī.} (complément d'état sur l'action de cuisiner). Variante attributive : \zh{我妈妈做的恩多莱真好吃。} « Le ndolé que fait ma mère est vraiment bon. »
    \item \zh{早上雅温得的路堵得很厉害。} \emph{Zǎoshang Yǎwēndé de lù dǔ de hěn lìhai.} Adjectif \zh{堵} (bouché) + 得 + degré.
    \item \zh{这个队员跑得很快。} \emph{Zhège duìyuán pǎo de hěn kuài.} « Ce joueur court très vite. »
    \item \zh{昨晚我睡得很晚。} \emph{Zuówǎn wǒ shuì de hěn wǎn.} Verbe séparable \zh{睡觉} : on répète \zh{睡} ; « tard » = \zh{晚}.
    \item \zh{这个问题很难理解。} \emph{Zhège wèntí hěn nán lǐjiě.} ou avec complément d'état : \zh{这个问题理解起来很难。} « Ce problème est difficile à comprendre. »
\end{enumerate}
\textbf{Conclusion :} en thème, choisir la structure selon le sens : qualité permanente $\rightarrow$ attribut avec \cle{的} ; performance ou résultat d'une action $\rightarrow$ complément d'état avec \cle{得} ; manière (style écrit) $\rightarrow$ \cle{地} avant le verbe.
\end{exoresolu}

\begin{attention}[Les 5 erreurs qui coûtent des points à l'examen — Leçon 32]
\begin{enumerate}
    \item \textbf{Confusion \cle{得} / \cle{的} / \cle{地}.} Réflexe : avant un nom $\rightarrow$ \cle{的} (\zh{白色的纸}) ; avant un verbe $\rightarrow$ \cle{地} (\zh{认真地写}) ; après un verbe ou un adjectif $\rightarrow$ \cle{得} (\zh{写得认真}).
    \item \textbf{Oubli de la répétition du verbe.} Avec un verbe transitif suivi d'un objet : \zh{我说中文说得好。} (verbe répété) ou \zh{我中文说得好。} (objet topicalisé). Verbes séparables (\zh{睡觉、游泳、吃饭}) : répéter le premier morphème : \zh{睡觉睡得很晚} ou \zh{睡得很晚}.
    \item \textbf{Objet coincé entre le verbe et \cle{得}.} Faute : \zh{他说得中文很好。} Correct : \zh{他说中文说得很好。} ou \zh{中文他说得很好。}
    \item \textbf{Négation mal formée (état contre possibilité).} État : \zh{他做得不好。} « Il ne le fait pas bien. » Possibilité : \zh{他做不好。} « Il n'arrive pas à bien le faire. » Ne pas mélanger les deux structures.
    \item \textbf{Tons et caractères proches.} \zh{得} / \zh{的} / \zh{地} se prononcent tous \emph{de} (ton neutre) : l'oreille ne les distingue pas, l'écrit si. \zh{累} (\emph{lèi}, fatigué) contre \zh{雷} (\emph{léi}, tonnerre) ; \zh{厉害} (\emph{lìhai}, redoutable, fort) : écrire le bon caractère, c'est valider le point.
\end{enumerate}
\end{attention}

\newpage

\coursec{Exercices}

\courssub{Je vérifie mes connaissances}

\begin{exercice}[QCM : Structure et particules]
Choisissez la bonne réponse parmi A, B, C, D.
\begin{enumerate}
\item La structure de base du complément d'état est : \textbf{Verbe + 得 + Adjectif}. Quel mot-outil introduit le complément d'état ?
      \begin{itemize}
      \item A. \zh{的}
      \item B. \zh{地}
      \item C. \zh{得}
      \item D. \zh{了}
      \end{itemize}
\item Dans la phrase \zh{他跑得很快} (\emph{Tā pǎo de hěn kuài}, « Il court très vite »), quel est le rôle de \zh{很} ?
      \begin{itemize}
      \item A. Verbe principal
      \item B. Adverbe de degré modifiant l'adjectif \zh{快}
      \item C. Complément d'objet
      \item D. Particule de phrase
      \end{itemize}
\item Laquelle des phrases suivantes est \textbf{incorrecte} (erreur de structure ou de particule) ?
      \begin{itemize}
      \item A. \zh{她唱歌唱得很好。}
      \item B. \zh{他写字写得很认真。}
      \item C. \zh{我们学习学得很努力地。}
      \item D. \zh{我睡得很晚。}
      \end{itemize}
\item Pour former la question sur le complément d'état avec \zh{怎么样}, on dit :
      \begin{itemize}
      \item A. \zh{他工作怎么样？}
      \item B. \zh{他工作得怎么样？}
      \item C. \zh{他怎么样工作？}
      \item D. \zh{他得怎么样工作？}
      \end{itemize}
\end{enumerate}
\end{exercice}

\begin{exercice}[Vrai ou Faux ? Justifiez]
Déterminez si les affirmations sont vraies (V) ou fausses (F) et justifiez brièvement en français.
\begin{enumerate}
\item Le complément d'état indique \emph{comment} une action est réalisée ou \emph{dans quel état} elle place le sujet.
\item Dans la structure \zh{Sujet + Verbe + Objet + Verbe + 得 + Complément}, le verbe est \textbf{répété} si le verbe principal est transitif et possède un objet.
\item La négation du complément d'état se fait toujours avec \zh{没} placé devant \zh{得}.
\item \zh{得}、\zh{的} et \zh{地} se prononcent tous \emph{de} (ton neutre) mais ont des fonctions grammaticales différentes.
\end{enumerate}
\end{exercice}

\begin{exercice}[Vocabulaire : Adjectifs courants du complément d'état]
Complétez le tableau en indiquant la nature grammaticale et la traduction de chaque mot.
\begin{center}
\begin{tabularx}{\linewidth}{|X|X|X|X|}
\hline
\rowcolor{popblueL} \textbf{Caractères} & \textbf{Pinyin} & \textbf{Nature} & \textbf{Traduction} \\ \hline
\zh{认真} & rènzhēn & & \\ \hline
\zh{仔细} & zǐxì & & \\ \hline
\zh{马虎} & mǎhu & & \\ \hline
\zh{流利} & liúlì & & \\ \hline
\zh{累} & lèi & & \\ \hline
\zh{清楚} & qīngchu & & \\ \hline
\end{tabularx}
\end{center}
\textit{Traductions proposées (dans le désordre) : sérieux, appliqué ; attentif, minutieux ; négligent ; courant (pour une langue) ; fatigué ; clair, distinct.}
\end{exercice}

\begin{exercice}[Pinyin : Association caractères / prononciation]
Reliez chaque groupe de caractères à son pinyin correct (tons inclus).
\begin{enumerate}
\item \zh{开车} \hspace{2cm} A. \emph{shuō Hànyǔ}
\item \zh{做饭} \hspace{2cm} B. \emph{kāichē}
\item \zh{说汉语} \hspace{2cm} C. \emph{zuòfàn}
\item \zh{写字} \hspace{2cm} D. \emph{xiězì}
\item \zh{睡觉} \hspace{2cm} E. \emph{shuìjiào}
\end{enumerate}
\end{exercice}

\courssub{Je m'entraîne}

\begin{exercice}[Traduction : Français $\rightarrow$ Chinois]
Traduisez les phrases suivantes en chinois (caractères + pinyin). Utilisez la structure du complément d'état (verbe répété si nécessaire).
\begin{enumerate}
\item Mon frère cadet conduit prudemment. (prudent = \zh{小心} xiǎoxīn)
\item Elle cuisine très bien.
\item Nous apprenons le chinois sérieusement. (sérieux = \zh{认真} rènzhēn)
\item Il dort très tard le soir.
\item Le professeur explique clairement. (clair = \zh{清楚} qīngchu)
\end{enumerate}
\end{exercice}

\begin{exercice}[Transformation : Négation, Question, \zh{怎么样}]
Partez de la phrase de base : \zh{妹妹跳舞跳得很美} (\emph{Mèimei tiàowǔ tiào de hěn měi}, « Ma petite sœur danse très bien »).
\begin{enumerate}
\item Transformez en phrase \textbf{négative} (avec \zh{不}).
\item Transformez en question \textbf{fermée} (avec \zh{吗}).
\item Transformez en question \textbf{ouverte} avec \zh{怎么样}.
\item Transformez en phrase \textbf{exclamative} avec l'adverbe \zh{真} (zhēn).
\end{enumerate}
\end{exercice}

\begin{exercice}[Texte à trous : Particules \zh{的} / \zh{地} / \zh{得}]
Complétez le texte ci-dessous avec la particule appropriée (\zh{的}、\zh{地} ou \zh{得}).
\begin{textebox}[Journal intime de Xiao Ming]
今天早上，我哥哥起\_\_\_\_ 很早。他穿衣服穿\_\_\_\_ 很快，吃饭吃\_\_\_\_ 也很快，然后开车开\_\_\_\_ 很小心，因为路上车很多。他工作\_\_\_\_ 很努力，下班回家回\_\_\_\_ 很累。但是他觉得汉语很有意思，所以晚上还学习\_\_\_\_ 很认真。\\
\emph{Jīntiān zǎoshang, wǒ gēge qǐ \_\_\_\_ hěn zǎo. Tā chuān yīfu chuān \_\_\_\_ hěn kuài, chīfàn chī \_\_\_\_ yě hěn kuài, ránhòu kāichē kāi \_\_\_\_ hěn xiǎoxīn, yīnwèi lùshang chē hěn duō. Tā gōngzuò \_\_\_\_ hěn nǔlì, xiàbān huí jiā huí \_\_\_\_ hěn lèi. Dànshì tā juéde Hànyǔ hěn yǒu yìsi, suǒyǐ wǎnshang hái xuéxí \_\_\_\_ hěn rènzhēn.}\\
« Ce matin, mon grand frère s'est levé très tôt. Il s'est habillé vite, a mangé vite, puis a conduit prudemment parce qu'il y avait beaucoup de voitures sur la route. Il travaille dur, et il rentre du travail très fatigué. Mais il trouve le chinois intéressant, alors le soir il étudie encore sérieusement. »
\end{textebox}
\end{exercice}

\begin{exercice}[Remise en ordre : Ordre des mots]
Remettez les mots dans l'ordre correct pour former une phrase grammaticale (caractères + pinyin).
\begin{enumerate}
\item \zh{很 / 说 / 他 / 中文 / 说 / 得 / 好}
\item \zh{跑 / 得 / 弟弟 / 步 / 很快 / 跑}
\item \zh{写 / 字 / 写 / 我 / 得 / 不 / 很 / 好}
\item \zh{听 / 老师 / 课 / 听 / 得 / 同学们 / 很 / 仔细}
\item \zh{怎么样 / 你 / 睡 / 觉 / 睡 / 得}
\end{enumerate}
\end{exercice}

\courssub{Je me prépare au Bac}

\begin{exercice}[Compréhension de texte et analyse linguistique]
Lisez le texte suivant, puis répondez aux questions.

\begin{textebox}[Témoignage : Apprendre le chinois à Yaoundé]
我叫保罗，住在喀麦隆的首都雅温得。我在大学学习汉语。我的汉语老师姓李，她讲课讲得非常清楚，所以我们听得都懂。课后，我经常和中国留学生练习口语。他们说话说得很快，开始我听不懂，现在听得懂多了。我也练习写汉字，虽然写得不太好，但是我写得很认真。我想明年去北京留学。\\
\emph{Wǒ jiào Bǎoluó, zhù zài Kāmàilóng de shǒudū Yǎwēndé. Wǒ zài dàxué xuéxí Hànyǔ. Wǒ de Hànyǔ lǎoshī xìng Lǐ, tā jiǎngkè jiǎng de fēicháng qīngchu, suǒyǐ wǒmen tīng de dōu dǒng. Kè hòu, wǒ jīngcháng hé Zhōngguó liúxuéshēng liànxí kǒuyǔ. Tāmen shuōhuà shuō de hěn kuài, kāishǐ wǒ tīng bu dǒng, xiànzài tīng de dǒng duō le. Wǒ yě liànxí xiě Hànzì, suīrán xiě de bú tài hǎo, dànshì wǒ xiě de hěn rènzhēn. Wǒ xiǎng míngnián qù Běijīng liúxué.}\\
« Je m'appelle Paul et j'habite à Yaoundé, la capitale du Cameroun. J'étudie le chinois à l'université. Mon professeur de chinois s'appelle Li ; elle explique ses cours très clairement, donc nous comprenons tout. Après les cours, je pratique souvent l'oral avec des étudiants chinois. Ils parlent très vite : au début je ne comprenais pas, maintenant je comprends beaucoup mieux. Je m'entraîne aussi à écrire les caractères ; même si je n'écris pas très bien, j'écris avec application. Je voudrais partir étudier à Pékin l'année prochaine. »
\end{textebox}

\textbf{Questions de compréhension (en français) :}
\begin{enumerate}
\item Où habite Paul et qu'étudie-t-il à l'université ?
\item Pourquoi Paul et ses camarades comprennent-ils bien le cours de M\textsuperscript{me} Li ?
\item Quelle difficulté Paul rencontre-t-il au début avec les étudiants chinois ? Comment la situation a-t-elle évolué ?
\item Quel est le projet de Paul pour l'année prochaine ?
\end{enumerate}

\textbf{Questions de langue :}
\begin{enumerate}
\setcounter{enumi}{4}
\item Relevez \textbf{trois} compléments d'état dans le texte. Pour chacun, indiquez : le verbe principal, la particule, le complément et la traduction.
\item Dans la phrase \zh{写得不太好}, analysez la structure de la négation. Pourquoi \zh{不} et non \zh{没} ?
\item Expliquez la différence de structure entre \zh{讲课讲得清楚} et \zh{听得懂} (complément de résultat potentiel).
\end{enumerate}
\end{exercice}

\begin{exercice}[Thème et Version]
\textbf{A. Thème (Français $\rightarrow$ Chinois) :} Traduisez en chinois (caractères + pinyin). Attention aux compléments d'état et à la répétition du verbe.
\begin{enumerate}
\item La professeure explique très clairement, nous comprenons bien.
\item Mon père conduit vite mais prudemment.
\item Ils font leurs devoirs soigneusement, donc ils les font très bien.
\item Elle parle français couramment, mais elle écrit assez lentement.
\end{enumerate}

\textbf{B. Version (Chinois $\rightarrow$ Français) :} Traduisez le texte suivant en français.
\begin{textebox}[Message de Wang Wei]
小张是我的同事。他开车开得真快，我坐在车上很害怕。上个星期，他请我吃饭。他做饭做得真好吃，特别是红烧肉。我们聊得很开心，喝得有点多。第二天他头疼得厉害，上班迟到了半个小时。\\
\emph{Xiǎo Zhāng shì wǒ de tóngshì. Tā kāichē kāi de zhēn kuài, wǒ zuò zài chēshang hěn hàipà. Shàng ge xīngqī, tā qǐng wǒ chīfàn. Tā zuòfàn zuò de zhēn hǎochī, tèbié shì hóngshāoròu. Wǒmen liáo de hěn kāixīn, hē de yǒudiǎn duō. Dì èr tiān tā tóuténg de lìhai, shàngbān chídào le bàn ge xiǎoshí.}\\
\end{textebox}
\end{exercice}

\begin{exercice}[Expression écrite : Mon entourage]
\textbf{Consigne :} Rédigez un paragraphe d'environ \textbf{80 à 100 caractères chinois} pour décrire comment un membre de votre famille (père, mère, frère, sœur, grand-père...) effectue ses activités quotidiennes ou professionnelles.

\textbf{Contraintes obligatoires :}
\begin{itemize}
\item Utilisez \textbf{au moins trois} compléments d'état distincts (structure V + 得 + Adj).
\item Intégrez au moins \textbf{trois} des verbes d'activité suivants : \zh{做饭} (cuisiner), \zh{开车} (conduire), \zh{说/讲} (parler), \zh{写} (écrire), \zh{工作} (travailler), \zh{学习} (étudier).
\item Utilisez des adverbes de degré (\zh{很}、\zh{真}、\zh{特别}、\zh{比较}、\zh{不太}).
\item Soignez l'ordre des mots (sujet + verbe (+ objet) + verbe + 得 + complément) et la ponctuation chinoise.
\end{itemize}

\textbf{Amorce suggérée :} \zh{我的爸爸……}
\end{exercice}



\newpage

\coursec{Corrigés des exercices}

\begin{corrige}[Exercice 1 — QCM : Structure et particules]
\begin{enumerate}
\item \textbf{C.} \zh{得} (\emph{de}, ton neutre) introduit le complément d'état. \zh{的} marque la détermination (complément du nom, possessif) et \zh{地} marque le complément circonstanciel de manière placé \emph{avant} le verbe.
\item \textbf{B.} \zh{很} est un adverbe de degré qui modifie l'adjectif \zh{快} ; il est presque obligatoire devant un adjectif monosyllabique employé comme complément, sinon la phrase prend une valeur comparative.
\item \textbf{C.} La phrase correcte est \zh{我们学习学得很努力。} : après \zh{得}, on met un adjectif ou un groupe adjectival, jamais la particule \zh{地} (qui ne s'emploie qu'avant un verbe).
\item \textbf{B.} \zh{他工作得怎么样？} (\emph{Tā gōngzuò de zěnmeyàng ?}, « Comment travaille-t-il ? »). \zh{怎么样} occupe la place du complément, après \zh{得}.
\end{enumerate}
\textbf{Barème indicatif :} 1 point par réponse juste, soit 4 points.
\end{corrige}

\begin{corrige}[Exercice 2 — Vrai ou Faux ? Justifiez]
\begin{enumerate}
\item \textbf{Vrai.} Le complément d'état évalue la manière dont l'action se déroule (\zh{他跑得很快}, « il court vite ») ou l'état qui en résulte (\zh{他累得说不出话来}, « il est fatigué au point de ne plus pouvoir parler »).
\item \textbf{Vrai.} Quand le verbe a un objet, on répète le verbe : \zh{她唱歌唱得很好} (\emph{Tā chànggē chàng de hěn hǎo}, « Elle chante bien »). On peut aussi thématiser l'objet en tête de phrase, mais la répétition reste la structure de base à maîtriser.
\item \textbf{Faux.} La négation du complément d'état se fait avec \zh{不} placé \emph{après} \zh{得}, devant l'adjectif : \zh{他跑得不快} (« Il ne court pas vite »). \zh{没} nie l'accomplissement d'une action, pas un état.
\item \textbf{Vrai.} Les trois particules se prononcent \emph{de} (ton neutre) : \zh{的} = détermination du nom, \zh{地} = manière avant le verbe, \zh{得} = complément d'état après le verbe. À l'écrit, il ne faut jamais les confondre.
\end{enumerate}
\textbf{Barème indicatif :} 0,5 point pour le jugement + 0,5 point pour la justification, soit 4 points.
\end{corrige}

\begin{corrige}[Exercice 3 — Vocabulaire : Adjectifs courants du complément d'état]
\begin{center}
\begin{tabularx}{\linewidth}{|X|X|X|X|}
\hline
\rowcolor{popblueL} \textbf{Caractères} & \textbf{Pinyin} & \textbf{Nature} & \textbf{Traduction} \\ \hline
\zh{认真} & rènzhēn & Adj. & sérieux, appliqué \\ \hline
\zh{仔细} & zǐxì & Adj. & attentif, minutieux \\ \hline
\zh{马虎} & mǎhu & Adj. & négligent, bâclé \\ \hline
\zh{流利} & liúlì & Adj. & courant (parler une langue) \\ \hline
\zh{累} & lèi & Adj. & fatigué \\ \hline
\zh{清楚} & qīngchu & Adj. & clair, distinct \\ \hline
\end{tabularx}
\end{center}
Tous ces mots sont des \textbf{adjectifs} en chinois ; ils peuvent donc fonctionner directement comme complément d'état après \zh{得}, sans verbe « être ».\par
\textbf{Point de vigilance :} \zh{马虎} se prononce \emph{mǎhu} (ton neutre sur la deuxième syllabe) ; c'est l'antonyme de \zh{认真}. On dit \zh{他做事做得很马虎} (« Il fait les choses de façon négligente »).
\end{corrige}

\begin{corrige}[Exercice 4 — Pinyin : Association caractères / prononciation]
1 -- B (\zh{开车} \emph{kāichē}, conduire) ; 2 -- C (\zh{做饭} \emph{zuòfàn}, cuisiner) ; 3 -- A (\zh{说汉语} \emph{shuō Hànyǔ}, parler chinois) ; 4 -- D (\zh{写字} \emph{xiězì}, écrire) ; 5 -- E (\zh{睡觉} \emph{shuìjiào}, dormir).\par
\textbf{Points de vigilance :}
\begin{itemize}
\item \zh{睡} \emph{shuì} (4\textsuperscript{e} ton) ne doit pas être confondu avec \zh{说} \emph{shuō} (1\textsuperscript{er} ton) : même initiale \emph{sh}, tons différents.
\item \zh{觉} est un caractère à deux lectures : \emph{jiào} dans \zh{睡觉} (dormir), \emph{jué} dans \zh{觉得} (penser, trouver).
\item \zh{饭} \emph{fàn} (4\textsuperscript{e} ton) : attention à ne pas écrire \emph{fǎn}.
\end{itemize}
\end{corrige}

\begin{corrige}[Exercice 5 — Traduction : Français $\rightarrow$ Chinois]
\begin{enumerate}
\item \zh{我弟弟开车开得很小心。} \emph{Wǒ dìdi kāichē kāi de hěn xiǎoxīn.} — Le verbe \zh{开} est répété parce qu'il a un objet \zh{车}.
\item \zh{她做饭做得很好吃。} \emph{Tā zuòfàn zuò de hěn hǎochī.} — Variante acceptée : \zh{她做菜做得很好吃。}
\item \zh{我们学中文学得很认真。} \emph{Wǒmen xué Zhōngwén xué de hěn rènzhēn.} — Variante : \zh{我们学习汉语学得很认真。}
\item \zh{他晚上睡得很晚。} \emph{Tā wǎnshang shuì de hěn wǎn.} — Le complément de temps \zh{晚上} se place avant le verbe, jamais après le complément d'état.
\item \zh{老师解释得很清楚。} \emph{Lǎoshī jiěshì de hěn qīngchu.} — Variante : \zh{老师讲得很清楚。}
\end{enumerate}
\textbf{Barème indicatif :} 2 points par phrase (1 point pour la structure V + \zh{得} + complément, 0,5 point pour l'ordre des mots, 0,5 point pour les caractères et le pinyin), soit 10 points.\par
\textbf{Points de vigilance :} ne jamais écrire \zh{地} après le verbe ; ne pas oublier la répétition du verbe quand il y a un objet ; ne pas traduire « très » par \zh{非常} partout, \zh{很} suffit.
\end{corrige}

\begin{corrige}[Exercice 6 — Transformation : Négation, Question, \zh{怎么样}]
\begin{enumerate}
\item Négation : \zh{妹妹跳舞跳得不美。} \emph{Mèimei tiàowǔ tiào de bù měi.} (« Ma petite sœur ne danse pas bien. ») — \zh{不} se place après \zh{得}, devant l'adjectif ; \zh{很} disparaît dans la négation.
\item Question fermée : \zh{妹妹跳舞跳得很美吗？} \emph{Mèimei tiàowǔ tiào de hěn měi ma ?} (« Ma petite sœur danse-t-elle bien ? »)
\item Question ouverte : \zh{妹妹跳舞跳得怎么样？} \emph{Mèimei tiàowǔ tiào de zěnmeyàng ?} (« Comment ma petite sœur danse-t-elle ? »)
\item Exclamation : \zh{妹妹跳舞跳得真美！} \emph{Mèimei tiàowǔ tiào de zhēn měi !} (« Comme ma petite sœur danse bien ! »)
\end{enumerate}
\textbf{Barème indicatif :} 1,5 point par transformation correcte, soit 6 points.\par
\textbf{Point de vigilance :} dans la question avec \zh{怎么样}, on ne garde ni \zh{很} ni l'adjectif : \zh{怎么样} remplace tout le complément.
\end{corrige}

\begin{corrige}[Exercice 7 — Texte à trous : Particules \zh{的} / \zh{地} / \zh{得}]
Tous les blancs se complètent avec \zh{得} : il s'agit chaque fois d'un complément d'état placé \emph{après} le verbe.\par
Texte complété : \zh{今天早上，我哥哥起得很早。他穿衣服穿得很快，吃饭吃得也很快，然后开车开得很小心，因为路上车很多。他工作得很努力，下班回家回得很累。但是他觉得汉语很有意思，所以晚上还学习得很认真。}\par
\textbf{Astuce de discrimination :}
\begin{itemize}
\item Après un verbe (évaluation de l'action) : \zh{得}.
\item Avant un verbe (manière) : \zh{地}, ex. \zh{他认真地学习}.
\item Avant un nom (détermination) : \zh{的}, ex. \zh{我的哥哥}.
\end{itemize}
\textbf{Barème indicatif :} 0,5 point par particule correcte (7 occurrences), soit 3,5 points.
\end{corrige}

\begin{corrige}[Exercice 8 — Remise en ordre : Ordre des mots]
\begin{enumerate}
\item \zh{他说中文说得很好。} \emph{Tā shuō Zhōngwén shuō de hěn hǎo.} (« Il parle très bien chinois. »)
\item \zh{弟弟跑步跑得很快。} \emph{Dìdi pǎobù pǎo de hěn kuài.} (« Mon petit frère court très vite. »)
\item \zh{我写字写得不是很好。} \emph{Wǒ xiězì xiě de bú shì hěn hǎo.} (« Je n'écris pas très bien. ») — La négation \zh{不} se place après \zh{得} ; on accepte aussi \zh{我写字写得不太好。}
\item \zh{同学们听老师讲课听得很仔细。} \emph{Tóngxuémen tīng lǎoshī jiǎngkè tīng de hěn zǐxì.} (« Les élèves écoutent le cours du professeur très attentivement. »)
\item \zh{你睡觉睡得怎么样？} \emph{Nǐ shuìjiào shuì de zěnmeyàng ?} (« Comment dors-tu ? »)
\end{enumerate}
\textbf{Barème indicatif :} 2 points par phrase (1 point pour la répétition du verbe, 1 point pour la place de \zh{得} et du complément), soit 10 points.\par
\textbf{Point de vigilance :} le schéma à respecter est Sujet + Verbe + Objet + Verbe + \zh{得} + Complément ; \zh{很} précède toujours l'adjectif du complément.
\end{corrige}

\begin{corrige}[Exercice 9 — Compréhension de texte et analyse linguistique]
\textbf{Questions de compréhension :}
\begin{enumerate}
\item Paul habite à Yaoundé, la capitale du Cameroun ; il étudie le chinois à l'université.
\item Parce que M\textsuperscript{me} Li explique ses cours très clairement (\zh{她讲课讲得非常清楚}).
\item Au début, les étudiants chinois parlaient trop vite et il ne comprenait pas (\zh{听不懂}) ; maintenant, il comprend beaucoup mieux (\zh{听得懂多了}).
\item Il veut partir étudier à Pékin l'année prochaine (\zh{他想明年去北京留学}).
\end{enumerate}
\textbf{Questions de langue :}
\begin{enumerate}
\setcounter{enumi}{4}
\item Trois compléments d'état du texte :
\begin{itemize}
\item \zh{讲得非常清楚} : verbe \zh{讲} + \zh{得} + \zh{非常清楚} — « explique très clairement » ;
\item \zh{说得很快} : verbe \zh{说} + \zh{得} + \zh{很快} — « parlent très vite » ;
\item \zh{写得很认真} : verbe \zh{写} + \zh{得} + \zh{很认真} — « écrit avec application ».
\end{itemize}
(On accepte aussi \zh{写得不太好}.)
\item \zh{写得不太好} : la négation \zh{不} se place \emph{après} \zh{得}, devant le groupe \zh{太好}. On utilise \zh{不} car on nie un état ou une qualité (écrire bien), et non l'accomplissement d'une action ; \zh{没} servirait à nier qu'une action a eu lieu (\zh{我没写}, « je n'ai pas écrit »).
\item \zh{讲课讲得清楚} est un \textbf{complément d'état} : verbe répété + \zh{得} + adjectif, qui évalue la manière dont l'action est faite. \zh{听得懂} est un \textbf{complément de résultat potentiel} : verbe + \zh{得} + verbe de résultat (\zh{懂}), qui exprime la \emph{capacité} d'atteindre le résultat (« pouvoir comprendre en écoutant ») ; sa négation est \zh{听不懂} (avec \zh{不}, sans \zh{得}).
\end{enumerate}
\textbf{Barème indicatif (sur 20) :} compréhension : 2 points par question (8 points) ; question 5 : 4 points ; question 6 : 4 points ; question 7 : 4 points.\par
\textbf{Points de vigilance :} bien distinguer \zh{得} du complément d'état (suivi d'un adjectif) et \zh{得} du potentiel (suivi d'un verbe de résultat) ; dans la réponse, citer toujours le verbe principal pour montrer la structure.
\end{corrige}

\begin{corrige}[Exercice 10 — Thème et Version]
\textbf{A. Thème :}
\begin{enumerate}
\item \zh{老师讲得非常清楚，我们听得很好。} \emph{Lǎoshī jiǎng de fēicháng qīngchu, wǒmen tīng de hěn hǎo.} — Variante : \zh{……我们都听懂了。}
\item \zh{我爸爸开车开得很快，但是很小心。} \emph{Wǒ bàba kāichē kāi de hěn kuài, dànshì hěn xiǎoxīn.}
\item \zh{他们做作业做得很仔细，所以做得很好。} \emph{Tāmen zuò zuòyè zuò de hěn zǐxì, suǒyǐ zuò de hěn hǎo.}
\item \zh{她说法语说得很流利，但是写字写得比较慢。} \emph{Tā shuō Fǎyǔ shuō de hěn liúlì, dànshì xiězì xiě de bǐjiào màn.}
\end{enumerate}
\textbf{B. Version :} « Xiao Zhang est mon collègue. Il conduit vraiment vite ; assis dans sa voiture, j'ai très peur. La semaine dernière, il m'a invité à manger. Il cuisine vraiment bien, surtout le porc braisé. Nous avons bavardé joyeusement et bu un peu trop. Le lendemain, il avait terriblement mal à la tête et il est arrivé au travail avec une demi-heure de retard. »\par
\textbf{Points de grammaire :}
\begin{itemize}
\item \zh{疼得厉害} : complément d'état de degré (« avoir terriblement mal ») ;
\item \zh{聊得很开心}、\zh{喝得有点多} : compléments d'état classiques ;
\item \zh{迟到了半个小时} : ici \zh{了} marque l'accomplissement et \zh{半个小时} est un complément de durée — ce n'est \emph{pas} un complément d'état, d'où l'absence de \zh{得}.
\end{itemize}
\textbf{Barème indicatif (sur 20) :} thème : 2,5 points par phrase (10 points) ; version : 10 points (exactitude du sens 6 points, qualité du français 4 points).\par
\textbf{Points de vigilance :} répéter le verbe devant l'objet (\zh{做作业做得……}) ; « couramment » pour une langue se rend par \zh{流利} ; ne pas confondre \zh{有点多} (« un peu trop ») et \zh{很多}.
\end{corrige}

\begin{corrige}[Exercice 11 — Expression écrite : Mon entourage]
\textbf{Proposition de rédaction modèle :}\par
\zh{我的妈妈在雅温得工作。她工作得非常努力，每天起得很早。她开车开得比较小心，因为路上车很多。晚上回家以后，她给我们做饭。她做饭做得特别好吃，我最喜欢她做的菜。周末，她还学习汉语。她学得很认真，说得也越来越流利。我很爱我的妈妈。}\par
\emph{Wǒ de māma zài Yǎwēndé gōngzuò. Tā gōngzuò de fēicháng nǔlì, měitiān qǐ de hěn zǎo. Tā kāichē kāi de bǐjiào xiǎoxīn, yīnwèi lùshang chē hěn duō. Wǎnshang huí jiā yǐhòu, tā gěi wǒmen zuòfàn. Tā zuòfàn zuò de tèbié hǎochī, wǒ zuì xǐhuan tā zuò de cài. Zhōumò, tā hái xuéxí Hànyǔ. Tā xué de hěn rènzhēn, shuō de yě yuè lái yuè liúlì. Wǒ hěn ài wǒ de māma.}\par
« Ma mère travaille à Yaoundé. Elle travaille très dur et se lève très tôt chaque jour. Elle conduit assez prudemment parce qu'il y a beaucoup de voitures sur la route. Le soir, en rentrant, elle nous cuisine. Elle cuisine particulièrement bien, j'adore ses plats. Le week-end, elle apprend aussi le chinois. Elle étudie sérieusement et parle de plus en plus couramment. J'aime beaucoup ma mère. »\par
\textbf{Compléments d'état utilisés :} \zh{工作得非常努力}、\zh{起得很早}、\zh{开得比较小心}、\zh{做得特别好吃}、\zh{学得很认真}、\zh{说得越来越流利} (six compléments, dont quatre avec verbe répété ou verbe seul).\par
\textbf{Barème indicatif (sur 20) :}
\begin{itemize}
\item Respect des contraintes (au moins 3 compléments d'état, 3 verbes d'activité, adverbes de degré, longueur 80--100 caractères) : 6 points ;
\item Correction grammaticale (structure V + \zh{得} + complément, répétition du verbe, ordre des mots) : 8 points ;
\item Exactitude des caractères et de la ponctuation chinoise : 4 points ;
\item Cohérence et richesse du texte : 2 points.
\end{itemize}
\textbf{Points de vigilance :}
\begin{itemize}
\item Ne jamais écrire \zh{地} à la place de \zh{得} après le verbe ;
\item Ne pas oublier de répéter le verbe quand il a un objet : \zh{做饭做得好吃}, et non \zh{做饭得很好吃} ;
\item Le complément de temps (\zh{每天}、\zh{周末}) se place avant le verbe, pas après le complément d'état ;
\item \zh{越来越} + adjectif exprime une progression : \zh{说得越来越流利} (« parle de plus en plus couramment »).
\end{itemize}
\end{corrige}

\newpage

\coursec{Fiche bilan}

\begin{popfigure}
\begin{tikzpicture}[mindmap, grow cyclic, every node/.style={concept, text=white, font=\small},
  root concept/.append style={concept color=popblue, font=\bfseries},
  level 1/.append style={level distance=3.2cm, sibling angle=51.4},
  level 1 concept/.append style={font=\footnotesize}]
\node[root concept, align=center]{Complément d'état\\ \zh{得} de}
  child[concept color=poporange]{ node[align=center]{Schéma\\ S + V + \zh{得} + Adj} }
  child[concept color=popgreen]{ node[align=center]{Emplois\\ manière, degré, résultat} }
  child[concept color=poppurple]{ node[align=center]{Négation\\ V + \zh{得} + \zh{不} + Adj} }
  child[concept color=poppink]{ node[align=center]{Question\\ V + \zh{得} + Adj + \zh{吗} ?\\ ou \zh{得怎么样}} }
  child[concept color=popgold]{ node[align=center]{Verbe + CO\\ répéter le verbe\\ \zh{他唱歌唱得很好}} }
  child[concept color=popblue!70]{ node[align=center]{Pièges\\ \zh{的}/\zh{得}/\zh{地}\\ pas de \zh{很} obligatoire} }
  child[concept color=popgreen!70!popblue]{ node[align=center]{Adverbes\\ \zh{得} + \zh{很}/\zh{非常}/\zh{太}…} };
\end{tikzpicture}
\legende{Carte mentale : le complément d'état avec \zh{得}}
\end{popfigure}

\begin{bilan}
\textbf{1. Lexique clé de la leçon}
\begin{center}
\begin{tabularx}{\linewidth}{|X|X|X|X|}
\hline
\rowcolor{popblueL} Caractères & Pinyin & Nature & Traduction \\
\hline
得 & de & particule & introduit le complément d'état \\
\hline
说 & shuō & verbe & parler, dire \\
\hline
写 & xiě & verbe & écrire \\
\hline
跑 & pǎo & verbe & courir \\
\hline
唱 & chàng & verbe & chanter \\
\hline
流利 & liúlì & adjectif & courant, fluide \\
\hline
清楚 & qīngchu & adjectif & clair \\
\hline
怎么样 & zěnmeyàng & pronom & comment \\
\hline
进步 & jìnbù & verbe/nom & progresser / progrès \\
\hline
\end{tabularx}
\end{center}

\textbf{2. Règles de grammaire à connaître par cœur}
\begin{center}
\begin{tabularx}{\linewidth}{|X|X|X|}
\hline
\rowcolor{popblueL} Structure & Exemple (caractères + pinyin) & Traduction \\
\hline
S + V + 得 + Adj & 他说得很快。Tā shuō de hěn kuài. & Il parle vite. \\
\hline
S + V + 得 + 不 + Adj (négation) & 她写得不好。Tā xiě de bù hǎo. & Elle écrit mal. \\
\hline
V + 得 + Adj + 吗 ? & 你跑得累吗？Nǐ pǎo de lèi ma ? & Es-tu fatigué d'avoir couru ? \\
\hline
V + 得 + 怎么样 ? & 他汉语说得怎么样？Tā Hànyǔ shuō de zěnmeyàng ? & Comment parle-t-il chinois ? \\
\hline
S + CO + V + 得 + Adj (verbe répété) & 他唱歌唱得很好。Tā chàng gē chàng de hěn hǎo. & Il chante bien (les chansons). \\
\hline
V + 得 + 很/非常/太 + Adj & 老师讲得很有意思。Lǎoshī jiǎng de hěn yǒu yìsi. & Le professeur explique de façon très intéressante. \\
\hline
\end{tabularx}
\end{center}

\textbf{3. Méthode express}
\begin{itemize}
\item Repérer le verbe d'action : le complément d'état évalue \emph{comment} l'action est faite.
\item Placer \zh{得} \textbf{immédiatement après le verbe}, jamais avant.
\item Si le verbe a un complément d'objet, répéter le verbe : \zh{他说汉语说得很流利。}
\item La négation porte sur l'adjectif : \zh{得} + \zh{不} + Adj.
\item Ne pas confondre : \zh{的} (de, possession/qualification), \zh{得} (de, complément d'état après verbe), \zh{地} (de, adverbe avant verbe).
\end{itemize}
\end{bilan}

\begin{aretenir}[Checklist avant le Bac]
\begin{itemize}
\item[$\square$] Je connais le schéma : Sujet + Verbe + \zh{得} + adjectif.
\item[$\square$] Je sais placer \zh{得} juste après le verbe, sans mot entre les deux.
\item[$\square$] Je sais former la négation : V + \zh{得} + \zh{不} + Adj.
\item[$\square$] Je sais poser la question avec \zh{吗} ou avec \zh{得怎么样}.
\item[$\square$] Je répète le verbe quand il a un complément d'objet : \zh{她唱歌唱得很好。}
\item[$\square$] Je distingue sans erreur \zh{的}, \zh{得} et \zh{地}.
\item[$\square$] Je sais renforcer le jugement avec \zh{很}, \zh{非常}, \zh{太}…\zh{了} après \zh{得}.
\item[$\square$] Je sais évaluer une action passée : \zh{他昨天跑得很快。}
\item[$\square$] Je sais traduire « Il parle bien le chinois » par \zh{他汉语说得很好。} (et non *\zh{他很好地说汉语}).
\item[$\square$] Je vérifie l'ordre des mots avant de rendre ma copie : verbe → \zh{得} → adjectif.
\end{itemize}
\end{aretenir}

\findecours

\end{document}
